% !TEX program = xelatex
% 运筹学 第四章 对偶理论与灵敏度分析 —— 可视化基础学习讲义
% 编译：xelatex 本文件（两次，生成目录与交叉引用）
\documentclass[11pt,a4paper,fontset=fandol,zihao=-4]{ctexart}
% ============================================================
%  导言区：宏包、配色、信息框、绘图样式
% ============================================================
\usepackage[a4paper,top=2.2cm,bottom=2.2cm,left=2.3cm,right=2.3cm,headheight=14pt]{geometry}
\usepackage{amsmath,amssymb,mathtools,bm}
\usepackage{booktabs,tabularx,array,multirow,longtable}
\usepackage[inline]{enumitem}
\usepackage[table]{xcolor}
\usepackage[most]{tcolorbox}
\usepackage{tikz}
\usetikzlibrary{arrows.meta,calc,positioning,shapes.geometric,decorations.pathreplacing,patterns,fit,backgrounds}
\usepackage{pgfplots}
\pgfplotsset{compat=1.18}
\usepgfplotslibrary{groupplots}
\usepackage{forest}
\usepackage{caption}
\usepackage{fancyhdr}
\usepackage{needspace}
\usepackage[colorlinks=true,linkcolor=navy,urlcolor=navy,citecolor=navy,bookmarksnumbered=true]{hyperref}

% ---------------- 配色 ----------------
\definecolor{navy}{HTML}{173A63}
\definecolor{auxblue}{HTML}{557A9E}
\definecolor{lightblue}{HTML}{F3F6F9}
\definecolor{formulablue}{HTML}{E8EFF7}
\definecolor{warm}{HTML}{FAF8F4}
\definecolor{warmline}{HTML}{C9B48A}
\definecolor{lightred}{HTML}{FFF5F4}
\definecolor{redline}{HTML}{C0504D}
\definecolor{lightgreen}{HTML}{F4F8F5}
\definecolor{greenline}{HTML}{5B8C6A}
\definecolor{grayblue}{HTML}{EEF2F6}
\definecolor{tagbase}{HTML}{DCE4EC}
\definecolor{tagup}{HTML}{E3EFE4}
\definecolor{tagcross}{HTML}{E9E5F1}
\definecolor{feas}{HTML}{DCEBDF}
\definecolor{optc}{HTML}{7A4FB5}
\definecolor{shiftc}{HTML}{C08A1E}
\definecolor{conA}{HTML}{2F6DAA}
\definecolor{badc}{HTML}{C0504D}

% ---------------- 章节格式 ----------------
\ctexset{
  section = {name={第,章},number=\arabic{section},format=\Large\bfseries\color{navy}\raggedright,beforeskip=1.2ex plus .3ex,afterskip=1ex},
  subsection = {format=\large\bfseries\color{navy}\raggedright,beforeskip=1.4ex plus .3ex,afterskip=.8ex},
  subsubsection = {format=\normalsize\bfseries\color{auxblue}\raggedright,beforeskip=1ex,afterskip=.5ex},
}
\setcounter{secnumdepth}{2}
\setcounter{tocdepth}{2}

\pagestyle{fancy}
\fancyhf{}
\fancyhead[L]{\small\color{auxblue}运筹学 · 对偶理论与灵敏度分析}
\fancyhead[R]{\small\color{auxblue}\leftmark}
\fancyfoot[C]{\small\thepage}
\renewcommand{\headrulewidth}{0.4pt}
\renewcommand{\headrule}{\hbox to\headwidth{\color{tagbase}\leaders\hrule height \headrulewidth\hfill}}
\renewcommand{\sectionmark}[1]{\markboth{#1}{}}

\setlist{itemsep=2pt,topsep=3pt,parsep=0pt,leftmargin=2em}
\captionsetup{font=small,labelfont={bf,color=navy},labelsep=quad,skip=4pt}
\setlength{\parskip}{3pt}
\linespread{1.18}

% ---------------- 信息框 ----------------
\tcbset{
  common/.style={enhanced,breakable,arc=2pt,boxrule=0.5pt,left=7pt,right=7pt,top=5pt,bottom=5pt,
    fonttitle=\bfseries\small,before skip=7pt,after skip=7pt,pad at break*=2pt},
}
% 一级：定义 / 定理 / 核心公式
\newtcolorbox{definitionbox}[1][定义]{common,colback=lightblue,colframe=navy,coltitle=white,colbacktitle=navy,title={#1}}
\newtcolorbox{theorembox}[1][定理]{common,colback=lightblue,colframe=navy,coltitle=white,colbacktitle=navy,title={#1}}
\newtcolorbox{formulabox}[1][核心公式]{common,colback=formulablue,colframe=formulablue,coltitle=navy,colbacktitle=formulablue,title={#1},boxrule=0pt,attach title to upper={\quad},}
% 二级：直观 / 推导 / 联系
\newtcolorbox{intuitionbox}[1][直观理解]{common,colback=grayblue,colframe=auxblue,boxrule=0pt,leftrule=2.5pt,arc=0pt,coltitle=auxblue,colbacktitle=grayblue,title={#1},fonttitle=\bfseries\small}
\newtcolorbox{proofbox}[1][证明]{common,colback=white,colframe=auxblue!60,boxrule=0pt,leftrule=1.5pt,arc=0pt,coltitle=auxblue,colbacktitle=white,title={#1}}
\newtcolorbox{connectionbox}[1][知识联系]{common,colback=grayblue,colframe=auxblue!70,boxrule=0pt,leftrule=2.5pt,arc=0pt,coltitle=navy,colbacktitle=grayblue,title={#1}}
\newtcolorbox{supplementbox}[1][教学补全]{common,colback=white,colframe=auxblue,boxrule=0.5pt,borderline={0.4pt}{1.5pt}{auxblue!40,dashed},coltitle=navy,colbacktitle=lightblue,title={#1}}
% 米色：示例 / 生活 / 历史 / 科普
\newtcolorbox{examplebox}[1][例题]{common,colback=warm,colframe=warmline,coltitle=black,colbacktitle=warm!70!warmline,title={#1}}
\newtcolorbox{lifebox}[1][生活中的它]{common,colback=warm,colframe=warmline,coltitle=black,colbacktitle=warm!60!warmline,title={\textbf{生活中的它}\quad #1}}
\newtcolorbox{historybox}[1][学科小史]{common,colback=warm,colframe=warmline,coltitle=black,colbacktitle=warm!70!warmline,title={#1}}
\newtcolorbox{sciencebox}[1][从课本到现实]{common,colback=warm,colframe=warmline,coltitle=black,colbacktitle=warm!70!warmline,title={#1}}
% 红：易错 / 纠错
\newtcolorbox{warningbox}[1][易错点]{common,colback=lightred,colframe=redline!70,coltitle=redline,colbacktitle=lightred,title={#1},boxrule=0.5pt}
\newtcolorbox{correctionbox}[1][原图纠错]{common,colback=lightred,colframe=redline!70,coltitle=white,colbacktitle=redline!80,title={#1}}
% 绿：总结 / 自检 / 方法
\newtcolorbox{summarybox}[1][方法总结]{common,colback=lightgreen,colframe=greenline,coltitle=greenline,colbacktitle=lightgreen,title={#1},boxrule=0.5pt}
\newtcolorbox{checkbox}[1][学完这里，你应该能够回答]{common,colback=lightgreen,colframe=greenline,coltitle=greenline,colbacktitle=lightgreen,title={#1},boxrule=0.5pt}

% 三级提示（细线 + 普通正文）
\newcommand{\tip}[2][提示]{\par\noindent{\color{auxblue}\rule[-.2ex]{1.5pt}{1.1em}}\hspace{4pt}\textbf{\color{auxblue}#1：}#2\par}

% 原始资料 / 补全 / 推测标签
\newcommand{\srcTag}[1]{\tikz[baseline=(t.base)]\node[fill=tagbase,rounded corners=1pt,inner sep=1.5pt,font=\scriptsize](t){原始资料·#1};}
\newcommand{\fillTag}{\tikz[baseline=(t.base)]\node[fill=formulablue,draw=auxblue,dashed,rounded corners=1pt,inner sep=1.5pt,font=\scriptsize](t){教学补全};}

% 练习区标签
\newcommand{\tagBase}[1]{\tikz[baseline=(t.base)]\node[fill=tagbase,rounded corners=2pt,inner sep=2pt,font=\small\bfseries,text=navy](t){#1};}
\newcommand{\tagUp}[1]{\tikz[baseline=(t.base)]\node[fill=tagup,rounded corners=2pt,inner sep=2pt,font=\small\bfseries,text=greenline!70!black](t){#1};}
\newcommand{\tagCross}[1]{\tikz[baseline=(t.base)]\node[fill=tagcross,rounded corners=2pt,inner sep=2pt,font=\small\bfseries,text=optc!70!black](t){#1};}
\newcommand{\diff}[1]{{\small\color{auxblue}[#1]}\ }

% 题号
\newcounter{qn}
\newcommand{\Q}[1][]{\par\needspace{4\baselineskip}\medskip\refstepcounter{qn}\noindent\textbf{\color{navy}\theqn.}\ #1}
\newcommand{\A}[1]{\par\needspace{5\baselineskip}\bigskip\noindent{\large\bfseries\color{navy}第 #1 题}\par\nobreak}
\newcommand{\opt}[4]{\par\noindent\begin{tabularx}{\linewidth}{@{}XX@{}}A.\ #1 & B.\ #2\\ C.\ #3 & D.\ #4\end{tabularx}\par}

% 图题
\newcommand{\fig}[1]{\par\nopagebreak[4]\captionof{figure}{#1}}

% 数学
\newcommand{\T}{^{\mathsf T}}
\newcommand{\Binv}{B^{-1}}
\newcommand{\st}{\text{s.t.}\quad}
\newcommand{\fr}[2]{\tfrac{#1}{#2}}

% ---------------- 绘图样式 ----------------
\tikzset{
  every picture/.append style={font=\small},
  ax/.style={-{Stealth[length=2mm]},thin,black!70},
  gridl/.style={black!8,very thin},
  con/.style={conA,thick},
  conB/.style={auxblue!70,thick},
  feasreg/.style={fill=feas,draw=none},
  optpt/.style={circle,fill=optc,draw=white,inner sep=2pt},
  vpt/.style={circle,fill=greenline,inner sep=1.4pt},
  badpt/.style={circle,fill=badc,draw=white,inner sep=2pt},
  shiftl/.style={shiftc,thick,dashed},
  iso/.style={black!55,thick,dashed},
  lab/.style={font=\scriptsize,fill=white,inner sep=1pt},
  box/.style={draw=navy,rounded corners=2pt,fill=lightblue,align=center,inner sep=4pt,font=\small},
  wbox/.style={draw=warmline,rounded corners=2pt,fill=warm,align=center,inner sep=4pt,font=\small},
  rbox/.style={draw=redline!70,rounded corners=2pt,fill=lightred,align=center,inner sep=4pt,font=\small},
  gbox/.style={draw=greenline,rounded corners=2pt,fill=lightgreen,align=center,inner sep=4pt,font=\small},
  arr/.style={-{Stealth[length=2mm]},thick,auxblue},
  hl/.style={fill=shiftc!25,draw=none},
}

% 坐标网格：\axes{xmax}{ymax}{xlabel}{ylabel}
\newcommand{\axes}[4]{%
  \draw[gridl] (0,0) grid (#1,#2);
  \draw[ax] (-0.3,0) -- (#1+0.4,0) node[right,font=\scriptsize]{$#3$};
  \draw[ax] (0,-0.3) -- (0,#2+0.4) node[above,font=\scriptsize]{$#4$};
  \foreach \i in {1,...,#1} \node[below,font=\tiny,black!60] at (\i,0) {\i};
  \foreach \i in {1,...,#2} \node[left,font=\tiny,black!60] at (0,\i) {\i};
  \node[below left,font=\tiny,black!60] at (0,0) {0};
}

% 讲义主例（材料 A、材料 B、设备台时）：max z=2x1+3x2
% 4x1<=16, 4x2<=12, x1+2x2<=8
\newcommand{\lecRegion}{\fill[feasreg] (0,0)--(4,0)--(4,2)--(2,3)--(0,3)--cycle;}
\newcommand{\lecLines}{%
  \draw[con] (4,-0.2)--(4,5.2) node[above,font=\scriptsize]{$4x_1=16$};
  \draw[con] (-0.2,3)--(9.3,3) node[right,font=\scriptsize]{$4x_2=12$};
  \draw[con] (-0.4,4.2)--(9.3,-0.65) node[right,font=\scriptsize]{$x_1+2x_2=8$};
}
% 练习模型 E：max z=3x1+2x2；x1+x2<=4, x1+3x2<=9, x1<=3
\newcommand{\exRegion}{\fill[feasreg] (0,0)--(3,0)--(3,1)--(1.5,2.5)--(0,3)--cycle;}
\newcommand{\exLines}{%
  \draw[con] (-0.2,4.2)--(4.5,-0.5) node[below,font=\scriptsize]{\ding1};
  \draw[con] (-0.3,3.1)--(6.3,0.9) node[right,font=\scriptsize]{\ding2};
  \draw[con] (3,-0.2)--(3,4.3) node[above,font=\scriptsize]{\ding3};
}
\newcommand{\ding}[1]{\textcircled{\scriptsize #1}}

% 解析区标题行：\SolHead{答案}{难度}{考查知识点}
\newcommand{\SolHead}[3]{\par\noindent\begin{tabularx}{\linewidth}{@{}>{\columncolor{lightblue}}l X@{}}
\textbf{正确答案} & \textbf{\color{navy}#1}\quad{\small 难度：#2}\quad{\small 考查：#3}\\\end{tabularx}\par}
\newcommand{\Lb}[1]{\par\noindent\textbf{\color{auxblue}#1}\ }
% 练习模型 E 的小图：\smallE{缩放}{附加绘图命令}
\newcommand{\smallE}[2]{\begin{tikzpicture}[x=#1cm,y=#1cm]\exRegion\draw[gridl] (0,0) grid (6,4);
  \draw[ax] (-0.2,0)--(6.3,0); \draw[ax] (0,-0.2)--(0,4.3);
  \draw[con] (-0.2,4.2)--(4.5,-0.5); \draw[con] (-0.3,3.1)--(6.3,0.9); \draw[con] (3,-0.2)--(3,4.3);#2\end{tikzpicture}}

% 特殊字符：带圈数字交给中文字体，勾叉用数学符号
\xeCJKsetcharclass{"2460}{"2473}{1}
\usepackage{newunicodechar}
\newunicodechar{✓}{\ensuremath{\checkmark}}
\newunicodechar{✗}{\ensuremath{\times}}


\hypersetup{pdftitle={对偶理论与灵敏度分析：可视化基础学习讲义},pdfauthor={运筹学专题讲义},pdfsubject={运筹学 第四章}}

\begin{document}
% ============================================================
%  封面、阅读说明、目录、专题定位
% ============================================================
\begin{titlepage}
\thispagestyle{empty}
\begin{tikzpicture}[remember picture,overlay]
  \fill[navy] (current page.north west) rectangle ([yshift=-0.9cm]current page.north east);
  \fill[auxblue!30] ([yshift=-0.9cm]current page.north west) rectangle ([yshift=-1.0cm]current page.north east);
\end{tikzpicture}
\vspace*{2.4cm}
{\color{auxblue}\large 运筹学 · 第四章专题讲义}\par\vspace{0.6cm}
{\fontsize{30}{38}\selectfont\bfseries\color{navy} 对偶理论与灵敏度分析}\par\vspace{0.4cm}
{\Large\color{navy!80} 修正单纯形法 · 对偶问题 · 影子价格 · 对偶单纯形法}\par\vspace{0.5cm}
{\large\color{black!70} 一份面向初学者的可视化基础学习讲义：看懂为什么，学会怎么算}\par
\vspace{1.2cm}
\begin{center}
\begin{tikzpicture}[x=0.72cm,y=0.72cm]
  % 原问题
  \begin{scope}
    \axes{6}{5}{x_1}{x_2}
    \lecRegion
    \draw[con] (4,-0.2)--(4,4.6);
    \draw[con] (-0.2,3)--(6.2,3);
    \draw[con] (-0.4,4.2)--(6.4,0.8);
    \node[optpt] at (4,2) {};
    \node[lab,above right] at (4,2) {$x^*=(4,2)$};
    \node[font=\small\bfseries,navy] at (3,-1.3) {原问题：$\max z$};
  \end{scope}
  % 对偶
  \begin{scope}[xshift=8.2cm]
    \fill[feas] (0,5)--(0,4)--(2,3)--(6,3)--(6,5)--cycle;
    \draw[gridl] (0,0) grid (6,5);
    \draw[ax] (-0.3,0) -- (6.4,0) node[right,font=\scriptsize]{$y_1$};
    \draw[ax] (0,-0.3) -- (0,5.4) node[above,font=\scriptsize]{$y_3$};
    \draw[con] (0,4)--(6.2,0.9);
    \draw[con] (-0.2,3)--(6.2,3);
    \node[optpt] at (2,3) {};
    \node[lab,above right] at (2,3) {$y^*=(\tfrac18,\tfrac32)$};
    \node[font=\small\bfseries,navy] at (3,-1.3) {对偶问题：$\min w$（$y_2=0$ 截面）};
  \end{scope}
  \draw[{Stealth[length=3mm]}-{Stealth[length=3mm]},very thick,shiftc] (7.2,2.5)--(10.2,2.5);
  \node[font=\small,align=center,shiftc!80!black] at (8.7,3.4) {同一组数据\\两种视角};
  \node[font=\small,align=center,navy] at (8.7,1.4) {$z^*=w^*$};
\end{tikzpicture}
\end{center}
\vspace{1.0cm}
\begin{tcolorbox}[enhanced,colback=lightblue,colframe=lightblue,arc=2pt,left=10pt,right=10pt]
\small
\begin{tabularx}{\linewidth}{@{}lX@{}}
\textbf{\color{navy}适用读者} & 第一次学习线性规划对偶理论、单纯形法基础尚不牢固的本科生\\
\textbf{\color{navy}前置知识} & 线性规划标准型、单纯形表与检验数、矩阵乘法与逆矩阵\\
\textbf{\color{navy}内容依据} & 课堂手写笔记（16 页）、课件《对偶单纯形法》（85 页）、课件《影子价格与灵敏度分析》（73 页），并按标准教材补全\\
\textbf{\color{navy}配套练习} & 38 道分层练习：基础 25 + 中等偏下 10 + 跨学科 3，附逐题、逐选项图解\\
\end{tabularx}
\end{tcolorbox}
\vfill
{\small\color{auxblue} 封面图：左为原问题可行域与最优顶点，右为对偶问题可行域（示意），二者最优值相等。}
\end{titlepage}

% ------------------------------------------------------------
\pagenumbering{roman}
\section*{阅读说明}
\addcontentsline{toc}{section}{阅读说明}
\markboth{阅读说明}{}

这份讲义把三份课堂资料（手写笔记、两套课件）重新组织成一条连续的学习路线。它不是课件的翻版：课件为了节省版面省略了不少证明、推导中间步骤和例子，这些内容在讲义里都已经补齐，并用“教学补全”标签单独标出，方便你区分“老师课上写过的”和“为了讲清楚而补上的”。

\paragraph{版面里的颜色是有含义的。}读几页之后你会自然习惯：
\begin{center}
\begin{tikzpicture}[every node/.style={minimum width=2.55cm,minimum height=0.85cm,font=\small,align=center}]
  \node[box] at (0,0) {定义 / 定理\\核心公式};
  \node[draw=auxblue,fill=grayblue,rounded corners=1pt] at (3,0) {为什么 / 直观\\推导 / 联系};
  \node[wbox] at (6,0) {例题 / 生活中的它\\学科史 / 科普};
  \node[rbox] at (9,0) {易错点\\原图纠错};
  \node[gbox] at (12,0) {方法总结\\学习自检};
\end{tikzpicture}
\fig{颜色语义：蓝 = 核心知识，蓝灰左线 = 理解与联系，米色 = 示例与应用，红 = 易错，绿 = 总结与检查}
\end{center}

\paragraph{三种标签。}\srcTag{笔记}\,/\,\srcTag{课件} 表示内容可以在原始资料中可靠找到；\fillTag\ 表示原始资料没有完整写出、依据标准教材补充的内容；正文中若出现“推测”字样，表示该处原图无法可靠辨认。

\paragraph{建议的使用方法。}
\begin{enumerate}
  \item 第一遍按顺序读正文，每读完一章做“学完这里，你应该能够回答”的自检；
  \item 第二遍只看图：每一章的图连起来就是这一章的骨架；
  \item 做练习时先独立完成全部 38 题，再查“答案速查”，最后读详细解析。解析区的图与题目区的图不同，解析图上画出了解题时添加的辅助信息。
\end{enumerate}

\tableofcontents
\clearpage
\pagenumbering{arabic}

% ------------------------------------------------------------
\section*{本专题学什么}
\addcontentsline{toc}{section}{本专题学什么、它在课程中的位置}
\markboth{本专题学什么}{}

前几章我们学会了用单纯形法求出一个线性规划的最优解。但一个真实的决策者拿到最优解之后，马上会追问三类问题：
\begin{enumerate}[label=\textbf{问题\arabic*},leftmargin=4.6em]
  \item \textbf{能不能算得更省？}单纯形表每一步都要把整张表重算一遍，很多数字其实用不上。（$\to$ 修正单纯形法）
  \item \textbf{资源到底值多少钱？}多给一单位原料、一小时设备，利润能多多少？有没有人愿意出价买走我的资源？（$\to$ 对偶问题、影子价格）
  \item \textbf{数据变了怎么办？}价格、库存、工艺都会变化，最优解还成立吗？能不能不从头算？（$\to$ 对偶单纯形法、灵敏度分析）
\end{enumerate}
本专题回答的就是这三类问题，而三类问题的答案都写在同一个对象里——\textbf{最优基的逆矩阵 $B^{-1}$}。这是贯穿全篇的一条主线。

\begin{center}
\begin{tikzpicture}[node distance=0.45cm and 0.5cm,
  k/.style={box,minimum width=2.6cm,minimum height=0.95cm},
  p/.style={draw=black!25,rounded corners=2pt,fill=white,align=center,font=\small,minimum width=2.4cm,minimum height=0.85cm}]
  \node[p] (lp) {线性规划\\标准型};
  \node[p,below=of lp] (sm) {单纯形法\\检验数 $\sigma_j$};
  \node[p,below=of sm] (mat) {矩阵乘法\\逆矩阵};
  \node[k,right=1.2cm of sm] (B) {\textbf{最优基的逆}\\$B^{-1}$};
  \node[k,above right=0.2cm and 1.1cm of B] (rev) {修正单纯形法};
  \node[k,right=1.1cm of B] (dual) {对偶问题\\影子价格};
  \node[k,below right=0.2cm and 1.1cm of B] (ds) {对偶单纯形法};
  \node[k,right=0.9cm of dual] (sa) {灵敏度分析};
  \draw[arr] (lp)--(B); \draw[arr] (sm)--(B); \draw[arr] (mat)--(B);
  \draw[arr] (B)--(rev); \draw[arr] (B)--(dual); \draw[arr] (B)--(ds);
  \draw[arr] (dual)--(sa); \draw[arr] (ds.east) -| (sa.south);
  \node[font=\scriptsize,auxblue,above=1pt of lp] {前置知识};
  \node[font=\scriptsize,auxblue,above=1pt of rev] {第 1 章};
  \node[font=\scriptsize,auxblue,above=1pt of sa] {第 6 章};
\end{tikzpicture}
\fig{前置知识地图：左列是需要提前会的内容，它们共同汇聚到 $B^{-1}$，再分出本专题的四个模块}
\end{center}

图中箭头表示“学习依赖”：例如对偶单纯形法既要用到单纯形表的旋转运算，也要用到对偶可行的概念，所以它排在对偶理论之后；灵敏度分析需要影子价格和对偶单纯形法两样工具，因此放在最后。

\begin{connectionbox}[本专题在课程体系中的位置]
\begin{tabularx}{\linewidth}{@{}>{\bfseries\color{navy}}l X@{}}
所属课程 & 运筹学（线性规划部分），第四章\\
前面学过 & 第一、二章：线性规划模型、图解法、单纯形法、大 $M$ 法与两阶段法\\
后面会用 & 整数规划（分支定界中的“加约束再求解”正是对偶单纯形法）、运输问题（位势法本质是对偶变量）、非线性规划（KKT 条件与拉格朗日乘子）\\
相关课程 & 微观经济学（边际收益、机会成本）、最优化方法（拉格朗日对偶）、电力系统分析（经济调度与节点电价）、博弈论（零和博弈与极小极大定理）\\
\end{tabularx}
\end{connectionbox}

% ============================================================
%  导论：学科小史、现实应用、原始资料重编码与纠错、知识总览
% ============================================================
\section*{导论：问题从哪里来}
\addcontentsline{toc}{section}{导论：学科小史、现实应用与原始资料}
\markboth{导论}{}

\subsection*{学科小史：这个知识是怎样出现的？}
\addcontentsline{toc}{subsection}{学科小史：这个知识是怎样出现的？}

对偶理论并不是先有数学、后找应用。它是在“怎样给稀缺资源定价”和“怎样快速求解大规模线性规划”这两个现实问题的推动下逐步形成的。

\begin{center}
\begin{tikzpicture}[x=1cm,y=1cm]
  \draw[very thick,navy] (0,0)--(15.6,0);
  \foreach \x/\yr/\t/\d in {
    0.8/1939/{康托罗维奇}/{生产计划中的“解乘数”\\为资源定出内部估价},
    4.0/1947/{丹齐格}/{提出单纯形法；与冯·诺伊曼\\交谈中出现对偶思想},
    7.2/1951/{盖尔、库恩、塔克}/{发表线性规划\\对偶定理的证明},
    10.4/1954/{莱姆克}/{提出对偶单纯形法\\（dual simplex method）},
    13.6/1950s--{今}/{修正单纯形法与计算机}/{RAND 等机构把只维护 $B^{-1}$\\的算法写成程序，延续至今}}{
    \fill[navy] (\x,0) circle (3pt);
    \node[font=\small\bfseries,navy,above] at (\x,0.15) {\yr};
    \node[font=\small,above,text width=3cm,align=center] at (\x,0.6) {\t};
    \node[font=\scriptsize,below,text width=3.1cm,align=center,black!75] at (\x,-0.2) {\d};
  }
\end{tikzpicture}
\fig{对偶理论发展的五个关键节点（时间轴只标出与本章直接相关的事件）}
\end{center}

\begin{historybox}[人们遇到了什么问题？]
\begin{itemize}
  \item \textbf{资源怎么估价。}1939 年，苏联数学家康托罗维奇（L.~V.~Kantorovich）研究胶合板厂的生产安排时，引入了一组“解乘数”，它们衡量每种资源对最优产出的贡献。这正是今天“影子价格”的前身。他与库普曼斯（T.~C.~Koopmans）因资源最优配置理论共同获得 1975 年诺贝尔经济学奖。
  \item \textbf{怎么求解。}1947 年，丹齐格（G.~B.~Dantzig）为美国空军的计划问题提出单纯形法。据丹齐格本人回忆，同年他向冯·诺伊曼（J.~von Neumann）介绍线性规划时，后者立刻把它与自己的博弈论联系起来，指出每个线性规划都伴随着另一个“对偶”的线性规划。
  \item \textbf{严格证明与新算法。}1951 年，盖尔、库恩、塔克（D.~Gale, H.~W.~Kuhn, A.~W.~Tucker）发表了对偶定理的证明；1954 年，莱姆克（C.~E.~Lemke）提出对偶单纯形法；同一时期，丹齐格与奥查德-海斯（W.~Orchard-Hays）等人在 RAND 公司发展了只维护基逆矩阵的修正单纯形法，以适应当时内存极小的计算机。
\end{itemize}
\textbf{为什么今天还要学？}因为现代求解器求解线性规划、整数规划时，大量时间都花在“改一点数据—从旧的最优表继续算”上，而这正是对偶单纯形法与灵敏度分析的用武之地。
\end{historybox}
\tip[说明]{“谁最早提出对偶”在文献中有不同叙述，上面采用较保守的说法：对偶思想在 1947 年丹齐格与冯·诺伊曼的交流中出现，严格证明于 1951 年发表。}

\needspace{14\baselineskip}
\subsection*{从课本到现实：这个知识到底有什么用？}
\addcontentsline{toc}{subsection}{从课本到现实：这个知识到底有什么用？}

\begin{sciencebox}[三个真实系统里的对偶变量]
\begin{minipage}[t]{0.58\linewidth}
\textbf{1. 电力市场的“节点电价”。}电网调度机构每隔一段时间求解一个“满足负荷、成本最低”的优化问题。每个节点“供需平衡”约束对应的对偶变量，就是该节点再多用 1 兆瓦时电能所增加的系统成本，许多电力市场直接把它作为结算电价（称为节点边际电价，LMP）。本讲义的跨学科第 36 题就是它的简化版。

\smallskip
\textbf{2. 航空公司的“票价门槛”。}一架航班的座位是有限资源。收益管理系统先求解一个线性规划，用各航段座位约束的对偶价格估计“一个座位的机会成本”。如果一张联程票的价格低于它占用的各段座位影子价格之和，系统就倾向于不卖。

\smallskip
\textbf{3. 求解器里的“重新优化”。}整数规划的分支定界法会反复往线性规划里加一条约束再求解。每加一条约束，旧的最优表就变成“对偶可行而原始不可行”的状态，用对偶单纯形法通常几步就能恢复最优——这正是第 6 章“新增约束”的内容。
\end{minipage}\hfill
\begin{minipage}[t]{0.38\linewidth}
\vspace{2pt}
\begin{tikzpicture}[x=1cm,y=1cm]
  \node[box,minimum width=3.6cm] (m) at (0,0) {线性规划\\（原问题）};
  \node[wbox,minimum width=3.6cm] (a) at (0,-1.7) {对偶变量 $y_i$};
  \node[gbox,minimum width=1.6cm,font=\scriptsize] (p) at (-1.4,-3.5) {节点电价};
  \node[gbox,minimum width=1.6cm,font=\scriptsize] (q) at (0.1,-3.5) {座位机会\\成本};
  \node[gbox,minimum width=1.6cm,font=\scriptsize] (r) at (1.6,-3.5) {重新优化\\的起点};
  \draw[arr] (m)--node[right,font=\scriptsize]{求解 / 读表}(a);
  \draw[arr] (a)--(p); \draw[arr] (a)--(q); \draw[arr] (a)--(r);
\end{tikzpicture}
\fig{课堂上的 $y_i$ 在三个系统中的身份}
\end{minipage}
\end{sciencebox}

本节讲的是整个专题在现实系统中的位置；每个知识点各自的生活例子放在相应章节的“生活中的它”模块里，二者不重复。

% ------------------------------------------------------------
\subsection*{原始资料重新编码}
\addcontentsline{toc}{subsection}{原始资料重新编码与纠错}

本讲义依据三份原始资料。下表说明每份资料讲了什么、对应讲义哪一章，以及讲义补充了什么。

\begin{center}
\small
\begin{tabularx}{\linewidth}{@{}>{\raggedright\arraybackslash}p{2.2cm}>{\raggedright\arraybackslash}X>{\raggedright\arraybackslash}p{1.5cm}>{\raggedright\arraybackslash}X@{}}
\toprule
\rowcolor{lightblue}\textbf{原始资料} & \textbf{可靠识别的内容} & \textbf{对应章节} & \textbf{讲义补充（教学补全）}\\
\midrule
手写笔记 P1--P4（10 月 14--15 日） & 修正单纯形法思路；初等行变换 = 左乘初等矩阵；例题 $\max z=5x_1+2x_2+3x_3-x_4+x_5$ 的单纯形表与修正单纯形法对照计算 & 第 1 章 & 检验数公式 $\sigma=c-c_B\T B^{-1}A$ 的推导；修正单纯形法完整步骤；$2\times2$ 逆矩阵求法\\
\addlinespace
手写笔记 P5--P11（10 月 19 日） & 煤电油引例与购买者模型；对偶问题对称形式；一般形式的对应规则表；例题 1、2；对偶的六条性质；对偶最优解与原问题终表检验数的关系 & 第 2、3 章 & 弱对偶性证明中 $\bar y\ge0$ 的必要性；强对偶定理证明；无界性的证明与反例；互补松弛定理的一般证明\\
\addlinespace
手写笔记 P12--P13（10 月 19--20 日） & 煤电油问题终表；影子价格的定义、数学表示、经济意义三条 & 第 4 章 & 影子价格 $=\partial z^*/\partial b_i$ 的推导；适用范围\\
\addlinespace
手写笔记 P14--P16（10 月 21 日） & 对偶单纯形法的特点、与单纯形法对照、五步操作；例题 $\min z=2x_1+3x_2+4x_3$ & 第 5 章 & 入基比值规则的推导；无可行解判据的证明\\
\addlinespace
课件《对偶单纯形法》85 页 & 例题 $\min w=2000y_1+1000y_2$ 的完整迭代；原始可行/对偶可行；出基、入基判据推导；几何路径；与原问题三张表逐格对照；新增约束 $x_2\le500$ 的重新优化 & 第 5、6 章 & 判据背后的对偶解释整合进正文\\
\addlinespace
课件《影子价格与灵敏度分析》73 页 & 例题 $\max z=2x_1+3x_2$（材料 $A$、$B$、设备台时）；$B^{-1}$ 的读取；影子价格计算与几何验证；互补松弛四步推导；灵敏度分析三种情形、两条判据、四个步骤；$b_1,b_2,b_3,c_1,c_2$ 的范围 & 第 4、6 章 & 新增变量、新增约束的公式；价值系数范围的几何解释统一\\
\bottomrule
\end{tabularx}
\end{center}

\paragraph{无法可靠辨认的内容。}手写笔记第 1 页“回顾”一行末尾被纸张边缘截断，可辨认出“人工变量法/退化/整数规划”，其后文字\textbf{原始资料此处无法可靠辨认}，不影响本章内容。第 5 页引例表格中“煤”的资源限制数字字迹有涂改，但后文模型统一写作 $9x_1+4x_2\le360$，终表 $x_3=84$ 也与 360 吻合，因此按 360 处理。

% ------------------------------------------------------------
\subsection*{原图纠错}

逐页验算原始资料后，发现以下问题。每一处都不影响课上得到的最终答案，但如果照抄，会在自己做题时出错。

\begin{correctionbox}[纠错 1　手写笔记 P2：单纯形表第三张表 $x_3$ 行的 $b$ 值]
\textbf{原始写法：}第三张表中 $x_3$ 行 $B^{-1}b=\frac{12}{5}$，右侧框出的 $B_3^{-1}$ 左上角为 $\frac25$。\par
\textbf{存在的问题：}由上一张表做旋转运算：$x_3$ 行 $=$ 原 $x_3$ 行 $-\frac12\times$（新 $x_1$ 行），得 $b=4-\frac12\cdot\frac65=\frac{17}{5}$，$x_4$ 列 $=\frac12-\frac12\cdot(-\frac15)=\frac35$。\par
\textbf{规范写法：}$x_3$ 行 $B^{-1}b=\frac{17}{5}$，$B_3^{-1}=\begin{pmatrix}\frac35&-\frac15\\[2pt]-\frac15&\frac25\end{pmatrix}$。\par
\textbf{为什么这样改：}最优值验算 $z^*=3\times\frac{17}{5}+5\times\frac65=\frac{51}{5}+\frac{30}{5}=\frac{81}{5}$，与笔记最终答案一致；若用 $\frac{12}{5}$ 则得 $\frac{66}{5}$，矛盾。\par
\textbf{是否影响后续：}不影响检验数行与最终答案，笔记下一页修正单纯形法求出的 $x^*=(\frac65,0,\frac{17}{5},0,0)$ 是正确的。
\end{correctionbox}

\begin{correctionbox}[纠错 2　手写笔记 P3：$B^{(3)-1}_{3,1}$ 求错]
\textbf{原始写法：}$B_{3,1}^{(3)}=\begin{pmatrix}2&1\\1&3\end{pmatrix}$，$B_{3,1}^{(3)-1}=\begin{pmatrix}-\frac25&\frac15\\[2pt]\frac15&-\frac35\end{pmatrix}$。\par
\textbf{存在的问题：}$\det B=2\times3-1\times1=5$，按 $2\times2$ 逆矩阵公式应为 $\frac15\begin{pmatrix}3&-1\\-1&2\end{pmatrix}$。可用乘法检验：原写法与 $B$ 相乘不得单位阵。\par
\textbf{规范写法：}$B^{-1}=\begin{pmatrix}\frac35&-\frac15\\[2pt]-\frac15&\frac25\end{pmatrix}$，于是 $x_B=B^{-1}b=\left(\frac{24-7}{5},\frac{-8+14}{5}\right)\T=\left(\frac{17}{5},\frac65\right)\T$。\par
\textbf{是否影响后续：}笔记 P4 的数值结果正确，说明计算时实际使用的是正确的逆矩阵，仅书写有误。
\end{correctionbox}

\begin{correctionbox}[纠错 3　手写笔记 P13：煤电油终表 $x_3$ 行]
\textbf{原始写法：}$x_3$ 行 $x_4$ 列为 $-0.32$。\quad
\textbf{规范写法：}$-3.12$。\par
\textbf{为什么：}由 $x_1=20-0.4x_4+0.2x_5$、$x_2=24+0.12x_4-0.16x_5$ 代入煤约束 $x_3=360-9x_1-4x_2$，得 $x_3=84+3.12x_4-1.16x_5$，表中系数为 $-3.12$ 与 $1.16$。检验数 $-1.36,-0.52$ 不受影响。
\end{correctionbox}

\begin{correctionbox}[纠错 4　手写笔记 P16：对偶单纯形法例题的最优值符号]
\textbf{原始写法：}“$z'^*=-\frac{28}{5}$，$\max z'=\min -z=\frac{28}{5}$”。\par
\textbf{存在的问题：}$z'=-z$，所以 $\max z'=-\min z$。由 $\max z'=-\frac{28}{5}$ 应得 $\min z=\frac{28}{5}$。原写法把两个量的符号写混了。\par
\textbf{规范写法：}$\max z'=-\frac{28}{5}$，故原问题 $\min z=\frac{28}{5}$。\par
\textbf{附带说明：}同页第二轮 $\theta$ 比值中 $x_5$ 一项被打叉，实际上 $a_{r5}=-\frac12<0$，比值 $\frac{-1}{-1/2}=2$ 应参与比较；因 $\frac85<2$，入基变量仍是 $x_2$，结论不变。
\end{correctionbox}

\begin{correctionbox}[纠错 5　手写笔记 P10--P11：对偶约束与对偶变量的记号]
\textbf{原始写法：}例题 1、2 的“套路”中写 $g_1(y)\le b_1$、$g_2(y),g_3(y)\ge b_2,b_3$；例题 2 第 4 步写“$f_1(x)\ge5\to x_1\ge0$，$f_2(x)\le4\to x_2\le0$，$f_3(x)=6\to x_3$ 无约束”。\par
\textbf{存在的问题：}对偶约束的右端是原问题的价值系数 $c_j$，不是 $b_j$；原问题约束对应的是对偶\textbf{变量} $y_i$，不是 $x_i$。\par
\textbf{规范写法：}$g_j(y)\ \square\ c_j$；例题 2 第 4 步应为 $y_1\ge0$，$y_2\le0$，$y_3$ 无约束（笔记 P10 上方列出的对偶模型本身是正确的）。
\end{correctionbox}

\begin{correctionbox}[纠错 6　表述不严谨的几处]
\begin{enumerate}[leftmargin=1.5em]
  \item \textbf{笔记 P2“初始可行基必须是单位阵”：}修正单纯形法只要求初始基 $B$ 可逆、$B^{-1}b\ge0$ 且 $B^{-1}$ 容易得到。单位阵最方便（$B^{-1}=I$），但不是必须。
  \item \textbf{笔记 P11 弱对偶性证明：}由 $A\bar x\le b$ 左乘 $\bar y$ 得 $\bar yA\bar x\le\bar yb$，这一步\textbf{需要 $\bar y\ge0$}（不等式两边乘负数要变号）。笔记写对偶模型时才给出 $y\ge0$，应在相乘前说明。
  \item \textbf{笔记 P13“影子价格 $=0$ 表示这种资源有剩余”：}反过来才总成立——资源有剩余 $\Rightarrow$ 影子价格为 0。退化时资源恰好用尽、影子价格也可能为 0（见第 3 章例子）。
  \item \textbf{笔记 P14--P15：}“保证初始可行解最优性”应为“保证初始\textbf{基解}满足最优性条件”——对偶单纯形法的初始基解恰恰不可行；“前提不满足，不可使用单纯形法”应为“不可使用\textbf{对偶}单纯形法”；单纯形法最优判据“所有 $\sigma_i<0$”应为 $\sigma_i\le0$。
  \item \textbf{课件 P69：}“$\Delta c_2\le\Delta c_2\le1$”为排版笔误，应为 $\Delta c_2\le1$。
  \item \textbf{课件 P52 与 P66 记号冲突：}P52 用 $\beta_r$ 表示 $B^{-1}$ 的第 $r$ \textbf{列}，P66 又用 $\beta_r\T$ 表示 $B^{-1}$ 的第 $r$ \textbf{行}。本讲义统一记 $\beta_r$ 为第 $r$ 列，第 $r$ 行记作 $(B^{-1})_{r\cdot}$。
  \item \textbf{课件 P13 红字“单纯形表的 $b$ 列的值，不是式子中的 $b$ 的值”：}含义应理解为：公式 $x_B=B^{-1}b$ 中的 $b$ 是\textbf{原模型的右端项} $(16,12,8)\T$，算出的结果才等于最优表 $b$ 列 $(4,2,4)\T$。二者不要混用。
\end{enumerate}
\end{correctionbox}

% ------------------------------------------------------------
\subsection*{知识点总览与关系图}
\addcontentsline{toc}{subsection}{知识点总览与关系图}

\begin{center}
\begin{tikzpicture}[node distance=0.55cm and 0.45cm,
  k/.style={box,minimum width=2.75cm,minimum height=1.0cm,text width=2.6cm}]
  \node[k] (K1) {\textbf{K1} 矩阵视角\\修正单纯形法};
  \node[k,right=of K1] (K2) {\textbf{K2} 对偶问题\\的构造};
  \node[k,right=of K2] (K3) {\textbf{K3} 对偶的\\基本性质};
  \node[k,right=of K3] (K4) {\textbf{K4} 互补松弛};
  \node[k,below=1.1cm of K2] (K5) {\textbf{K5} 影子价格};
  \node[k,below=1.1cm of K3] (K6) {\textbf{K6} 对偶单纯形法};
  \node[k,below=1.1cm of K4] (K7) {\textbf{K7--K9} 灵敏度分析\\$b$、$c$、新增};
  \node[k,below=1.1cm of K1] (Y) {$y\T=c_B\T B^{-1}$\\（全篇枢纽）};
  \draw[arr] (K1)--(K2); \draw[arr] (K2)--(K3); \draw[arr] (K3)--(K4);
  \draw[arr] (K1)--(Y); \draw[arr] (Y)--(K5); \draw[arr] (K3)--(K5); \draw[arr] (K4)--(K5);
  \draw[arr] (K3)--(K6); \draw[arr] (K5)--(K6); \draw[arr] (K6)--(K7); \draw[arr] (K4)--(K7);
  \draw[arr] (K5) to[bend right=25] (K7);
\end{tikzpicture}
\fig{九个核心知识点之间的依赖关系。箭头 $A\to B$ 表示学 $B$ 时要用到 $A$}
\end{center}

整份讲义按图中从左上到右下的顺序展开。读者可以把这张图当作“地铁线路图”：每学完一站，回来看看自己走到了哪里。

% ============================================================
%  第 1 章  矩阵视角与修正单纯形法（K1）
% ============================================================
\section{矩阵视角：单纯形表里藏着什么}

前两章我们一直在单纯形表上“画圈、换基、消元”。这套操作能算出答案，但它像一台黑箱：每一轮整张表都被重写一遍，我们说不清表里每个数字到底是什么。本章回答三个问题：
\begin{enumerate}
  \item 单纯形表中的每一个数，能不能用原始数据 $A,b,c$ 直接写出来？
  \item 如果能，是不是只要记住一个小矩阵，就可以不再维护整张大表？
  \item 这个小矩阵为什么会成为后面对偶理论、灵敏度分析的“总钥匙”？
\end{enumerate}
\noindent\textbf{本章知识路线：}初等行变换 $\to$ 左乘矩阵 $\to$ 单纯形表的矩阵表示 $\to$ 修正单纯形法。

\subsection{初等行变换就是左乘一个矩阵}

\srcTag{笔记 P1}\ 笔记里写道：“任何一种初等行变换都可以被若干次特定顺序的矩阵乘法来表示”，并给了一个交换两行的例子。先看图。

\begin{center}
\begin{tikzpicture}[every node/.style={font=\small}]
  \node (A) at (0,0) {$\begin{pmatrix}1&2&3\\4&5&6\\7&8&9\end{pmatrix}$};
  \node (B) at (5.4,0) {$\begin{pmatrix}4&5&6\\1&2&3\\7&8&9\end{pmatrix}$};
  \draw[arr] (A.east)--node[above]{交换第 1、2 行}(B.west);
  \node (E) at (0,-2.5) {$\underbrace{\begin{pmatrix}0&1&0\\1&0&0\\0&0&1\end{pmatrix}}_{E\text{：把单位阵也交换 1、2 行}}$};
  \node (A2) at (3.0,-2.35) {$\begin{pmatrix}1&2&3\\4&5&6\\7&8&9\end{pmatrix}$};
  \node at (4.6,-2.35) {$=$};
  \node (B2) at (6.2,-2.35) {$\begin{pmatrix}4&5&6\\1&2&3\\7&8&9\end{pmatrix}$};
  \node[font=\scriptsize,shiftc!80!black,align=left] at (9.4,-2.35) {$EA$ 第 1 行 $=0\cdot(\text{第1行})$\\$+1\cdot(\text{第2行})+0\cdot(\text{第3行})$};
  \draw[arr,shiftc] (B.south)--(B2.north);
  \node[font=\scriptsize,auxblue,right] at (5.6,-1.0) {结果相同};
\end{tikzpicture}
\fig{初等行变换 $=$ 左乘初等矩阵：先对单位阵 $I$ 做同样的行变换得到 $E$，再用 $E$ 左乘 $A$}
\end{center}

图上半部分是我们熟悉的“直接交换两行”；下半部分是同一件事的矩阵写法：$E$ 由单位阵交换 1、2 行得到，$EA$ 的第 1 行等于 $E$ 的第 1 行 $(0,1,0)$ 去“挑选”$A$ 的第 2 行。右侧注释写出了第 1 行的“挑选”过程。

\begin{definitionbox}[结论 1.1　初等矩阵（笔记中称“左行右列”）]
对矩阵 $A$ 做一次初等\textbf{行}变换，等于用对应的初等矩阵 $E$ \textbf{左}乘 $A$；做一次初等\textbf{列}变换，等于\textbf{右}乘。三种初等行变换对应的 $E$ 都由单位阵 $I$ 做同样的变换得到：交换两行；某行乘非零数 $k$；某行加上另一行的 $k$ 倍。
\end{definitionbox}

\begin{proofbox}[为什么成立 \fillTag]
矩阵乘法的规则是：$EA$ 的第 $i$ 行 $=\sum_k e_{ik}\,(A\text{ 的第 }k\text{ 行})$，即 $E$ 的第 $i$ 行给出“用 $A$ 的哪些行、各乘多少”。单位阵的第 $i$ 行是 $e_i\T$，只挑出 $A$ 的第 $i$ 行，所以 $IA=A$。把 $I$ 的行按某种初等变换改动，挑选方式就跟着改动，$EA$ 的行便恰好按同样方式改动。由于单纯形法中每次旋转都是若干次初等行变换，所以\textbf{一整轮旋转等于左乘若干个初等矩阵的乘积}。
\end{proofbox}

\tip[笔记原话]{“尽量做行变换，不要做列变换”——因为列对应变量，交换列会打乱变量与列的对应关系；而行变换只是在重写同一组方程。}

\subsection{单纯形表的矩阵表示}

有了“旋转 $=$ 左乘”，我们就能回答本章第一个问题。设标准型为
\[
\max z=c\T x,\qquad Ax=b,\quad x\ge0,
\]
其中 $A$ 是 $m\times n$ 矩阵。某一轮选定了 $m$ 个基变量，把 $A$ 的列重新分成两块：基变量的列组成\textbf{基矩阵} $B$（$m\times m$，可逆），其余列组成 $N$。相应地 $x=(x_B,x_N)$，$c=(c_B,c_N)$。

\begin{proofbox}[推导：从 $Ax=b$ 到单纯形表 \fillTag]
\textbf{第一步（解出基变量）。}$Ax=b$ 写成 $Bx_B+Nx_N=b$。两边左乘 $B^{-1}$（$B$ 可逆）：
\begin{align}
x_B+B^{-1}Nx_N&=B^{-1}b. \label{eq:xb}
\end{align}
这正是单纯形表中“约束行”的内容：基变量的列变成单位阵，非基变量的列是 $B^{-1}N$，右端是 $B^{-1}b$。

\textbf{第二步（把目标函数写成只含非基变量）。}由 \eqref{eq:xb} 得 $x_B=B^{-1}b-B^{-1}Nx_N$，代入目标函数：
\begin{align}
z&=c_B\T x_B+c_N\T x_N
 = c_B\T B^{-1}b+\bigl(c_N\T-c_B\T B^{-1}N\bigr)x_N. \label{eq:z}
\end{align}
\textbf{第三步（读出检验数）。}非基变量 $x_j$ 前的系数就是它的检验数；基变量在 \eqref{eq:z} 中已被消去，检验数为 0。两种情况可以统一写成一条公式（对基变量，$B^{-1}P_j$ 是单位向量 $e_r$，$c_j-c_B\T e_r=c_j-c_j=0$）。
\end{proofbox}

\begin{formulabox}[核心公式 1.1]
\[
\boxed{\;x_B=B^{-1}b,\qquad z=c_B\T B^{-1}b,\qquad \sigma\T=c\T-c_B\T B^{-1}A\;}\quad
\text{即}\quad\sigma_j=c_j-c_B\T B^{-1}P_j .
\]
\end{formulabox}

\begin{center}
\begin{tikzpicture}[x=1cm,y=1cm]
  \fill[lightblue] (0,0) rectangle (1.6,2.4);
  \fill[formulablue] (1.6,0) rectangle (8,2.4);
  \fill[lightgreen] (1.6,-0.8) rectangle (8,0);
  \fill[shiftc!20] (5.6,0) rectangle (8,2.4);
  \draw[navy] (0,-0.8) rectangle (8,2.4);
  \draw[navy] (1.6,-0.8)--(1.6,2.4);
  \draw[navy] (0,0)--(8,0);
  \node at (0.8,1.2) {$B^{-1}b$};
  \node at (3.6,1.2) {$B^{-1}A$};
  \node at (6.8,1.2) {$B^{-1}$};
  \node at (4.8,-0.4) {$\sigma\T=c\T-c_B\T B^{-1}A$};
  \node[font=\scriptsize,above] at (0.8,2.4) {$b$ 列};
  \node[font=\scriptsize,above] at (4.8,2.4) {$x_1\ \cdots\ x_n$（含松弛变量）};
  \node[font=\scriptsize,below,align=center,shiftc!80!black] at (6.8,-0.85) {初始单位阵所在的列};
  % 注释
  \draw[arr] (9.2,1.9) node[right,font=\scriptsize,align=left]{基变量取值\\$x_B=B^{-1}b$} -- (0.8,2.0);
  \draw[arr] (9.2,0.9) node[right,font=\scriptsize,align=left]{约束系数} -- (4.4,0.9);
  \draw[arr] (9.2,-0.4) node[right,font=\scriptsize,align=left]{检验数行} -- (7.6,-0.4);
  \draw[decorate,decoration={brace,amplitude=5pt,mirror},thick,auxblue] (-0.2,2.4)--(-0.2,0) node[midway,left=6pt,font=\scriptsize,align=right]{$m$ 行};
\end{tikzpicture}
\fig{标注公式图：任意一张单纯形表都可以由 $B^{-1}$ 和原始数据 $A,b,c$ 写出}
\end{center}

图中蓝色块是约束部分，绿色块是检验数行。最值得注意的是右侧橙色块：如果初始表中某几列恰好是单位阵（例如松弛变量所在的列），那么在任何一张后续表中，这几列都变成 $B^{-1}I=B^{-1}$。也就是说，\textbf{$B^{-1}$ 不用另算，它就印在当前表上}。这是课件反复强调的结论，后面读影子价格时还要用到。

\paragraph{符号说明。}
\begin{center}\small
\begin{tabularx}{\linewidth}{@{}>{\bfseries}l X l@{}}
\toprule
符号 & 含义 & 维数\\\midrule
$B$ & 当前基变量在原系数矩阵 $A$ 中的列，按\textbf{基变量在表中的行次序}排列 & $m\times m$\\
$P_j$ & $A$ 的第 $j$ 列，即变量 $x_j$ 在各约束中的系数 & $m\times1$\\
$c_B$ & 基变量的价值系数，排列次序与 $B$ 的列次序一致 & $m\times1$\\
$B^{-1}b$ & 当前基解中基变量的取值（$b$ 是\textbf{原模型}的右端项） & $m\times1$\\
$\sigma_j$ & 检验数：非基变量 $x_j$ 增加 1 单位时目标函数的变化率 & 数\\
\bottomrule
\end{tabularx}
\end{center}

\paragraph{使用条件。}公式要求 $B$ 可逆（$m$ 个基变量的列线性无关）。若 $B^{-1}b\ge0$，这组基解可行；若再有 $\sigma\le0$（max 问题），就是最优解。这两条判据在后面会被反复使用，第 5 章称之为“原始可行”与“对偶可行”。

\begin{warningbox}
\textbf{错误想法：}“$x_B=B^{-1}b$ 里的 $b$ 就是当前表的 $b$ 列。”\quad
\textbf{错在哪里：}当前表的 $b$ 列本身已经是 $B^{-1}b$，再乘一次就错了。\quad
\textbf{正确想法：}$b$ 永远是原模型右端项，$B^{-1}$ 从当前表的初始单位阵列读出。

\textbf{错误想法：}“$B$ 的列随便排。”\quad
\textbf{正确想法：}$B$ 的第 $r$ 列必须是第 $r$ 行基变量的系数列，$c_B$ 也按同样次序排；否则 $B^{-1}b$ 算出的数会对应到错误的变量上。
\end{warningbox}

\begin{center}
\begin{tikzpicture}[x=1cm,y=1cm]
  \node[gbox,text width=6.2cm,align=left] (ok) at (0,0) {\textbf{正确}：基按行次序 $\{x_3,x_1\}$\\[2pt]
  $B=(P_3,P_1)=\begin{pmatrix}2&1\\1&3\end{pmatrix}$，$c_B=(3,5)\T$\\[2pt]
  $B^{-1}b=(\tfrac{17}{5},\tfrac65)\T\Rightarrow x_3=\tfrac{17}{5},\ x_1=\tfrac65$};
  \node[rbox,text width=6.2cm,align=left] (bad) at (7.2,0) {\textbf{错误}：列次序与行次序不一致\\[2pt]
  $B=(P_1,P_3)=\begin{pmatrix}1&2\\3&1\end{pmatrix}$，$c_B=(3,5)\T$\\[2pt]
  $c_B$ 把价值 3 配给了 $x_1$，$z$ 与 $\sigma$ 全部算错};
\end{tikzpicture}
\fig{易错对照：$B$ 的列次序、$c_B$ 的次序、表中基变量的行次序，三者必须一致}
\end{center}

\subsection{修正单纯形法：只维护 $B^{-1}$}

核心公式 1.1 说明：只要知道 $B^{-1}$，表中任何一个数都能随用随算。而单纯形法每一轮真正\textbf{用到}的数只有三组：
\begin{enumerate}
  \item 检验数 $\sigma_j$（决定谁入基）；
  \item 入基列 $B^{-1}P_k$ 与右端 $B^{-1}b$（做比值、决定谁出基）；
  \item 新的 $B^{-1}$（下一轮用）。
\end{enumerate}
其余的列（非入基变量的 $B^{-1}P_j$）每轮都被重算，却从来没被看过。修正单纯形法就是“只算用得上的数”。

\begin{definitionbox}[算法 1.1　修正单纯形法（max 问题）\srcTag{笔记 P1--P4}]
给定一个可行基 $B$ 及其逆 $B^{-1}$：
\begin{enumerate}[label=\textbf{第\arabic*步}，,leftmargin=4.2em]
  \item 计算 $x_B=B^{-1}b$ 和单纯形乘子 $y\T=c_B\T B^{-1}$；
  \item 对每个非基变量计算 $\sigma_j=c_j-y\T P_j$。若全部 $\sigma_j\le0$，当前解最优，停止；
  \item 取 $\sigma_k$ 最大的 $x_k$ 入基，计算入基列 $\bar P_k=B^{-1}P_k$。若 $\bar P_k\le0$，问题无界，停止；
  \item 比值 $\theta=\min\{(B^{-1}b)_i/\bar a_{ik}:\bar a_{ik}>0\}$，取到最小值的第 $r$ 行基变量出基；
  \item 用 $P_k$ 替换 $B$ 的第 $r$ 列得新基，求新的 $B^{-1}$，回到第 1 步。
\end{enumerate}
\end{definitionbox}

\begin{center}
\begin{tikzpicture}[node distance=0.42cm,
  s/.style={box,minimum width=4.6cm,minimum height=0.8cm},
  d/.style={diamond,aspect=2.6,draw=navy,fill=white,font=\small,inner sep=1pt,align=center}]
  \node[s] (a) {已知可行基 $B$ 与 $B^{-1}$};
  \node[s,below=of a] (b) {$x_B=B^{-1}b$，\ $y\T=c_B\T B^{-1}$};
  \node[s,below=of b] (c) {$\sigma_j=c_j-y\T P_j$（只算非基列）};
  \node[d,below=of c] (d) {全部\\$\sigma_j\le0$?};
  \node[gbox,right=1.2cm of d] (opt) {最优，停止};
  \node[s,below=of d] (e) {$x_k$ 入基，算 $\bar P_k=B^{-1}P_k$};
  \node[s,below=of e] (f) {$\theta$ 比值定出基行 $r$};
  \node[s,below=of f] (g) {换列得新 $B$，更新 $B^{-1}$};
  \draw[arr] (a)--(b); \draw[arr] (b)--(c); \draw[arr] (c)--(d);
  \draw[arr] (d)--node[above,font=\scriptsize]{是}(opt);
  \draw[arr] (d)--node[right,font=\scriptsize]{否}(e);
  \draw[arr] (e)--(f); \draw[arr] (f)--(g);
  \draw[arr] (g.west) -- ++(-1.2,0) |- (b.west);
  \node[font=\scriptsize,auxblue,align=left,right] at (3.2,-6.1) {每轮只算：\\一行 $y\T$、若干个数 $\sigma_j$、\\一列 $\bar P_k$、一个 $B^{-1}$};
\end{tikzpicture}
\fig{修正单纯形法流程图：与普通单纯形法的判据完全相同，只是不再维护整张表}
\end{center}

流程图左侧的回路就是一轮迭代。与普通单纯形法相比，入基、出基的\textbf{判据一字未改}，改变的只是“算哪些数”。笔记 P2 页眉写的“计算机：用特定的矩阵乘法运算（微观）”，指的就是这一点：计算机程序里实际执行的是修正单纯形法。

\begin{supplementbox}[$2\times2$ 逆矩阵的快速求法 \fillTag]
\begin{minipage}{0.52\linewidth}
本章例题与练习常遇到 $2\times2$ 基矩阵。若 $B=\begin{pmatrix}a&b\\c&d\end{pmatrix}$ 且 $ad-bc\ne0$，则
\[
B^{-1}=\frac{1}{ad-bc}\begin{pmatrix}d&-b\\-c&a\end{pmatrix}.
\]
口诀：\textbf{主对角互换，副对角变号，再除以行列式}。求完务必用 $BB^{-1}=I$ 检验——纠错 2 就是漏了这一步。
\end{minipage}\hfill
\begin{minipage}{0.44\linewidth}\centering
\begin{tikzpicture}[x=0.75cm,y=0.75cm]
  \node at (0,0) {$\begin{pmatrix}2&1\\1&3\end{pmatrix}$};
  \draw[{Stealth}-{Stealth},shiftc,thick] (-0.45,0.35) to[bend left=40] (0.45,-0.35);
  \node[font=\scriptsize,shiftc!80!black] at (0,1.0) {2 与 3 互换};
  \draw[arr] (1.3,0)--(2.4,0);
  \node at (4.2,0) {$\frac15\begin{pmatrix}3&-1\\-1&2\end{pmatrix}$};
  \node[font=\scriptsize,auxblue] at (4.2,-1.0) {$\det=2\cdot3-1\cdot1=5$};
  \node[font=\scriptsize,redline] at (4.4,1.0) {副对角取负};
\end{tikzpicture}
\fig{$2\times2$ 求逆的三步操作}
\end{minipage}
\end{supplementbox}

\subsection{例题：同一道题，两种算法}

\begin{examplebox}[例 1.1　\srcTag{笔记 P2--P4}]
用修正单纯形法求解
\[
\max z=5x_1+2x_2+3x_3-x_4+x_5,\qquad
\begin{cases}x_1+2x_2+2x_3+x_4=8,\\ 3x_1+4x_2+x_3+x_5=7,\\ x_j\ge0.\end{cases}
\]
\textbf{已知：}$A=(P_1,\dots,P_5)=\begin{pmatrix}1&2&2&1&0\\3&4&1&0&1\end{pmatrix}$，$b=\binom87$，$c=(5,2,3,-1,1)$。\quad
\textbf{求：}最优解与最优值。\par
\textbf{思路：}$x_4,x_5$ 的列恰为单位阵，取 $B^{(1)}=(P_4,P_5)=I$ 作初始基，$c_B=(-1,1)\T$。注意 $x_4,x_5$ 在目标函数中有非零系数，所以 $c_B\ne0$，初始检验数不等于 $c_j$。

\textbf{第一轮。}$B^{-1}=I$，$x_B=b=(8,7)\T$，$y\T=c_B\T B^{-1}=(-1,1)$。
\[
\sigma_1=5-(-1\cdot1+1\cdot3)=3,\quad \sigma_2=2-(-2+4)=0,\quad \sigma_3=3-(-2+1)=4 .
\]
$\sigma_3=4$ 最大，$x_3$ 入基；$\bar P_3=B^{-1}P_3=(2,1)\T$，$\theta=\min\{\frac82,\frac71\}=4$，第 1 行的 $x_4$ 出基。

\textbf{第二轮。}新基 $\{x_3,x_5\}$：$B=(P_3,P_5)=\begin{pmatrix}2&0\\1&1\end{pmatrix}$，$B^{-1}=\begin{pmatrix}\frac12&0\\-\frac12&1\end{pmatrix}$，$x_B=B^{-1}b=(4,3)\T$，$c_B=(3,1)\T$，
$y\T=(3,1)B^{-1}=(1,1)$。对非基变量 $x_1,x_2,x_4$：
\[
\sigma_1=5-(1+3)=1,\quad \sigma_2=2-(2+4)=-4,\quad \sigma_4=-1-(1+0)=-2 .
\]
$x_1$ 入基；$\bar P_1=B^{-1}P_1=(\frac12,\frac52)\T$，$\theta=\min\{\frac{4}{1/2},\frac{3}{5/2}\}=\min\{8,\frac65\}=\frac65$，$x_5$ 出基。

\textbf{第三轮。}新基 $\{x_3,x_1\}$：$B=(P_3,P_1)=\begin{pmatrix}2&1\\1&3\end{pmatrix}$，$B^{-1}=\frac15\begin{pmatrix}3&-1\\-1&2\end{pmatrix}$，$x_B=(\frac{17}{5},\frac65)\T$，$c_B=(3,5)\T$，$y\T=(\frac45,\frac75)$：
\[
\sigma_2=2-\bigl(\tfrac85+\tfrac{28}{5}\bigr)=-\tfrac{26}{5},\quad
\sigma_4=-1-\tfrac45=-\tfrac95,\quad \sigma_5=1-\tfrac75=-\tfrac25 .
\]
全部 $\le0$，最优。\quad\textbf{结果：}$x^*=(\frac65,0,\frac{17}{5},0,0)\T$，$z^*=5\cdot\frac65+3\cdot\frac{17}{5}=\frac{81}{5}$。\par
\textbf{检查：}代回约束：$\frac65+\frac{34}{5}=8$ ✓，$\frac{18}{5}+\frac{17}{5}=7$ ✓。\par
\textbf{本题真正想说明什么：}三轮中我们一次也没有写出完整的 $2\times5$ 表，每轮只算了一个 $2\times2$ 逆矩阵、一行 $y\T$ 和一列 $\bar P_k$。
\end{examplebox}

\begin{center}
\begin{tikzpicture}[x=1cm,y=1cm]
  \foreach \i/\bs/\z/\sig in {0/{\{x_4,x_5\}}/-1/{(3,0,4)},1/{\{x_3,x_5\}}/15/{(1,-4,-2)},2/{\{x_3,x_1\}}/{\tfrac{81}{5}}/{(-\tfrac{26}{5},-\tfrac95,-\tfrac25)}}{
    \node[box,minimum width=3.6cm,minimum height=1.35cm,align=center] (n\i) at (\i*4.8,0) {基 $\bs$\\$z=\z$\\[1pt]{\scriptsize 非基 $\sigma=\sig$}};
  }
  \draw[arr] (n0)--node[above,font=\scriptsize]{$x_3$ 入，$x_4$ 出}node[below,font=\scriptsize]{$\theta=4$}(n1);
  \draw[arr] (n1)--node[above,font=\scriptsize]{$x_1$ 入，$x_5$ 出}node[below,font=\scriptsize]{$\theta=\tfrac65$}(n2);
  % 目标值柱
  \begin{scope}[yshift=-3.3cm]
    \draw[ax] (-1.8,0)--(11.6,0);
    \foreach \i/\h/\t in {0/-0.08/{-1},1/1.2/{15},2/1.296/{16.2}}{
      \fill[auxblue!70] (\i*4.8-0.4,0) rectangle (\i*4.8+0.4,\h);
      \node[font=\scriptsize,above] at (\i*4.8,{max(\h,0)}) {$z=\t$};
    }
    \node[font=\scriptsize,left] at (-1.8,0.6) {目标值};
  \end{scope}
\end{tikzpicture}
\fig{例 1.1 的迭代路径：三个基、两次换基，目标值单调上升直至检验数全部非正}
\end{center}

上图上半部分把三轮迭代压缩成三个“状态盒”，箭头上写出换基依据；下半部分的柱高表示每个基对应的目标值。读图可以看出：每换一次基，目标值就上升一次（$-1\to15\to16.2$），而换基的理由都来自检验数。笔记 P2 右栏用普通单纯形表算了同一道题，两种方法每一轮的基、检验数完全相同（见纠错 1）。

\begin{lifebox}[查菜谱只翻用得到的那一页]
\begin{minipage}[t]{0.5\linewidth}
\textbf{生活场景：}做一道新菜时，你不会把整本菜谱抄一遍，而是先查目录，再只翻到需要的那一页。\par
\textbf{对应关系：}整本菜谱 $\leftrightarrow$ 完整单纯形表；目录 $\leftrightarrow$ $B^{-1}$；翻到的那一页 $\leftrightarrow$ 入基列 $B^{-1}P_k$。\par
\textbf{小估算：}例 1.1 有 $m=2$ 行、$n=5$ 列。普通单纯形法每轮更新 $2\times(5+1)=12$ 个数；修正单纯形法每轮只需 $B^{-1}$（4 个）、$\bar P_k$（2 个）、$x_B$（2 个），共 8 个。约束越多、变量越多，差距越大。\par
\textbf{边界说明：}实际软件并不显式存储 $B^{-1}$，而是存储它的分解形式（如 LU 分解），以减少舍入误差；“只算用得到的数”的思想不变。
\end{minipage}\hfill
\begin{minipage}[t]{0.46\linewidth}\vspace{0pt}\centering
\begin{tikzpicture}[x=1cm,y=1cm]
  % 书
  \draw[warmline,fill=white] (0,0) rectangle (1.6,2.2);
  \foreach \y in {0.4,0.7,...,1.9} \draw[black!25] (0.2,\y)--(1.4,\y);
  \node[font=\scriptsize] at (0.8,2.45) {整本菜谱};
  \fill[shiftc!30] (0.15,1.25) rectangle (1.45,1.45);
  \node[font=\scriptsize] at (0.8,-0.3) {目录 + 一页};
  % 表
  \begin{scope}[xshift=2.9cm]
    \draw[navy] (0,0) grid[step=0.36] (1.8,2.16);
    \fill[shiftc!30] (1.08,0) rectangle (1.44,2.16);
    \node[font=\scriptsize] at (0.9,2.45) {单纯形表};
    \node[font=\scriptsize] at (0.9,-0.3) {$B^{-1}$ + 入基列};
  \end{scope}
  \draw[{Stealth}-{Stealth},auxblue,thick] (1.7,1.35)--(3.95,1.1);
  \draw[{Stealth}-{Stealth},auxblue,thick,dashed] (1.7,0.5)--(2.85,0.5);
\end{tikzpicture}
\fig{生活场景 $\leftrightarrow$ 抽象模型}
\end{minipage}
\end{lifebox}

\begin{checkbox}
\begin{enumerate}[leftmargin=1.5em]
  \item 为什么说单纯形法的一轮旋转等于左乘一个矩阵？
  \item 写出 $x_B$、$z$、$\sigma$ 用 $B^{-1}$ 表示的三个公式，并说明 $b$ 指的是哪一个 $b$。
  \item 在当前单纯形表上，$B^{-1}$ 印在哪几列？为什么？
  \item 修正单纯形法每一轮必须算的是哪几组数？
  \item 若基变量按行次序是 $\{x_3,x_1\}$，$B$ 和 $c_B$ 应怎样排列？
\end{enumerate}
\end{checkbox}

\paragraph{到这里，我们已经解决了什么？}我们把单纯形表“拆开”了：表中每个数都由 $B^{-1}$ 和原始数据决定，检验数可以写成 $\sigma_j=c_j-y\T P_j$，其中 $y\T=c_B\T B^{-1}$。
\paragraph{下一步还需要解决什么？}这个 $y$ 只是计算中的一个中间量，还是有独立含义？下一章我们将看到：$y$ 恰好是另一个线性规划——对偶问题——的解，而且它有非常直观的经济含义：资源的价格。

\begin{center}
\begin{tikzpicture}[node distance=0.35cm and 0.35cm,k/.style={box,minimum height=0.8cm,font=\small}]
  \node[k] (a) {初等行变换};
  \node[k,right=of a] (b) {左乘 $E$};
  \node[k,right=of b] (c) {一轮旋转 $=$ 左乘 $B^{-1}$};
  \node[k,right=of c] (d) {表 $=B^{-1}[\,b\mid A\,]$};
  \node[k,below=0.6cm of c] (e) {$\sigma=c-c_B\T B^{-1}A$};
  \node[k,below=0.6cm of d] (f) {修正单纯形法};
  \node[gbox,below=0.6cm of b] (g) {$y\T=c_B\T B^{-1}$\\$\to$ 第 2 章};
  \draw[arr] (a)--(b); \draw[arr] (b)--(c); \draw[arr] (c)--(d);
  \draw[arr] (d)--(f); \draw[arr] (c)--(e); \draw[arr] (e)--(f); \draw[arr] (e)--(g);
\end{tikzpicture}
\fig{本章知识地图}
\end{center}

% ============================================================
%  第 2 章  对偶问题的构造（K2）
% ============================================================
\section{对偶问题：同一组数据的另一种读法}

上一章结尾留下了一个问题：计算检验数时出现的 $y\T=c_B\T B^{-1}$ 有没有独立的含义？本章先把它放一放，从一个经济学问题出发：如果有人想把工厂的资源整体买走，他应该出什么价？我们会发现，这个“买方定价问题”本身也是一个线性规划，它和原来的生产问题用的是同一组数字，只是读法不同。

\noindent\textbf{本章回答三个问题：}(1) 对偶问题是怎样从实际问题中“长”出来的？(2) 对称形式的对偶怎样写？(3) 约束带 $\ge$、$=$，变量带 $\le0$、无约束时，对偶怎样写？\par
\noindent\textbf{本章知识路线：}生产者模型 $\to$ 购买者模型 $\to$ 对称形式 $\to$ 一般形式的对应规则 $\to$ 例题。

\subsection{问题引入：一家工厂的两种算账方式}

\begin{examplebox}[引例　生产计划问题（原问题 $P$）\srcTag{笔记 P5--P6}]
\begin{minipage}[t]{0.5\linewidth}
某工厂生产甲、乙两种产品，消耗煤、电、油三种资源，数据如下表。求总收入最大的生产方案。
\begin{center}\small
\begin{tabular}{@{}lccc@{}}
\toprule
资源 & 甲 & 乙 & 资源限制\\\midrule
煤 & 9 & 4 & 360\\
电 & 4 & 5 & 200\\
油 & 3 & 10 & 300\\\midrule
单价 & 7 & 12 & \\\bottomrule
\end{tabular}
\end{center}
设生产甲 $x_1$、乙 $x_2$ 单位：
\[
\begin{aligned}
\max z&=7x_1+12x_2\\
\st &9x_1+4x_2\le360\quad(\text{煤})\\
&4x_1+5x_2\le200\quad(\text{电})\\
&3x_1+10x_2\le300\quad(\text{油})\\
&x_1,x_2\ge0
\end{aligned}
\]
单纯形法解得 $x^*=(20,24)$，$z^*=7\times20+12\times24=428$。
\end{minipage}\hfill
\begin{minipage}[t]{0.46\linewidth}\vspace{0pt}\centering
\begin{tikzpicture}[x=0.95cm,y=0.95cm]
  \fill[feasreg] (0,0)--(4,0)--(3.448,1.241)--(2,2.4)--(0,3)--cycle;
  \draw[gridl] (0,0) grid (6,5);
  \draw[ax] (-0.2,0)--(6.3,0) node[right,font=\scriptsize]{$x_1$};
  \draw[ax] (0,-0.2)--(0,5.3) node[above,font=\scriptsize]{$x_2$};
  \foreach \i in {1,...,6} \node[below,font=\tiny,black!60] at (\i,0) {\the\numexpr\i*10\relax};
  \foreach \i in {1,...,5} \node[left,font=\tiny,black!60] at (0,\i) {\the\numexpr\i*10\relax};
  \draw[con] (4,0)--(1.78,5) node[above,font=\scriptsize]{煤};
  \draw[con] (5,0)--(0,4) node[left,font=\scriptsize]{电};
  \draw[con] (0,3)--(6,1.2) node[right,font=\scriptsize]{油};
  \draw[iso] (0,3.567)--(5.1,0.592);
  \node[optpt] at (2,2.4) {};
  \node[lab,above right] at (2,2.45) {$(20,24)$};
  \node[font=\scriptsize,black!55,rotate=-30] at (4.4,1.3) {$z=428$};
\end{tikzpicture}
\fig{原问题可行域：电、油两条约束在最优点 $(20,24)$ 交汇，煤约束不经过该点}
\end{minipage}
\end{examplebox}

图中绿色区域是可行域，虚线是等值线 $7x_1+12x_2=428$。最优点由“电”和“油”两条边界线交出；“煤”的边界离最优点还有一段距离，说明煤没有用完：$9\times20+4\times24=276<360$，剩余 84（这就是终表中 $x_3=84$ 的来历）。

\paragraph{现在换一个角度。}\srcTag{笔记 P5}\ 另一家厂想把这家工厂的煤、电、油\textbf{全部买走}。它要给三种资源定单价 $y_1,y_2,y_3$。它会怎样思考？
\begin{enumerate}
  \item \textbf{买方希望总花费最少：}$\min w=360y_1+200y_2+300y_3$；
  \item \textbf{但卖方必须愿意卖：}对每一种产品，资源按新价格计算的“售价”不能低于自己生产这种产品的收入，否则卖方宁可自己生产。
\end{enumerate}

\begin{center}
\begin{tikzpicture}[x=1cm,y=1cm]
  \node[wbox,minimum width=3.4cm] (jia) at (0,0) {生产 1 件甲\\需煤 9、电 4、油 3};
  \node[box,minimum width=3.6cm] (sell) at (5.6,0.8) {卖掉这些资源得\\$9y_1+4y_2+3y_3$};
  \node[box,minimum width=3.6cm] (make) at (5.6,-0.8) {自己生产甲得\\收入 $7$};
  \draw[arr] (jia.east)--(sell.west);
  \draw[arr] (jia.east)--(make.west);
  \node[gbox,minimum width=3.6cm] (c) at (11.0,0) {卖方愿意卖的条件\\$9y_1+4y_2+3y_3\ge7$};
  \draw[arr] (sell.east)--(c.west); \draw[arr] (make.east)--(c.west);
  \node[font=\scriptsize,auxblue] at (8.3,1.5) {比较};
\end{tikzpicture}
\fig{购买者模型中一条约束的来历：站在产品甲的角度比较“卖资源”与“自己生产”}
\end{center}

这幅图说明了对偶约束的经济来源：产品甲占用的那一篮子资源（9 份煤、4 份电、3 份油），按买方给的价格折算后，不能低于甲的售价 7。产品乙同理得 $4y_1+5y_2+10y_3\ge12$。于是得到\textbf{对偶问题 $D$}：
\[
\begin{aligned}
\min w&=360y_1+200y_2+300y_3\\
\st &9y_1+4y_2+3y_3\ge7\quad(\text{甲})\\
&4y_1+5y_2+10y_3\ge12\quad(\text{乙})\\
&y_1,y_2,y_3\ge0
\end{aligned}
\]
笔记中的红字概括为一句话：\textbf{生产 $x_1$（$x_2$）甲（乙）产品的原料售价，不低于自己用这些原料生产所获得的利益。}

\subsection{对称形式的对偶问题}

把原问题和对偶问题并排写成矩阵形式，规律一目了然。

\begin{definitionbox}[定义 2.1　对称形式的对偶问题 \srcTag{笔记 P6--P7}]
\[
\text{(P)}\quad
\begin{aligned}
\max z&=c\T x\\ \st &Ax\le b\\ &x\ge0
\end{aligned}
\qquad\Longleftrightarrow\qquad
\text{(D)}\quad
\begin{aligned}
\min w&=b\T y\\ \st &A\T y\ge c\\ &y\ge0
\end{aligned}
\]
其中 $A$ 为 $m\times n$ 矩阵，$x\in\mathbb R^n$，$y\in\mathbb R^m$。两者互称对偶问题。
\end{definitionbox}

\begin{center}
\begin{tikzpicture}[x=0.9cm,y=0.9cm]
  % 原问题数据块
  \fill[formulablue] (0,0) rectangle (3,2.4); \node at (1.5,1.2) {$A$};
  \fill[warm] (3.2,0) rectangle (3.8,2.4); \node at (3.5,1.2) {$b$};
  \fill[lightgreen] (0,2.6) rectangle (3,3.2); \node at (1.5,2.9) {$c\T$};
  \draw[navy] (0,0) rectangle (3,2.4); \draw[navy] (3.2,0) rectangle (3.8,2.4); \draw[navy] (0,2.6) rectangle (3,3.2);
  \node[font=\scriptsize] at (1.5,-0.35) {$n$ 列 $=n$ 个变量 $x_j$};
  \node[font=\scriptsize,rotate=90] at (-0.35,1.2) {$m$ 行 $=m$ 条约束};
  \node[font=\small\bfseries,navy] at (1.9,3.75) {原问题（max）};
  % 箭头
  \draw[-{Stealth[length=3mm]},very thick,shiftc] (4.5,1.6)--(6.3,1.6) node[midway,above,font=\scriptsize,shiftc!80!black]{整体转置};
  % 对偶数据块
  \begin{scope}[xshift=7cm]
  \fill[formulablue] (0,0) rectangle (2.4,3); \node at (1.2,1.5) {$A\T$};
  \fill[lightgreen] (2.6,0) rectangle (3.2,3); \node at (2.9,1.5) {$c$};
  \fill[warm] (0,3.2) rectangle (2.4,3.8); \node at (1.2,3.5) {$b\T$};
  \draw[navy] (0,0) rectangle (2.4,3); \draw[navy] (2.6,0) rectangle (3.2,3); \draw[navy] (0,3.2) rectangle (2.4,3.8);
  \node[font=\scriptsize] at (1.2,-0.35) {$m$ 列 $=m$ 个变量 $y_i$};
  \node[font=\scriptsize,rotate=90] at (-0.35,1.5) {$n$ 行 $=n$ 条约束};
  \node[font=\small\bfseries,navy] at (1.6,4.35) {对偶问题（min）};
  \end{scope}
  \node[font=\scriptsize,align=left,auxblue] at (12.5,2.6) {右端项 $b$\\$\to$ 目标系数};
  \node[font=\scriptsize,align=left,auxblue] at (12.5,1.0) {目标系数 $c$\\$\to$ 右端项};
\end{tikzpicture}
\fig{对称形式的对偶 $=$ 把“数据矩阵”整体转置：$A\to A\T$，$b$ 与 $c$ 交换位置，max 变 min，$\le$ 变 $\ge$}
\end{center}

图中三种颜色分别标出系数矩阵、右端项、价值系数。转置之后：原问题的\textbf{每一条约束}变成对偶的\textbf{一个变量}（行 $\to$ 列），原问题的\textbf{每一个变量}变成对偶的\textbf{一条约束}。笔记把它类比为“矩阵转置”，非常准确。

\paragraph{为什么要这样定义？}从引例可见：对偶变量 $y_i$ 是第 $i$ 种资源的单价，所以有 $m$ 个；对偶约束逐个产品比较“资源售价 $\ge$ 产品收入”，所以有 $n$ 条，系数正好是 $A$ 的列，即 $A\T$ 的行。

\begin{proofbox}[对称性：对偶的对偶就是原问题 \fillTag]
把 (D) 改写成 max 型与 $\le$ 型：$\max\,(-b)\T y$，$\st (-A\T)y\le -c$，$y\ge0$。按定义 2.1 写它的对偶（数据为 $-A\T,\,-c,\,-b$）：
$\min\,(-c)\T x$，$\st (-A\T)\T x\ge -b$，$x\ge0$，即 $\max c\T x$，$\st Ax\le b$，$x\ge0$。这正是 (P)。\quad 所以笔记写“$P$ 与 $D$ 互为对偶（类比矩阵转置 $(A\T)\T=A$）”。
\end{proofbox}

\subsection{一般形式：约束和变量不“规矩”时怎么办}

实际模型里常常有 $\ge$ 约束、等式约束、$x_j\le0$ 或无符号限制的变量。笔记把对应关系整理成一张表（P7--P8），我们先给出规则，再解释为什么。

\begin{theorembox}[对应规则 2.2　原问题（max）$\longleftrightarrow$ 对偶问题（min）\srcTag{笔记 P7--P8}]
\begin{center}
\begin{tabular}{@{}l@{\qquad}c@{\qquad}l@{}}
\toprule
\textbf{原问题（max）} & & \textbf{对偶问题（min）}\\\midrule
目标函数 $\max z$ & $\longleftrightarrow$ & 目标函数 $\min w$\\
约束条件右端项 $b_i$ & $\longleftrightarrow$ & 目标函数系数\\
目标函数系数 $c_j$ & $\longleftrightarrow$ & 约束条件右端项\\\midrule
第 $i$ 个约束 $\le b_i$ & $\longleftrightarrow$ & 第 $i$ 个变量 $y_i\ge0$\\
第 $i$ 个约束 $\ge b_i$ & $\longleftrightarrow$ & 第 $i$ 个变量 $y_i\le0$\\
第 $i$ 个约束 $=b_i$ & $\longleftrightarrow$ & 第 $i$ 个变量 $y_i$ 无约束\\\midrule
第 $j$ 个变量 $x_j\ge0$ & $\longleftrightarrow$ & 第 $j$ 个约束 $\ge c_j$\\
第 $j$ 个变量 $x_j\le0$ & $\longleftrightarrow$ & 第 $j$ 个约束 $\le c_j$\\
第 $j$ 个变量 $x_j$ 无约束 & $\longleftrightarrow$ & 第 $j$ 个约束 $=c_j$\\\bottomrule
\end{tabular}
\end{center}
原问题是 min 时，把表\textbf{从右往左}读（笔记：“原 $\min z$，则对偶 $\max$（右往左看）；原 $\max z$，则对偶 $\min$（左 $\to$ 右）”）。
\end{theorembox}

\begin{center}
\begin{tikzpicture}[x=1cm,y=1cm,
  n/.style={draw=navy,rounded corners=2pt,fill=white,minimum width=2.6cm,minimum height=0.62cm,font=\small},
  g/.style={draw=greenline,rounded corners=2pt,fill=lightgreen,minimum width=2.6cm,minimum height=0.62cm,font=\small},
  r/.style={draw=redline!70,rounded corners=2pt,fill=lightred,minimum width=2.6cm,minimum height=0.62cm,font=\small},
  e/.style={draw=auxblue,rounded corners=2pt,fill=grayblue,minimum width=2.6cm,minimum height=0.62cm,font=\small}]
  \node[font=\small\bfseries,navy] at (0,1.0) {max 问题的约束};
  \node[font=\small\bfseries,navy] at (4.6,1.0) {min 问题的变量};
  \node[g] (a1) at (0,0) {$\le$（正常）}; \node[g] (b1) at (4.6,0) {$y_i\ge0$（正常）};
  \node[r] (a2) at (0,-0.9) {$\ge$（反常）}; \node[r] (b2) at (4.6,-0.9) {$y_i\le0$（反常）};
  \node[e] (a3) at (0,-1.8) {$=$}; \node[e] (b3) at (4.6,-1.8) {无约束};
  \draw[arr] (a1)--(b1); \draw[arr] (a2)--(b2); \draw[arr] (a3)--(b3);
  \begin{scope}[xshift=8.4cm]
  \node[font=\small\bfseries,navy] at (0,1.0) {max 问题的变量};
  \node[font=\small\bfseries,navy] at (4.6,1.0) {min 问题的约束};
  \node[g] (c1) at (0,0) {$x_j\ge0$（正常）}; \node[g] (d1) at (4.6,0) {$\ge c_j$（正常）};
  \node[r] (c2) at (0,-0.9) {$x_j\le0$（反常）}; \node[r] (d2) at (4.6,-0.9) {$\le c_j$（反常）};
  \node[e] (c3) at (0,-1.8) {无约束}; \node[e] (d3) at (4.6,-1.8) {$=c_j$};
  \draw[arr] (c1)--(d1); \draw[arr] (c2)--(d2); \draw[arr] (c3)--(d3);
  \end{scope}
\end{tikzpicture}
\fig{记忆图：“正常对正常、反常对反常、等号对无约束”。max 问题的正常约束是 $\le$，min 问题的正常约束是 $\ge$，正常变量都是 $\ge0$}
\end{center}

图中绿色是“正常”情形，即对称形式里出现的情形；红色是“反常”情形；蓝灰色是等号与无约束的配对。只要记住对称形式（定义 2.1），其余六种情形都能由“正常对正常、反常对反常”推出，不必死记整张表。

\begin{proofbox}[为什么是这样对应？\fillTag]
思路：把一般形式先化成对称形式，用定义 2.1 写出对偶，再把变量换回来。
\begin{enumerate}[leftmargin=1.5em]
  \item \textbf{约束 $a_i\T x\ge b_i$（反常）。}两边乘 $-1$ 化为正常约束 $-a_i\T x\le-b_i$，它对应的对偶变量 $y_i'\ge0$，在对偶中贡献系数 $-a_i$ 与目标项 $-b_iy_i'$。令 $y_i=-y_i'$，贡献变回 $a_i$ 与 $b_iy_i$，而 $y_i=-y_i'\le0$。
  \item \textbf{约束 $a_i\T x=b_i$。}拆成 $a_i\T x\le b_i$ 与 $a_i\T x\ge b_i$，分别对应 $y_i'\ge0$ 与 $y_i''\le0$；二者在对偶中系数相同，可合并为 $y_i=y_i'+y_i''$，它可正可负，即无约束。
  \item \textbf{变量 $x_j\le0$（反常）。}令 $x_j=-x_j'$，$x_j'\ge0$，则这一列变为 $-P_j$、价值系数变为 $-c_j$。对偶约束 $-P_j\T y\ge-c_j$ 等价于 $P_j\T y\le c_j$。
  \item \textbf{变量 $x_j$ 无约束。}令 $x_j=x_j'-x_j''$，两者都 $\ge0$，得到两条对偶约束 $P_j\T y\ge c_j$ 与 $-P_j\T y\ge-c_j$，合起来就是 $P_j\T y=c_j$。
\end{enumerate}
\textbf{一句话直观解释：}“反常”就是对称形式乘了一个 $-1$，这个 $-1$ 会原封不动地传到对偶的另一侧。
\end{proofbox}

\subsection{例题：按四步写出对偶}

笔记把写对偶的过程总结为四步：\textbf{1° 明确目标}（max 还是 min）；\textbf{2° 确定查表方向}（max 左 $\to$ 右，min 右 $\to$ 左）；\textbf{3° 变量 $\to$ 约束}；\textbf{4° 约束 $\to$ 变量}。

\begin{examplebox}[例 2.1　\srcTag{笔记 P8--P9}]
写出下列问题的对偶问题：
\[
\max z=2x_1+3x_2-5x_3+x_4,\quad
\begin{cases}
x_1+x_2-3x_3+x_4\ge5\\
2x_1+2x_3-x_4\le4\\
x_2+x_3+x_4=6\\
x_1\le0,\ x_2,x_3\ge0,\ x_4\text{ 无约束}
\end{cases}
\]
\textbf{1° 目标：}原问题 max $\Rightarrow$ 对偶 $\min w=5y_1+4y_2+6y_3$（右端项 $5,4,6$ 成为目标系数）。\par
\textbf{2° 方向：}max 问题，规则表从左往右读。\par
\textbf{3° 变量 $\to$ 约束：}系数矩阵转置后，第 $j$ 列给出第 $j$ 条对偶约束，右端为 $c_j$：
\[
A=\begin{pmatrix}1&1&-3&1\\2&0&2&-1\\0&1&1&1\end{pmatrix},\qquad
\begin{array}{lll}
x_1\le0\ (\text{反常}) &\Rightarrow& y_1+2y_2\le2\\
x_2\ge0\ (\text{正常}) &\Rightarrow& y_1+y_3\ge3\\
x_3\ge0\ (\text{正常}) &\Rightarrow& -3y_1+2y_2+y_3\ge-5\\
x_4\text{ 无约束} &\Rightarrow& y_1-y_2+y_3=1
\end{array}
\]
\textbf{4° 约束 $\to$ 变量：}第 1 个约束 $\ge$（反常）$\Rightarrow y_1\le0$；第 2 个 $\le$（正常）$\Rightarrow y_2\ge0$；第 3 个 $=\Rightarrow y_3$ 无约束。\par
\textbf{结果：}
\[
\min w=5y_1+4y_2+6y_3,\quad
\begin{cases}
y_1+2y_2\le2\\ y_1+y_3\ge3\\ -3y_1+2y_2+y_3\ge-5\\ y_1-y_2+y_3=1\\ y_1\le0,\ y_2\ge0,\ y_3\text{ 无约束}
\end{cases}
\]
\textbf{检查：}原问题 4 个变量、3 个约束 $\Rightarrow$ 对偶 4 个约束、3 个变量 ✓。
\end{examplebox}

\begin{center}
\begin{tikzpicture}[x=1cm,y=1cm,
  pv/.style={draw=navy,fill=lightblue,rounded corners=2pt,minimum width=1.9cm,minimum height=0.55cm,font=\small},
  dv/.style={draw=greenline,fill=lightgreen,rounded corners=2pt,minimum width=3.6cm,minimum height=0.55cm,font=\small}]
  \node[font=\small\bfseries,navy] at (0,0.75) {原问题的变量};
  \node[font=\small\bfseries,greenline] at (4.2,0.75) {对偶的约束};
  \foreach \i/\a/\b in {0/{x_1\le0}/{y_1+2y_2\le2},1/{x_2\ge0}/{y_1+y_3\ge3},2/{x_3\ge0}/{\cdots\ge-5},3/{x_4\text{ 自由}}/{y_1-y_2+y_3=1}}{
    \node[pv] (p\i) at (0,-\i*0.8) {$\a$};
    \node[dv] (d\i) at (4.2,-\i*0.8) {$\b$};
    \draw[arr] (p\i)--(d\i);
  }
  \begin{scope}[xshift=8.6cm]
  \node[font=\small\bfseries,navy] at (0,0.75) {原问题的约束};
  \node[font=\small\bfseries,greenline] at (3.6,0.75) {对偶的变量};
  \foreach \i/\a/\b in {0/{\cdots\ge5}/{y_1\le0},1/{\cdots\le4}/{y_2\ge0},2/{\cdots=6}/{y_3\text{ 自由}}}{
    \node[pv,minimum width=2.1cm] (q\i) at (0,-\i*0.8) {$\a$};
    \node[dv,minimum width=2.1cm] (e\i) at (3.6,-\i*0.8) {$\b$};
    \draw[arr] (q\i)--(e\i);
  }
  \end{scope}
\end{tikzpicture}
\fig{例 2.1 的推导路线图：左半是第 3 步（变量 $\to$ 约束），右半是第 4 步（约束 $\to$ 变量）}
\end{center}

图中每一根箭头都是规则表的一行。这样逐项连线，可以避免漏写或把符号配错。

\begin{examplebox}[例 2.2　同样的数据，目标改为 min　\srcTag{笔记 P9--P11}]
把例 2.1 的目标改成 $\min z=2x_1+3x_2-5x_3+x_4$，约束不变。\par
\textbf{1°} 原问题 min $\Rightarrow$ 对偶 $\max w=5y_1+4y_2+6y_3$。\quad
\textbf{2°} 规则表\textbf{从右往左}读：此时原问题处在表的右栏。\par
\textbf{3° 变量 $\to$ 约束}（右栏“变量”对左栏“约束”）：$x_1\le0\Rightarrow$ 对偶约束 $\ge$；$x_2,x_3\ge0\Rightarrow\le$；$x_4$ 无约束 $\Rightarrow=$。\par
\textbf{4° 约束 $\to$ 变量}（右栏“约束”对左栏“变量”）：$\ge5\Rightarrow y_1\ge0$；$\le4\Rightarrow y_2\le0$；$=6\Rightarrow y_3$ 无约束。\par
\textbf{结果：}
$\max w=5y_1+4y_2+6y_3$，\ $y_1+2y_2\ge2$，\ $y_1+y_3\le3$，\ $-3y_1+2y_2+y_3\le-5$，\ $y_1-y_2+y_3=1$，\ $y_1\ge0,\ y_2\le0,\ y_3$ 无约束。\par
\textbf{本题真正想说明什么：}与例 2.1 比较，所有不等号和符号限制都\textbf{恰好反过来}，而等式与无约束不变。原因是：对 min 问题，“正常”的约束是 $\ge$。
\end{examplebox}

\begin{center}
\begin{tikzpicture}[x=1cm,y=1cm]
  \node[box,text width=5.4cm,align=left] (l) at (0,0) {\textbf{例 2.1}（原 max）\\[1pt]
  $y_1+2y_2\ \mathbf{\le}\ 2$\quad $y_1+y_3\ \mathbf{\ge}\ 3$\\
  $\cdots\ \mathbf{\ge}\ -5$\quad $\cdots=1$\\
  $y_1\mathbf{\le}0$，\ $y_2\mathbf{\ge}0$，\ $y_3$ 自由};
  \node[gbox,text width=5.4cm,align=left] (r) at (9.0,0) {\textbf{例 2.2}（原 min）\\[1pt]
  $y_1+2y_2\ \mathbf{\ge}\ 2$\quad $y_1+y_3\ \mathbf{\le}\ 3$\\
  $\cdots\ \mathbf{\le}\ -5$\quad $\cdots=1$\\
  $y_1\mathbf{\ge}0$，\ $y_2\mathbf{\le}0$，\ $y_3$ 自由};
  \draw[{Stealth}-{Stealth},thick,shiftc] (l)--node[above,font=\scriptsize,shiftc!80!black]{只改 max/min}node[below,font=\scriptsize,shiftc!80!black]{不等号全部反向}(r);
\end{tikzpicture}
\fig{相近情形辨析：同一组约束，原问题由 max 改为 min 后，对偶中所有“$\le/\ge$”互换，“$=$/自由”不变}
\end{center}

\begin{warningbox}
\begin{enumerate}[leftmargin=1.5em]
  \item \textbf{把 $b$ 和 $c$ 放反。}对偶目标系数来自原问题的\textbf{右端项} $b$，对偶约束右端来自原问题的\textbf{目标系数} $c$。
  \item \textbf{只转置了一半。}对偶约束的系数是 $A$ 的\textbf{列}。检查方法：对偶第 $j$ 条约束的系数个数 $=$ 原问题约束条数 $m$。
  \item \textbf{查表方向错。}原问题是 min 时，必须从右栏读起（纠错 5 中笔记把对偶变量写成 $x_i$，就是查表时串了栏）。
  \item \textbf{先化标准型再写对偶。}不必！对偶直接由原模型写出；先加松弛变量反而会多出一堆约束。
\end{enumerate}
\end{warningbox}

\begin{lifebox}[闲置食材：卖给邻居，还是自己做点心？]
\begin{minipage}[t]{0.54\linewidth}
\textbf{生活场景：}你周末在家烤饼干出售。邻居想整包买走你的面粉和黄油，问你最低多少钱肯卖。\par
\textbf{对应关系：}见右图——你的“生产问题”是原问题，邻居的“出价问题”是对偶问题；食材数量是 $b$，食材单价是 $y$，每份饼干的收入是 $c$。\par
\textbf{小判断：}一份饼干用面粉 2 杯、黄油 1 块，能卖 10 元。若邻居出价面粉 3 元/杯、黄油 3 元/块，则这一篮子食材值 $2\times3+1\times3=9<10$ 元，你会拒绝——这正是对偶约束 $2y_1+y_2\ge10$ 不满足。\par
\textbf{边界说明：}真实交易还涉及讨价还价、运输与时间成本，线性模型只抓住了“不能低于自己生产的收益”这一条底线。
\end{minipage}\hfill
\begin{minipage}[t]{0.42\linewidth}\vspace{0pt}\centering
\begin{tikzpicture}[x=1cm,y=1cm]
  \node[wbox,minimum width=2.2cm] (f) at (0,1.8) {面粉 2 杯};
  \node[wbox,minimum width=2.2cm] (b) at (0,0.9) {黄油 1 块};
  \node[wbox,minimum width=2.2cm] (k) at (0,-0.2) {饼干卖 10 元};
  \node[box,minimum width=2.0cm] (a) at (3.4,1.8) {$a_{1j}=2$};
  \node[box,minimum width=2.0cm] (bb) at (3.4,0.9) {$a_{2j}=1$};
  \node[box,minimum width=2.0cm] (c) at (3.4,-0.2) {$c_j=10$};
  \draw[{Stealth}-{Stealth},auxblue] (f)--(a); \draw[{Stealth}-{Stealth},auxblue] (b)--(bb); \draw[{Stealth}-{Stealth},auxblue] (k)--(c);
  \node[gbox,minimum width=4.4cm] at (1.7,-1.35) {$2y_1+y_2\ge10$};
  \node[font=\scriptsize] at (0,2.55) {生活场景};
  \node[font=\scriptsize] at (3.4,2.55) {对偶模型};
\end{tikzpicture}
\fig{生活场景 $\leftrightarrow$ 对偶约束}
\end{minipage}
\end{lifebox}

\begin{checkbox}
\begin{enumerate}[leftmargin=1.5em]
  \item 用引例说明对偶变量 $y_i$ 和对偶约束各代表什么。
  \item 对称形式的对偶中，$A$、$b$、$c$ 分别去了哪里？
  \item 原问题（max）的 $\ge$ 约束对应怎样的对偶变量？为什么？
  \item 原问题是 min 时，规则表应怎样读？
  \item 原问题有 5 个变量、3 个约束，对偶问题有几个变量、几个约束？
\end{enumerate}
\end{checkbox}

\paragraph{到这里，我们已经解决了什么？}任给一个线性规划，我们都能机械地写出它的对偶问题，并能说出对偶变量的经济意义（资源单价）。
\paragraph{下一步还需要解决什么？}写出来的对偶问题和原问题到底有什么数量关系？买方最低的总出价 $w^*$ 与卖方最高的总收入 $z^*$ 谁大？下一章将证明：它们恰好相等。

\begin{center}
\begin{tikzpicture}[node distance=0.35cm and 0.35cm,k/.style={box,minimum height=0.8cm,font=\small}]
  \node[k] (a) {生产者模型 $\max c\T x$};
  \node[k,right=of a] (b) {购买者模型 $\min b\T y$};
  \node[k,right=of b] (c) {对称形式：转置};
  \node[k,below=0.55cm of b] (d) {一般形式：正常/反常};
  \node[k,below=0.55cm of c] (e) {四步写对偶};
  \node[gbox,below=0.55cm of a,xshift=-1.2cm] (f) {对偶的对偶 $=$ 原问题};
  \draw[arr] (a)--(b); \draw[arr] (b)--(c); \draw[arr] (c)--(d); \draw[arr] (d)--(e); \draw[arr] (b.south west)--(f.north east);
\end{tikzpicture}
\fig{本章知识地图}
\end{center}

% ============================================================
%  第 3 章  对偶的基本性质与互补松弛（K3、K4）
% ============================================================
\section{对偶的基本性质：两个问题如何相互牵制}

上一章我们会写对偶问题了，但还不知道它和原问题之间有什么\textbf{数量关系}。回到煤电油的例子：卖方自己生产最多能收入 $z^*=428$；买方要让卖方愿意卖，至少要出多少钱？直觉告诉我们，买方的出价不可能低于 428，否则卖方宁可自己生产。本章把这个直觉变成定理，并由此得到一个非常实用的工具——互补松弛条件。

\noindent\textbf{本章回答三个问题：}(1) 原问题与对偶问题的目标值谁大？(2) 什么时候二者相等？(3) 已知一个问题的最优解，能不能不解另一个问题就直接写出它的最优解？\par
\noindent\textbf{本章知识路线：}弱对偶性 $\to$ 最优性判别、无界性 $\to$ 强对偶定理 $\to$ 终表中的对偶解 $\to$ 互补松弛。

\paragraph{约定。}本章以对称形式为准：(P) $\max c\T x$，$Ax\le b$，$x\ge0$；(D) $\min b\T y$，$A\T y\ge c$，$y\ge0$。笔记中把 $y$ 写成行向量（$YA\ge C$），二者只是书写习惯不同。一般形式的结论完全相同。

\subsection{弱对偶性：卖方的收入永远不超过买方的出价}

\begin{theorembox}[定理 3.1　弱对偶性 \srcTag{笔记 P11}]
若 $\bar x$ 是 (P) 的任一可行解，$\bar y$ 是 (D) 的任一可行解，则
\[
\boxed{\;c\T\bar x\ \le\ b\T\bar y\;}
\]
\end{theorembox}

\begin{proofbox}[证明（补全了笔记中省略的条件）]
\textbf{要证明：}任一原可行目标值 $\le$ 任一对偶可行目标值。\quad
\textbf{已知：}① $A\bar x\le b$，② $\bar x\ge0$，③ $A\T\bar y\ge c$，④ $\bar y\ge0$。\par
\textbf{思路：}从两边“夹”出中间量 $\bar y\T A\bar x$。\par
\textbf{详细证明：}
\begin{align*}
\bar y\T A\bar x\ &\le\ \bar y\T b && \text{由①④：}\bar y_i\ge0\text{ 乘 }(A\bar x)_i\le b_i\text{，不变号}\\
\bar y\T A\bar x\ =\ (A\T\bar y)\T\bar x\ &\ge\ c\T\bar x && \text{由③②：}\bar x_j\ge0\text{ 乘 }(A\T\bar y)_j\ge c_j\text{，不变号}
\end{align*}
两式合起来：$c\T\bar x\le\bar y\T A\bar x\le b\T\bar y$。\hfill$\square$\par
\textbf{一句话直观解释：}原问题的收入用“资源的实际消耗”去估，对偶的出价用“资源的总量”去估，而单价非负，所以总量估出来的一定更多。\par
\textbf{条件为什么不能删：}若允许 $\bar y_i<0$，第一步乘负数要变号，结论就不成立。这正是纠错 6 第 2 条强调 $\bar y\ge0$ 的原因。
\end{proofbox}

\begin{figure}[!htbp]\centering
\begin{tikzpicture}[x=0.84cm,y=1cm]
  \draw[ax] (-0.3,0)--(14.5,0) node[right,font=\scriptsize]{目标值};
  \foreach \x/\v in {1.9/190, 4.28/{z^*=w^*=428}, 8.6/860} {\draw (\x,0.1)--(\x,-0.1);}
  \fill[navy!70] (0.6,0) circle (2pt); \fill[navy!70] (1.3,0) circle (2pt); \fill[navy!70] (1.9,0) circle (2pt);
  \fill[greenline] (8.6,0) circle (2pt); \fill[greenline] (10.5,0) circle (2pt); \fill[greenline] (12.6,0) circle (2pt);
  \node[font=\scriptsize,below] at (1.9,-0.15) {$c\T\bar x=190$};
  \node[font=\scriptsize,below] at (8.6,-0.15) {$b\T\bar y=860$};
  \node[optpt] at (4.28,0) {};
  \node[font=\scriptsize,below] at (4.28,-0.15) {$z^*=w^*=428$};
  \draw[decorate,decoration={brace,amplitude=5pt},thick,navy] (0.3,0.25)--(4.2,0.25) node[midway,above=5pt,font=\scriptsize,align=center]{原问题可行解的目标值\\（都在左边）};
  \draw[decorate,decoration={brace,amplitude=5pt},thick,greenline] (4.36,0.25)--(14,0.25) node[midway,above=5pt,font=\scriptsize,align=center]{对偶可行解的目标值\\（都在右边）};
  \fill[shiftc!25] (1.9,-0.85) rectangle (8.6,-0.65);
  \node[font=\scriptsize,shiftc!80!black] at (5.25,-1.1) {仅凭两个可行解就能断定 $190\le z^*\le860$};
\end{tikzpicture}
\fig{弱对偶性的数轴图（煤电油例）：所有原可行值落在左侧，所有对偶可行值落在右侧，二者在最优值处相遇}\label{fig:weak}
\end{figure}

图中蓝点是原问题若干可行解的目标值，绿点是对偶问题若干可行解的目标值。例如 $\bar x=(10,10)$ 原可行，$c\T\bar x=190$；$\bar y=(1,1,1)$ 满足 $9+4+3=16\ge7$、$4+5+10=19\ge12$，对偶可行，$b\T\bar y=860$。橙色区间说明：\textbf{任意一对可行解都给出最优值的一个上下界}，这在大型问题中用来判断“离最优还有多远”。

\subsection{弱对偶性的三个推论}

\begin{theorembox}[推论 3.2　最优性判别 \srcTag{笔记 P11 性质 4}]
若 $\hat x$ 原可行，$\hat y$ 对偶可行，且 $c\T\hat x=b\T\hat y$，则 $\hat x$、$\hat y$ 分别是两个问题的最优解。
\end{theorembox}
\noindent\textbf{证明：}对任意原可行 $x$，由弱对偶 $c\T x\le b\T\hat y=c\T\hat x$，故 $\hat x$ 最优；对偶同理。\hfill$\square$\quad
在图~\ref{fig:weak} 中，这就是“蓝点与绿点重合”的情形。

\begin{theorembox}[推论 3.3　无界性 \srcTag{笔记 P11 性质 3}]
若原问题（对偶问题）有可行解但目标值无界，则其对偶问题（原问题）没有可行解。
\end{theorembox}
\noindent\textbf{证明 \fillTag：}反证。设对偶有可行解 $\bar y$，则由弱对偶，所有原可行解满足 $c\T x\le b\T\bar y$，目标值有上界，与“无界”矛盾。\hfill$\square$

\noindent\textbf{反过来不一定成立}（笔记红字“反过来不一定”）：对偶无可行解时，原问题可能无界，也可能同样无可行解。

\begin{figure}[!htbp]\centering
\begin{tikzpicture}[x=0.62cm,y=0.62cm]
  \begin{scope}
    \fill[feasreg] (0,0)--(1,0)--(5,4)--(5,5)--(0,5)--cycle;
    \axes{5}{5}{x_1}{x_2}
    \draw[con] (1,0)--(5.4,4.4) node[right,font=\scriptsize]{$x_1-x_2=1$};
    \draw[iso] (0,1)--(1,0); \draw[iso] (0,3)--(3,0); \draw[iso] (0,5)--(5,0);
    \draw[arr,shiftc] (2.6,2.6)--(4,4) node[above left,font=\scriptsize]{$z\to+\infty$};
    \node[font=\small,align=center] at (2.6,-1.8) {(a) 原问题无界：\\$\max x_1+x_2,\ x_1-x_2\le1$};
    \node[font=\scriptsize,align=center] at (2.6,-3.3) {其对偶 $\min y$，$y\ge1$，$-y\ge1$\\无可行解};
  \end{scope}
  \begin{scope}[xshift=9cm]
    \fill[feasreg,opacity=0.8] (0,0)--(0,5)--(5,5)--(5,4)--(1,0)--cycle;
    \fill[pattern=north east lines,pattern color=redline!60] (2,0)--(5,3)--(5,0)--cycle;
    \axes{5}{5}{x_1}{x_2}
    \draw[con] (1,0)--(5.4,4.4) node[right,font=\scriptsize]{$x_1-x_2=1$};
    \draw[con,redline] (2,0)--(5.4,3.4) node[right,font=\scriptsize]{$x_1-x_2=2$};
    \node[font=\scriptsize,lab] at (1.4,3.4) {$x_1-x_2\le1$};
    \node[font=\scriptsize,lab,text=redline] at (4.2,0.8) {$x_1-x_2\ge2$};
    \node[font=\small,align=center] at (2.6,-1.8) {(b) 两侧都不可行：\\$\max 2x_1-x_2$};
    \node[font=\scriptsize,align=center] at (2.6,-3.3) {两个区域没有公共部分；\\其对偶 $y_1-y_2\ge2$ 与 $y_1-y_2\le1$ 同样矛盾};
  \end{scope}
\end{tikzpicture}
\fig{无界性推论与它的反例：(a) 原问题无界 $\Rightarrow$ 对偶无可行解；(b) 对偶无可行解时，原问题也可能无可行解，而不是无界}\label{fig:unb}
\end{figure}

图 (a) 中可行域向右上方无限延伸，等值线（虚线）可以一直往外推，目标值无界；它的对偶由两条互相矛盾的约束组成，没有可行解。图 (b) 中原问题的两条约束 $x_1-x_2\le1$（绿色一侧）与 $-x_1+x_2\le-2$ 即 $x_1-x_2\ge2$（红色斜线一侧）没有公共部分，对偶也矛盾——这说明“对偶无可行解”推不出“原问题无界”。

\subsection{强对偶定理：两者的最优值恰好相等}

弱对偶只说“左边 $\le$ 右边”。最重要的结论是：只要有最优解，两边一定能碰上。

\begin{theorembox}[定理 3.4　强对偶定理（对偶定理）\srcTag{笔记 P11 性质 5}]
若原问题有最优解，则对偶问题也有最优解，且两者最优值相等：$z^*=w^*$。
\end{theorembox}

\begin{proofbox}[证明（借助单纯形法的最优表）\fillTag]
\textbf{思路：}直接从原问题的最优表中“读出”一个对偶解，验证它对偶可行、且目标值等于 $z^*$，再用推论 3.2。\par
\textbf{第一步。}把 (P) 加松弛变量 $x_s$ 化为 $Ax+Ix_s=b$。设单纯形法得到最优基 $B$，令 $y\T=c_B\T B^{-1}$。\par
\textbf{第二步（$y$ 满足对偶约束）。}最优表中决策变量的检验数非正：$\sigma_j=c_j-y\T P_j\le0$，即 $A\T y\ge c$。\par
\textbf{第三步（$y\ge0$）。}松弛变量 $s_i$ 的列是单位向量 $e_i$、价值系数为 0，其检验数 $\sigma(s_i)=0-y\T e_i=-y_i\le0$，即 $y_i\ge0$。\par
\textbf{第四步（目标值相等）。}$b\T y=y\T b=c_B\T B^{-1}b=c_B\T x_B=z^*$（核心公式 1.1）。\par
\textbf{结论：}$y$ 对偶可行且 $b\T y=z^*$，由推论 3.2，$y$ 是对偶最优解，$w^*=z^*$。\hfill$\square$\par
\textbf{一句话直观解释：}单纯形法在找原问题最优解的同时，顺手把对偶最优解 $y\T=c_B\T B^{-1}$ 也找出来了。\par
\textbf{严格性说明：}上述证明默认单纯形法能在有限步内终止于最优表（需要防循环规则，前面章节已讲）。不借助单纯形法的证明要用到 Farkas 引理，属于后续“最优化方法”课程的内容。
\end{proofbox}

\begin{figure}[!htbp]\centering
\begin{tikzpicture}[node distance=0.4cm and 0.55cm,s/.style={box,minimum height=0.95cm,text width=2.7cm}]
  \node[s] (a) {原问题最优表\\基 $B$};
  \node[s,right=of a] (b) {令\\$y\T=c_B\T B^{-1}$};
  \node[s,right=of b,yshift=0.9cm] (c) {决策列 $\sigma_j\le0$\\$\Rightarrow A\T y\ge c$};
  \node[s,right=of b,yshift=-0.9cm] (d) {松弛列 $\sigma=-y_i\le0$\\$\Rightarrow y\ge0$};
  \node[s,right=of c,yshift=-0.9cm] (e) {$b\T y=c_B\T B^{-1}b$\\$=z^*$};
  \node[gbox,right=of e,text width=1.9cm] (f) {推论 3.2\\$\Rightarrow w^*=z^*$};
  \draw[arr] (a)--(b); \draw[arr] (b)--(c); \draw[arr] (b)--(d); \draw[arr] (c)--(e); \draw[arr] (d)--(e); \draw[arr] (e)--(f);
\end{tikzpicture}
\fig{强对偶定理的推导路线图：每一步的依据都来自最优表}
\end{figure}

\begin{examplebox}[例 3.1　验证煤电油问题 $z^*=w^*$ \srcTag{笔记 P13}]
\begin{minipage}[t]{0.5\linewidth}
原问题终表（已按纠错 3 修正）：
\begin{center}\small
\begin{tabular}{@{}ccc|ccccc@{}}
\toprule
$c_B$ & $x_B$ & $b$ & $x_1$ & $x_2$ & $x_3$ & $x_4$ & $x_5$\\\midrule
0 & $x_3$ & 84 & 0 & 0 & 1 & $-3.12$ & $1.16$\\
7 & $x_1$ & 20 & 1 & 0 & 0 & $0.4$ & $-0.2$\\
12 & $x_2$ & 24 & 0 & 1 & 0 & $-0.12$ & $0.16$\\\midrule
\multicolumn{3}{c|}{$\sigma_j$} & 0 & 0 & \cellcolor{shiftc!25}0 & \cellcolor{shiftc!25}$-1.36$ & \cellcolor{shiftc!25}$-0.52$\\\bottomrule
\end{tabular}
\end{center}
\end{minipage}\hfill
\begin{minipage}[t]{0.47\linewidth}
松弛变量 $x_3,x_4,x_5$ 的检验数取相反数：
\[y^*=(0,\ 1.36,\ 0.52).\]
\textbf{对偶可行：}$9\cdot0+4\cdot1.36+3\cdot0.52=5.44+1.56=7\ge7$ ✓；$4\cdot0+5\cdot1.36+10\cdot0.52=6.8+5.2=12\ge12$ ✓。\par
\textbf{目标值：}$w=360\cdot0+200\cdot1.36+300\cdot0.52=272+156=428=z^*$ ✓。
\end{minipage}
\end{examplebox}
表中橙色单元格是松弛变量的检验数，它们取相反数恰好就是对偶最优解。这就是笔记 P12 红字：“\textbf{对偶问题的最优解就是原问题松弛变量检验数的负值}。”

\begin{figure}[!htbp]\centering
\begin{tikzpicture}[x=1cm,y=1cm,c/.style={draw=navy,minimum width=3.3cm,minimum height=0.7cm,font=\small}]
  \node[c,fill=lightblue] (a) at (0,0) {$x_B$（基变量）};
  \node[c,fill=lightblue] (b) at (3.4,0) {$x_N$（非基决策变量）};
  \node[c,fill=shiftc!20] (c) at (6.8,0) {$x_s$（松弛变量）};
  \node[c,fill=lightgreen] (d) at (0,-0.75) {$\sigma=0$};
  \node[c,fill=lightgreen] (e) at (3.4,-0.75) {$c_N-c_B\T B^{-1}N$};
  \node[c,fill=lightgreen] (f) at (6.8,-0.75) {$-c_B\T B^{-1}=-y\T$};
  \node[c,fill=warm] (g) at (0,-1.8) {$-y_{s1}$};
  \node[c,fill=warm] (h) at (3.4,-1.8) {$-y_{s2}$};
  \node[c,fill=warm] (i) at (6.8,-1.8) {$-y$（对偶最优解）};
  \draw[arr] (d)--(g); \draw[arr] (e)--(h); \draw[arr] (f)--(i);
  \node[font=\scriptsize,left] at (-1.8,0) {原问题各列};
  \node[font=\scriptsize,left] at (-1.8,-0.75) {检验数行};
  \node[font=\scriptsize,left] at (-1.8,-1.8) {对偶问题};
  \node[font=\scriptsize,auxblue,align=left] at (10.4,-1.3) {$y_s$：对偶约束的\\剩余变量};
\end{tikzpicture}
\fig{原问题最优表的检验数行与对偶最优解的对应（笔记 P12 第 7 条）：检验数取负号就是对偶的变量值}
\end{figure}

图中第三行是对偶问题各变量的取值：松弛变量列对应对偶决策变量 $y$；决策变量列对应对偶约束的剩余变量 $y_s$（如 $x_j$ 的检验数为 $-\frac{434}{33}$ 时，对偶第 $j$ 条约束多出 $\frac{434}{33}$，见第 5 章）。笔记还提醒：这一读法要求原问题是“$Ax\le b$ 加松弛变量”的形式；若约束是 $\ge$ 或 $=$，要按第 2 章的符号规则调整正负号。

\begin{figure}[!htbp]\centering
\small
\begin{tikzpicture}[x=1cm,y=1cm]
  \foreach \i/\t in {0/有最优解,1/无界,2/无可行解} {
    \node[font=\scriptsize\bfseries,navy] at (2.2+\i*2.4,0.55) {对偶：\t};
    \node[font=\scriptsize\bfseries,navy,anchor=east] at (0.8,-\i*0.8) {原问题：\t};
  }
  \foreach \r/\c/\t/\s in {0/0/{可能（$z^*=w^*$）}/lightgreen,0/1/不可能/lightred,0/2/不可能/lightred,
                         1/0/不可能/lightred,1/1/不可能/lightred,1/2/可能/lightgreen,
                         2/0/不可能/lightred,2/1/可能/lightgreen,2/2/{可能（图~\ref{fig:unb}b）}/lightgreen} {
    \node[draw=white,fill=\s,minimum width=2.35cm,minimum height=0.75cm,font=\scriptsize] at (2.2+\c*2.4,-\r*0.8) {\t};
  }
\end{tikzpicture}
\fig{原问题与对偶问题解的状态只有四种组合可能出现（绿色），其余五种由弱对偶性与强对偶定理排除（红色）}
\end{figure}

这张 $3\times3$ 表是本节所有结论的总结：对角线左上角“都有最优解且相等”是强对偶定理；“一个无界、另一个无可行解”是推论 3.3；右下角“都不可行”由图~\ref{fig:unb}(b) 给出例子。

\begin{lifebox}[搬家装书：至少要几个箱子？]
\begin{minipage}[t]{0.54\linewidth}
\textbf{生活场景：}搬家时你想用最少的纸箱装完所有书。\par
\textbf{对应关系：}任何一种实际装法用了 $k$ 个箱子 $\leftrightarrow$ 原问题（min）的一个可行解，给出\textbf{上界}；“书的总体积 $\div$ 一个箱子的容积”$\leftrightarrow$ 一个对偶可行解，给出\textbf{下界}。\par
\textbf{小估算：}书的总体积相当于 2.3 箱，所以至少要 3 个箱子；你刚好找到一种装法只用 3 个箱子——上下界相等，由推论 3.2，这已经是最好的装法，不必再试。\par
\textbf{边界说明：}箱子个数必须是整数、书的形状也会浪费空间，所以下界 2.3 要向上取整，有时上下界碰不上（这是整数规划的难点）。
\end{minipage}\hfill
\begin{minipage}[t]{0.42\linewidth}\vspace{0pt}\centering
\begin{tikzpicture}[x=1cm,y=1cm]
  \draw[ax] (0,0)--(4.6,0) node[right,font=\scriptsize]{箱数};
  \foreach \x in {1,2,3,4} {\draw (\x,0.08)--(\x,-0.08); \node[below,font=\tiny] at (\x,-0.08) {\x};}
  \fill[greenline] (2.3,0) circle (2.5pt) node[above,font=\scriptsize,align=center,text=greenline]{体积下界\\2.3};
  \fill[navy] (3,0) circle (2.5pt) node[below=8pt,font=\scriptsize,align=center,text=navy]{实际装法\\3 箱};
  \fill[navy!40] (4,0) circle (2.5pt) node[above,font=\scriptsize,text=navy!60]{4 箱};
  \draw[decorate,decoration={brace,amplitude=4pt,mirror},thick,shiftc] (2.3,-1.1)--(3,-1.1) node[midway,below=4pt,font=\scriptsize,shiftc!80!black]{最优值在这里};
  \node[wbox,font=\scriptsize] at (1.0,1.6) {总体积};
  \node[box,font=\scriptsize] at (3.4,1.6) {对偶可行解};
  \draw[{Stealth}-{Stealth},auxblue] (1.6,1.6)--(2.55,1.6);
\end{tikzpicture}
\fig{生活场景 $\leftrightarrow$ 上下界}
\end{minipage}
\end{lifebox}

\subsection{互补松弛定理：资源“有剩余”与“有价值”不能同时成立}

强对偶说 $z^*=w^*$。回头看弱对偶证明中的不等式链 $c\T x\le y\T Ax\le y\T b$：既然两端相等，中间两个“$\le$”都必须取等号。这就是互补松弛。

\begin{theorembox}[定理 3.5　互补松弛定理 \srcTag{笔记 P12 性质 6}\ \srcTag{课件 P33--P40}]
设 $\hat x$、$\hat y$ 分别是 (P)、(D) 的可行解，$\hat x_s=b-A\hat x$ 为原问题松弛变量，$\hat y_s=A\T\hat y-c$ 为对偶剩余变量。则 $\hat x$、$\hat y$ 都是最优解的\textbf{充要条件}是
\[
\boxed{\;\hat y\T\hat x_s=0\quad\text{且}\quad\hat y_s\T\hat x=0\;}
\qquad\Longleftrightarrow\qquad
\begin{cases}
\hat y_i\,(b_i-a_i\T\hat x)=0,& i=1,\dots,m\quad(\text{资源侧})\\
\hat x_j\,(P_j\T\hat y-c_j)=0,& j=1,\dots,n\quad(\text{产品侧})
\end{cases}
\]
\end{theorembox}

\begin{proofbox}[证明 \fillTag（课件用最优表给出了一种证明，这里给出不依赖单纯形表的一般证明）]
\textbf{已知：}$\hat x,\hat y$ 可行，所以 $\hat x_s\ge0$、$\hat y_s\ge0$、$\hat x\ge0$、$\hat y\ge0$。\par
\textbf{思路：}把 $b\T\hat y-c\T\hat x$ 拆成两个非负和。
\begin{align*}
b\T\hat y-c\T\hat x
&=\hat y\T(A\hat x+\hat x_s)-(A\T\hat y-\hat y_s)\T\hat x && \text{代入 } b=A\hat x+\hat x_s,\ c=A\T\hat y-\hat y_s\\
&=\hat y\T\hat x_s+\hat y_s\T\hat x && \text{$\hat y\T A\hat x$ 两项抵消}\\
&=\textstyle\sum_i \hat y_i(\hat x_s)_i+\sum_j(\hat y_s)_j\hat x_j . && \text{每一项都是两个非负数之积}
\end{align*}
\textbf{必要性：}若二者最优，由强对偶 $b\T\hat y=c\T\hat x$，左边为 0；右边是一堆非负数之和，和为 0 则每一项都为 0。\par
\textbf{充分性：}若每一项都为 0，则 $b\T\hat y=c\T\hat x$，由推论 3.2，二者最优。\hfill$\square$\par
\textbf{一句话直观解释：}对偶间隙 $b\T y-c\T x$ 恰好等于“资源剩余量 $\times$ 资源单价”与“产品的亏损额 $\times$ 产量”之和；最优时间隙为 0，所以每一对乘积都只能为 0。
\end{proofbox}

\begin{figure}[!htbp]\centering
\begin{tikzpicture}[x=1cm,y=1cm]
  \foreach \i/\n/\s/\y in {0/煤/84/0, 1/电/0/1.36, 2/油/0/0.52} {
    \begin{scope}[xshift=\i*4.4cm]
      \node[font=\small\bfseries] at (1.3,2.75) {\n};
      \draw[ax] (0,0)--(2.7,0);
      \pgfmathsetmacro{\hs}{\s/40}
      \pgfmathsetmacro{\hy}{\y/0.6}
      \fill[auxblue!70] (0.35,0) rectangle (1.05,\hs);
      \fill[shiftc!80] (1.55,0) rectangle (2.25,\hy);
      \node[font=\scriptsize,above] at (0.7,\hs) {$\s$};
      \node[font=\scriptsize,above] at (1.9,\hy) {$\y$};
      \node[font=\scriptsize,below] at (0.7,0) {剩余};
      \node[font=\scriptsize,below] at (1.9,0) {影子价};
      \node[font=\scriptsize,gbox,inner sep=2pt] at (1.3,-0.85) {乘积 $=0$};
    \end{scope}
  }
\end{tikzpicture}
\fig{互补松弛的“成对柱状图”（煤电油例）：每一种资源的两根柱子中至少有一根高度为 0}
\end{figure}

图中蓝柱是资源剩余量 $(x_3,x_4,x_5)=(84,0,0)$，橙柱是对偶最优解 $(y_1,y_2,y_3)=(0,1.36,0.52)$（两组柱子用了不同的纵向比例尺，只用来看“是否为零”）。煤有剩余，影子价格为 0；电和油用完了，才可能有正的价格。产品侧同理：$x_1=20>0$、$x_2=24>0$，所以两条对偶约束都取等号（例 3.1 中两条约束恰好等于 7 和 12）。

\begin{examplebox}[例 3.2　已知原问题最优解，用互补松弛直接求对偶最优解]
\textbf{已知：}煤电油问题 $x^*=(20,24)$。\quad\textbf{求：}对偶最优解 $y^*$。\par
\textbf{思路：}产品侧：$x_j>0\Rightarrow$ 第 $j$ 条对偶约束取等；资源侧：资源有剩余 $\Rightarrow y_i=0$。\par
\textbf{第一步。}$x_1=20>0$，$x_2=24>0$，两条对偶约束取等号：$9y_1+4y_2+3y_3=7$，$4y_1+5y_2+10y_3=12$。\par
\textbf{第二步。}煤的剩余 $360-276=84>0\Rightarrow y_1=0$。代入得 $4y_2+3y_3=7$，$5y_2+10y_3=12$。\par
\textbf{第三步。}由第一式 $y_2=\frac{7-3y_3}{4}$，代入第二式：$\frac{35-15y_3}{4}+10y_3=12\Rightarrow 25y_3=13\Rightarrow y_3=0.52$，$y_2=\frac{7-1.56}{4}=1.36$。\par
\textbf{结果：}$y^*=(0,1.36,0.52)$，与例 3.1 从终表读出的完全一致。\quad\textbf{检查：}$y^*\ge0$，$w=428=z^*$ ✓。\par
\textbf{本题真正想说明什么：}知道原问题最优解后，对偶最优解只需解一个小方程组，而不必再做一次单纯形法。
\end{examplebox}

\begin{warningbox}
\textbf{错误想法 1：}“影子价格为 0，所以这种资源一定有剩余。”\par
\textbf{错在哪里：}互补松弛只保证两个因子\textbf{至少一个}为 0，不保证“恰好一个”。退化时资源恰好用尽，影子价格也可能是 0（见图~\ref{fig:degen}）。\par
\textbf{错误想法 2：}“任意一对可行解都满足互补松弛。”\par
\textbf{错在哪里：}定理说的是“二者\textbf{都最优}当且仅当满足互补松弛”。对于非最优的可行解对，对偶间隙 $>0$，必有某一项乘积为正。
\end{warningbox}

\begin{figure}[!htbp]\centering
\begin{tikzpicture}[x=0.6cm,y=0.6cm]
  \begin{scope}
  \fill[feasreg] (0,0)--(4,0)--(4,2)--(0,2)--cycle;
  \axes{6}{4}{x_1}{x_2}
  \draw[con] (4,-0.2)--(4,4.3) node[above,font=\scriptsize]{材料 $A$};
  \draw[con,redline] (-0.2,2)--(6.3,2) node[right,font=\scriptsize]{材料 $B$：$4x_2\le8$};
  \draw[con] (-0.4,4.2)--(6.3,0.85) node[right,font=\scriptsize]{台时};
  \node[optpt] at (4,2) {};
  \node[lab,above right] at (4,2.1) {$(4,2)$：三线共点};
  \end{scope}
  \node[font=\small,align=left,text width=7.2cm] at (14.6,2) {把讲义主例中材料 $B$ 的可用量由 12 改为 8，三条边界线交于同一点 $(4,2)$，出现退化。此时一个最优基为 $\{x_4,x_2,x_1\}$，其中 $x_4=0$ 是取值为零的基变量：材料 $B$ \textbf{恰好用尽}，但它的检验数为 0，影子价格 $y_2=0$。两个因子同时为零，互补松弛照样成立（$0\times0=0$），只是不再“恰有一个为零”。};
\end{tikzpicture}
\fig{退化时的互补松弛（对应课件 P42“退化时的互补松弛”）}\label{fig:degen}
\end{figure}

\begin{lifebox}[冰箱里剩下的鸡蛋与用光的酱油]
\begin{minipage}[t]{0.56\linewidth}
\textbf{生活场景：}按一周菜单做完饭，冰箱里还剩 4 个鸡蛋，而酱油恰好用完。\par
\textbf{对应关系：}剩下的鸡蛋 $\leftrightarrow$ 松弛变量 $s_i>0$；“再多买一盒鸡蛋，这周的饭也不会更好”$\leftrightarrow$ $y_i=0$；酱油用光 $\leftrightarrow$ $s_i=0$；“多一瓶酱油就能多做一道红烧菜”$\leftrightarrow$ $y_i>0$。\par
\textbf{小判断：}鸡蛋：$s=4$，$y=0$，乘积 0 ✓；酱油：$s=0$，$y>0$，乘积 0 ✓。“有剩余又很值钱”的组合不会在最优安排中出现：如果鸡蛋多余还很值钱，说明菜单没排好，还能改进。\par
\textbf{边界说明：}真实生活中剩余的鸡蛋下周还能用，有“跨期价值”；单期线性模型不考虑这一点。
\end{minipage}\hfill
\begin{minipage}[t]{0.40\linewidth}\vspace{0pt}\centering
\begin{tikzpicture}[x=1cm,y=1cm]
  \node[wbox,minimum width=1.9cm] (e) at (0,1.2) {剩 4 个鸡蛋};
  \node[wbox,minimum width=1.9cm] (s) at (0,0) {酱油用光};
  \node[box,minimum width=2.3cm] (e2) at (3.1,1.2) {$s_i>0,\ y_i=0$};
  \node[box,minimum width=2.3cm] (s2) at (3.1,0) {$s_i=0,\ y_i>0$};
  \draw[{Stealth}-{Stealth},auxblue] (e)--(e2); \draw[{Stealth}-{Stealth},auxblue] (s)--(s2);
  \node[rbox,minimum width=4.6cm,font=\scriptsize] at (1.55,-1.2) {$s_i>0$ 且 $y_i>0$：最优时不可能};
\end{tikzpicture}
\fig{生活场景 $\leftrightarrow$ 互补松弛}
\end{minipage}
\end{lifebox}

\begin{table}[!htbp]\centering
\small
\begin{tabularx}{\linewidth}{@{}>{\bfseries}lXX@{}}
\toprule
\rowcolor{lightblue}比较维度 & 弱对偶性 & 强对偶定理\\\midrule
适用对象 & \textbf{任意}一对可行解 & 两个问题的\textbf{最优}解\\
结论 & $c\T\bar x\le b\T\bar y$（不等式） & $z^*=w^*$（等式）\\
证明工具 & 两次“非负数乘不等式” & 单纯形法最优表 $y\T=c_B\T B^{-1}$\\
图像 & 数轴上蓝点在左、绿点在右（图~\ref{fig:weak}） & 蓝点与绿点在最优值处重合\\
典型用途 & 给出最优值的上下界；推出无界性 & 从终表读对偶解；导出互补松弛\\
常见错误 & 忘记 $\bar y\ge0$ 的条件 & 以为“对偶无可行解 $\Rightarrow$ 原问题无界”\\
\bottomrule
\end{tabularx}\par\nopagebreak[4]
\captionof{table}{相近概念辨析：弱对偶性与强对偶定理}
\end{table}

\begin{checkbox}
\begin{enumerate}[leftmargin=1.5em]
  \item 弱对偶性证明中，哪两个地方用到了变量的非负性？
  \item 已知原（max）可行解目标值 30、对偶可行解目标值 34，最优值在什么范围？
  \item 为什么说“原问题无界 $\Rightarrow$ 对偶无可行解”，反过来却不成立？
  \item 怎样从原问题最优表中读出对偶最优解？
  \item 写出互补松弛条件的两组等式，并各用一句话解释。
  \item 影子价格为 0 时，资源一定有剩余吗？
\end{enumerate}
\end{checkbox}

\paragraph{到这里，我们已经解决了什么？}原问题与对偶问题的目标值“左小右大、最优时相等”，对偶最优解可以从终表直接读出，也可以用互补松弛求出。
\paragraph{下一步还需要解决什么？}$y^*=(0,1.36,0.52)$ 这组数在经济上究竟意味着什么？“电”的价格 1.36 能指导什么决策？下一章专门讨论这组数——影子价格。

\begin{figure}[!htbp]\centering
\begin{tikzpicture}[node distance=0.35cm and 0.4cm,k/.style={box,minimum height=0.8cm,font=\small}]
  \node[k] (a) {弱对偶性 $c\T\bar x\le b\T\bar y$};
  \node[k,right=of a] (b) {最优性判别};
  \node[k,right=of b] (c) {无界性推论};
  \node[k,below=0.55cm of a] (d) {强对偶定理 $z^*=w^*$};
  \node[k,right=of d] (e) {终表读 $y=c_B\T B^{-1}$};
  \node[k,right=of e] (f) {互补松弛};
  \draw[arr] (a)--(b); \draw[arr] (b)--(c); \draw[arr] (a)--(d); \draw[arr] (b)--(d); \draw[arr] (d)--(e); \draw[arr] (d)--(f); \draw[arr] (e)--(f);
\end{tikzpicture}
\fig{本章知识地图}
\end{figure}

% ============================================================
%  第 4 章  影子价格（K5）
% ============================================================
\section{影子价格：资源在企业内部值多少钱}

上一章得到：原问题最优表中松弛变量检验数的相反数就是对偶最优解 $y^*$。这一章换一个更贴近决策的问法：\textbf{某种资源多给一个单位，最优利润能增加多少？}我们会看到，答案就是 $y_i^*$，它被称为第 $i$ 种资源的\textbf{影子价格}（shadow price）。

\noindent\textbf{本章回答三个问题：}(1) 影子价格怎样从最优表中读出？(2) 它为什么等于“资源多一单位时最优值的增量”？(3) 它能指导哪些决策，又在什么范围内有效？\par
\noindent\textbf{本章知识路线：}最优表 $\to$ 读出 $B^{-1}$ $\to$ $y\T=c_B\T B^{-1}$ $\to$ 几何验证 $\to$ 经济意义 $\to$ 使用范围。

\subsection{本章主例与它的最优表}

\begin{figure}[!htbp]\centering
\begin{tikzpicture}[x=0.5cm,y=0.6cm]
  \lecRegion
  \axes{9}{5}{x_1}{x_2}
  \lecLines
  \draw[iso] (0,4.667)--(7,0) node[below,font=\scriptsize,black!60]{$z=14$};
  \foreach \p in {(0,0),(4,0),(2,3),(0,3)} \node[vpt] at \p {};
  \node[optpt] at (4,2) {};
  \node[lab,above right] at (4,2) {$(4,2)$};
  \node[font=\small,green!40!black] at (1.5,1.2) {可行域};
\end{tikzpicture}
\fig{主例可行域：最优顶点 $(4,2)$ 由材料 $A$ 与台时两条约束交出}\label{fig:lec}
\end{figure}
\begin{examplebox}[主例　两种原材料与一台设备 \srcTag{课件 P5--P10}]
某厂用材料 $A$、材料 $B$ 和一台设备生产两种产品，单位利润分别为 2 与 3：
\begin{center}\small
\begin{tabular}{@{}lccc@{}}\toprule
资源 & 产品一 & 产品二 & 可用量\\\midrule
材料 $A$ & 4 & 0 & 16 kg\\ 材料 $B$ & 0 & 4 & 12 kg\\ 设备台时 & 1 & 2 & 8 台时\\\midrule
单位利润 & 2 & 3 & ---\\\bottomrule
\end{tabular}
\end{center}
\[
\begin{aligned}\max z&=2x_1+3x_2\\ \st 4x_1+x_3&=16\quad(A)\\ 4x_2+x_4&=12\quad(B)\\ x_1+2x_2+x_5&=8\quad(\text{台时}),\quad x_j\ge0\end{aligned}
\]
松弛变量 $x_3,x_4,x_5$ 分别是三种资源的剩余量。

\end{examplebox}

单纯形法经过 3 次换基得到最优表（课件 P8）：

\begin{table}[!htbp]\centering
\small
\begin{tabular}{@{}ccc|ccccc@{}}
\toprule
& $c_j$ & & 2 & 3 & 0 & 0 & 0\\
$c_B$ & $x_B$ & $b$ & $x_1$ & $x_2$ & $x_3$ & $x_4$ & $x_5$\\\midrule
0 & $x_4$ & 4 & 0 & 0 & \cellcolor{formulablue}$\frac12$ & \cellcolor{formulablue}1 & \cellcolor{formulablue}$-2$\\
3 & $x_2$ & 2 & 0 & 1 & \cellcolor{formulablue}$-\frac18$ & \cellcolor{formulablue}0 & \cellcolor{formulablue}$\frac12$\\
2 & $x_1$ & 4 & 1 & 0 & \cellcolor{formulablue}$\frac14$ & \cellcolor{formulablue}0 & \cellcolor{formulablue}0\\\midrule
\multicolumn{3}{c|}{$\sigma_j$} & 0 & 0 & \cellcolor{shiftc!25}$-\frac18$ & \cellcolor{shiftc!25}0 & \cellcolor{shiftc!25}$-\frac32$\\\bottomrule
\end{tabular}\par\nopagebreak[4]
\captionof{table}{主例最优表。蓝色块即 $B^{-1}$，橙色块是松弛变量的检验数}\label{tab:lec}
\end{table}

读表：$x_1=4$，$x_2=2$，$x_4=4$（材料 $B$ 剩 4 kg），$x_3=x_5=0$（材料 $A$ 与台时用尽），最优值 $z^*=2\times4+3\times2=14$。图~\ref{fig:lec} 中，最优点 $(4,2)$ 恰好落在“材料 $A$”竖线与“台时”斜线的交点上，而“材料 $B$”水平线 $x_2=3$ 在它上方——这与表中“材料 $B$ 有剩余”完全吻合。

\subsection{从最优表读出 $B^{-1}$}

\begin{theorembox}[结论 4.1　$B^{-1}$ 就印在最优表上 \srcTag{课件 P11}]
三个松弛变量在初始方程组中的列并排是单位阵 $I$；最优表 $=B^{-1}\times$ 初始表，所以这三列在最优表中变成 $B^{-1}I=B^{-1}$。行次序是基变量次序 $x_4,x_2,x_1$，列次序是三条约束（材料 $A$、材料 $B$、台时）：
\[
B^{-1}=\begin{pmatrix}\frac12&1&-2\\[2pt]-\frac18&0&\frac12\\[2pt]\frac14&0&0\end{pmatrix}.
\]
\end{theorembox}

\begin{figure}[!htbp]\centering
\begin{tikzpicture}[x=1cm,y=1cm]
  \node (M) at (0,0) {$\begin{pmatrix}\frac12&1&-2\\[2pt]-\frac18&0&\frac12\\[2pt]\frac14&0&0\end{pmatrix}$};
  \node (b) at (2.4,0) {$\begin{pmatrix}16\\12\\8\end{pmatrix}$};
  \node at (3.5,0) {$=$};
  \node (r) at (4.5,0) {$\begin{pmatrix}4\\2\\4\end{pmatrix}$};
  \fill[shiftc!25,opacity=0.6] (-1.45,0.32) rectangle (1.45,0.78);
  \node[font=\scriptsize,align=left] at (8.4,0.55) {第 1 行：$\frac12\cdot16+1\cdot12+(-2)\cdot8=4=x_4$ ✓};
  \node[font=\scriptsize,align=left] at (8.4,0) {第 2 行：$-\frac18\cdot16+0+\frac12\cdot8=2=x_2$ ✓};
  \node[font=\scriptsize,align=left] at (8.4,-0.55) {第 3 行：$\frac14\cdot16+0+0=4=x_1$ ✓};
  \node[font=\scriptsize,auxblue] at (2.4,-1.15) {原模型右端 $b$};
  \node[font=\scriptsize,auxblue] at (4.5,-1.15) {最优表 $b$ 列};
  \node[font=\scriptsize,auxblue] at (0,-1.15) {最优表中 $x_3,x_4,x_5$ 三列};
\end{tikzpicture}
\fig{验算 $x_B=B^{-1}b$：用原模型右端项 $(16,12,8)\T$ 去乘，应还原出最优表的 $b$ 列（课件 P13）}
\end{figure}

图中高亮的是第 1 行的计算过程。三行结果 $(4,2,4)$ 与最优表 $b$ 列逐一相同，说明 $B^{-1}$ 读对了。注意：被乘的是\textbf{原模型}的 $b=(16,12,8)\T$，不是最优表的 $b$ 列——这就是课件红字想提醒的事（纠错 6 第 7 条）。

\subsection{影子价格的定义与计算}

\begin{definitionbox}[定义 4.2　影子价格 \srcTag{课件 P15--P17}\ \srcTag{笔记 P13}]
设 $B$ 为最优基，称
\[
\boxed{\;y\T=c_B\T B^{-1}\;}
\]
的第 $i$ 个分量 $y_i$ 为第 $i$ 种资源的\textbf{影子价格}。同一组数在不同语境下有不同名称：矩阵表达中称\textbf{单纯形乘子}，对偶理论中称\textbf{对偶变量}（对偶最优解），经济分析中称\textbf{影子价格}。\par
\textbf{前提}（笔记 P13）：原问题为 $\max z,\ Ax\le b,\ x\ge0$ 型，对偶为 min 型。
\end{definitionbox}

\begin{examplebox}[例 4.1　计算主例的影子价格]
\textbf{已知：}基为 $\{x_4,x_2,x_1\}$，价值系数 $c_B=(0,3,2)\T$，$B^{-1}$ 如上。\quad\textbf{求：}$y$。\par
\textbf{思路：}$y_i=c_B\T\times(B^{-1}\text{ 的第 }i\text{ 列})$。
\[
y_1=0\cdot\tfrac12+3\cdot(-\tfrac18)+2\cdot\tfrac14=\tfrac18,\qquad
y_2=0\cdot1+3\cdot0+2\cdot0=0,\qquad
y_3=0\cdot(-2)+3\cdot\tfrac12+2\cdot0=\tfrac32 .
\]
\textbf{结果：}$y=(\frac18,\,0,\,\frac32)$。\quad
\textbf{检查：}与表~\ref{tab:lec} 橙色块 $(-\frac18,0,-\frac32)$ 逐项互为相反数 ✓；$y\T b=\frac18\cdot16+0+\frac32\cdot8=2+12=14=z^*$ ✓。
\end{examplebox}

\begin{proofbox}[为什么“松弛变量检验数的相反数”就是影子价格（课件 P19）]
第 $i$ 条约束的松弛变量记作 $s_i$，它在原方程组中的列是单位向量 $e_i$，价值系数为 0。代入检验数公式 $\sigma_j=c_j-y\T P_j$：
\[
\sigma(s_i)=0-y\T e_i=-y_i .
\]
所以影子价格不必另算，\textbf{读松弛变量所在列的检验数、取相反数即可}。又因最优表满足 $\sigma\le0$，得 $y_i=-\sigma(s_i)\ge0$：在 $Ax\le b$ 型问题中，影子价格非负。
\end{proofbox}

\subsection{影子价格的含义：资源的边际贡献}

\begin{theorembox}[定理 4.3　影子价格 $=$ 最优值对资源量的变化率]
若右端项 $b_i$ 改变后最优基 $B$ 不变，则
\[
z^*=c_B\T B^{-1}b=y\T b=\sum_i y_ib_i,\qquad\text{从而}\qquad
\boxed{\;y_i=\frac{\partial z^*}{\partial b_i}\;}
\]
即第 $i$ 种资源每增加 1 单位，最优值增加 $y_i$。
\end{theorembox}

\begin{proofbox}[推导 \fillTag]
\textbf{第一步。}由核心公式 1.1，$z^*=c_B\T x_B=c_B\T B^{-1}b$。\par
\textbf{第二步。}$b$ 变为 $b+\Delta b$ 时，只要最优基 $B$ 不变（第 6 章会求出这个范围），$c_B$ 和 $B^{-1}$ 都不变，于是新的最优值为 $c_B\T B^{-1}(b+\Delta b)=z^*+y\T\Delta b$。\par
\textbf{第三步。}取 $\Delta b=e_i$（只把第 $i$ 种资源加 1），增量恰为 $y_i$。在最优基不变的范围内 $z^*$ 是 $b$ 的线性函数，所以 $y_i$ 就是偏导数 $\partial z^*/\partial b_i$。\par
\textbf{一句话直观解释：}在最优基不变的范围内，最优值随资源线性变化，斜率就是影子价格。
\end{proofbox}

课件用三幅图逐一验证了这个结论。我们把它们并排画在一起，使用相同的坐标范围，方便比较。

\begin{figure}[!htbp]\centering
\begin{tikzpicture}[x=0.53cm,y=0.6cm]
  % (a) 材料 A +1
  \begin{scope}
    \fill[feasreg] (0,0)--(4.25,0)--(4.25,1.875)--(2,3)--(0,3)--cycle;
    \axes{8}{4}{x_1}{x_2}
    \draw[con,black!30] (4,-0.2)--(4,4.3);
    \draw[shiftl] (4.25,-0.2)--(4.25,4.3) node[above,font=\scriptsize]{$4x_1=17$};
    \draw[con] (-0.2,3)--(8.3,3);
    \draw[con] (-0.4,4.2)--(8.3,-0.15);
    \node[badpt] at (4,2) {}; \node[optpt,fill=shiftc] at (4.25,1.875) {};
    \node[font=\scriptsize,align=center] at (4,-1.6) {(a) 材料 $A$：$16\to17$\\新顶点 $(\frac{17}{4},\frac{15}{8})$\\$z=\frac{113}{8}$，增 $\frac18=y_1$};
  \end{scope}
  % (b) 材料 B +1
  \begin{scope}[xshift=5.35cm]
    \fill[feasreg] (0,0)--(4,0)--(4,2)--(1.5,3.25)--(0,3.25)--cycle;
    \axes{8}{4}{x_1}{x_2}
    \draw[con] (4,-0.2)--(4,4.3);
    \draw[con,black!30] (-0.2,3)--(8.3,3);
    \draw[shiftl] (-0.2,3.25)--(8.3,3.25) node[above left,font=\scriptsize]{$4x_2=13$};
    \draw[con] (-0.4,4.2)--(8.3,-0.15);
    \node[optpt] at (4,2) {};
    \node[font=\scriptsize,align=center] at (4,-1.6) {(b) 材料 $B$：$12\to13$\\最优点仍是 $(4,2)$\\$z=14$，增 $0=y_2$};
  \end{scope}
  % (c) 台时 +1
  \begin{scope}[xshift=10.7cm]
    \fill[feasreg] (0,0)--(4,0)--(4,2.5)--(3,3)--(0,3)--cycle;
    \axes{8}{4}{x_1}{x_2}
    \draw[con] (4,-0.2)--(4,4.3);
    \draw[con] (-0.2,3)--(8.3,3);
    \draw[con,black!30] (-0.4,4.2)--(8.3,-0.15);
    \draw[shiftl] (-0.4,4.7)--(8.3,0.35) node[above left,font=\scriptsize]{$x_1+2x_2=9$};
    \node[badpt] at (4,2) {}; \node[optpt,fill=shiftc] at (4,2.5) {};
    \node[font=\scriptsize,align=center] at (4,-1.6) {(c) 台时：$8\to9$\\新顶点 $(4,\frac52)$\\$z=\frac{31}{2}$，增 $\frac32=y_3$};
  \end{scope}
\end{tikzpicture}
\fig{影子价格的几何验证（课件 P21--P26）：灰线为原边界，橙色虚线为资源加 1 后的新边界，红点为原最优点，橙点为新最优点}\label{fig:sp3}
\end{figure}

三幅图读法相同：灰色实线是原来的边界，橙色虚线是该资源增加 1 单位后的新边界，可行域随之扩大一小条。
\begin{itemize}
  \item 图 (a)：材料 $A$ 的边界右移，最优顶点沿台时斜线移到 $(\frac{17}{4},\frac{15}{8})$，$z=2\cdot\frac{17}{4}+3\cdot\frac{15}{8}=\frac{68+45}{8}=\frac{113}{8}$，比 14 多 $\frac18$；
  \item 图 (b)：材料 $B$ 的边界上移，但最优顶点 $(4,2)$ 根本不在这条线上，可行域虽然扩大，最优点不动，增量为 0；
  \item 图 (c)：台时边界整体外推，最优顶点沿材料 $A$ 竖线上移到 $(4,\frac52)$，$z=8+\frac{15}{2}=\frac{31}{2}$，多 $\frac32$。
\end{itemize}
三个增量恰好是 $(\frac18,0,\frac32)$，与例 4.1 算出的影子价格完全一致。

\begin{supplementbox}[影子价格只在一段范围内是常数 \fillTag]
\begin{minipage}{0.47\linewidth}
把材料 $A$ 的可用量 $b_1$ 从 0 一直增加，最优值 $z^*(b_1)$ 的变化如右图（第 6 章将求出中间一段的端点）：
\begin{itemize}[leftmargin=1.2em]
  \item $0\le b_1\le8$：最优点在 $(\frac{b_1}{4},3)$，$z^*=\frac{b_1}{2}+9$，斜率 $\frac12$；
  \item $8\le b_1\le32$：最优基即本章的 $\{x_4,x_2,x_1\}$，斜率 $\frac18$；
  \item $b_1\ge32$：材料 $A$ 不再紧缺，$z^*=16$，斜率 0。
\end{itemize}
斜率就是该段的影子价格。它随资源增多而\textbf{逐段变小}——资源越充裕，边际价值越低。
\end{minipage}\hfill
\begin{minipage}{0.5\linewidth}\centering
\begin{tikzpicture}
\begin{axis}[width=7.6cm,height=5.2cm,xmin=0,xmax=40,ymin=8,ymax=17.5,axis lines=left,
  xlabel={材料 $A$ 可用量 $b_1$},ylabel={$z^*$},xtick={0,8,16,32,40},ytick={9,13,14,16},
  grid=major,grid style={black!8},label style={font=\scriptsize},tick label style={font=\scriptsize}]
  \addplot[navy,very thick] coordinates {(0,9)(8,13)(32,16)(40,16)};
  \addplot[only marks,mark=*,mark options={fill=optc,draw=white},mark size=2.5pt] coordinates {(16,14)};
  \node[font=\scriptsize,anchor=south east] at (axis cs:7,11.6) {斜率 $\frac12$};
  \node[font=\scriptsize,anchor=north west] at (axis cs:18,14.1) {斜率 $\frac18=y_1$};
  \node[font=\scriptsize,anchor=south] at (axis cs:36,16) {斜率 0};
  \draw[shiftc,dashed] (axis cs:8,8)--(axis cs:8,13);
  \draw[shiftc,dashed] (axis cs:32,8)--(axis cs:32,16);
  \fill[shiftc!12] (axis cs:8,8) rectangle (axis cs:32,8.6);
  \node[font=\scriptsize,shiftc!80!black] at (axis cs:20,8.9) {最优基不变的范围};
\end{axis}
\end{tikzpicture}
\fig{$z^*$ 随 $b_1$ 变化的折线：每一段的斜率是该段的影子价格}\label{fig:zb1}
\end{minipage}
\end{supplementbox}

图~\ref{fig:zb1} 的横轴是材料 $A$ 的可用量，纵轴是最优值；紫点是当前状态 $b_1=16$、$z^*=14$。橙色区间 $[8,32]$ 内折线是一条直线，斜率 $\frac18$ 就是影子价格；越过端点后斜率改变，“每多 1 kg 增利 $\frac18$”就不再成立。这是使用影子价格时最容易忽视的一点（课件 P61“影子价格为常数的边界”）。

\subsection{影子价格的经济意义与用法}

\srcTag{笔记 P13}\ 笔记把影子价格的意义总结为三条：
\begin{enumerate}
  \item \textbf{资源的边际贡献：}资源转换成经济效益的效率；
  \item \textbf{资源的稀缺程度：}影子价格 $=0$ 多半意味着资源有剩余（严格地说是“有剩余 $\Rightarrow$ 影子价格为 0”，见纠错 6）；影子价格 $>0$ 说明资源在最优方案中已用尽，值越大越稀缺；
  \item \textbf{资源的边际使用价值：}买方愿意为这种资源支付的最低合理出价（笔记 P13：“$Y^*$：买主的最低出价”）。
\end{enumerate}

\begin{table}[!htbp]\centering
\small
\begin{tabularx}{\linewidth}{@{}>{\bfseries}lXX@{}}
\toprule
\rowcolor{lightblue}比较维度 & 市场价格 & 影子价格\\\midrule
由什么决定 & 市场供求关系 & 企业自身的技术条件、资源结构与产品价格\\
对谁相同 & 所有企业面对同一价格 & 不同企业不同；同一企业在不同时期也不同\\
怎样得到 & 市场报价 & 求解线性规划，$y\T=c_B\T B^{-1}$\\
会不会变 & 随行情波动 & 最优基改变时随之改变（图~\ref{fig:zb1} 的折点）\\
经济含义 & 交易价格 & 企业内部的机会成本\\\bottomrule
\end{tabularx}\par\nopagebreak[4]
\captionof{table}{相近概念辨析：影子价格不是市场价格（课件 P28）}
\end{table}

\begin{figure}[!htbp]\centering
\begin{tikzpicture}[x=1cm,y=1cm]
  \foreach \i/\t/\d in {0/{决定是否购入}/{市场价 $<$ 影子价 $\Rightarrow$ 买入有利},1/{对外报价下限}/{卖出报价不应低于影子价},2/{识别生产瓶颈}/{影子价最高者最稀缺，优先扩充}} {
    \node[gbox,text width=4.1cm,minimum height=1.3cm] at (\i*4.8,0) {\textbf{\t}\\[2pt]{\scriptsize \d}};
  }
  \begin{scope}[yshift=-2.3cm]
    \draw[ax] (-1.8,0)--(12.2,0) node[right,font=\scriptsize]{元/单位};
    \foreach \x/\l in {0/0,4/1,8/{1.5}} {\draw (\x,0.08)--(\x,-0.08); \node[below,font=\tiny] at (\x,-0.1) {\l};}
    \fill[shiftc] (8,0) circle (2.5pt) node[above,font=\scriptsize]{台时影子价 $\frac32$};
    \fill[greenline] (4,0) circle (2.5pt) node[above,font=\scriptsize,align=center]{市场租金 1：\\租入有利};
    \fill[redline] (10.4,0) circle (2.5pt) node[above,font=\scriptsize,align=center]{市场租金 1.8：\\不租};
    \fill[auxblue] (0.5,0) circle (2.5pt) node[above,font=\scriptsize]{材料 $A$ $\frac18$};
  \end{scope}
\end{tikzpicture}
\fig{影子价格的三种用法（课件 P29），下方数轴以主例的设备台时为例：市场租金低于 $\frac32$ 才值得租入}
\end{figure}

图下半部分是一条价格数轴：橙点是设备台时的影子价格 $\frac32$。若市场上租 1 台时只要 1 元（绿点，在左侧），租入后每台时净赚 $\frac32-1=\frac12$；若要 1.8 元（红点，在右侧），就不划算。材料 $A$ 的影子价格只有 $\frac18$，说明三种资源中台时最紧缺，是生产的瓶颈。

\begin{warningbox}
\begin{enumerate}[leftmargin=1.5em]
  \item \textbf{超出范围仍按影子价格外推。}“多买 100 kg 材料 $A$ 能多赚 $100\times\frac18=12.5$ 元”是错的：$b_1$ 超过 32 后影子价格变为 0，实际最多多赚 $16\times\frac18=2$ 元（课件 P61）。
  \item \textbf{把影子价格当成市场价格。}影子价格是“在本企业内部值多少”，不是“在市场上卖多少钱”。
  \item \textbf{读表时取错符号或列。}影子价格是\textbf{松弛变量}列的检验数的\textbf{相反数}，不是决策变量列，也不是检验数本身。
\end{enumerate}
\end{warningbox}

\begin{lifebox}[柠檬水摊：多一份冰块值多少钱？]
\begin{minipage}[t]{0.55\linewidth}
\textbf{生活场景：}周末摆摊卖柠檬水（每杯赚 2 元，用 1 份冰、1 个柠檬）和冰茶（每杯赚 3 元，用 2 份冰）。今天只有 10 份冰、4 个柠檬。\par
\textbf{对应关系：}冰、柠檬的数量 $\leftrightarrow b$；每杯利润 $\leftrightarrow c$；最优安排是柠檬水 4 杯、冰茶 3 杯，赚 17 元；冰和柠檬都用完了。由互补松弛，$y_{\text{冰}}+y_{\text{柠檬}}=2$、$2y_{\text{冰}}=3$，得冰的影子价格 1.5 元/份，柠檬 0.5 元/个。\par
\textbf{小估算：}邻摊愿以 1 元/份卖给你冰块。多 1 份冰可多做半杯冰茶，增利 1.5 元，扣掉成本净赚 0.5 元——值得买。\par
\textbf{边界说明：}真实摆摊还受顾客需求限制，冰多到一定程度卖不完，影子价格就会降为 0，这正是“只在范围内有效”。
\end{minipage}\hfill
\begin{minipage}[t]{0.41\linewidth}\vspace{0pt}\centering
\begin{tikzpicture}[x=0.5cm,y=0.5cm]
  \fill[feasreg] (0,0)--(4,0)--(4,3)--(0,5)--cycle;
  \axes{6}{5}{L}{T}
  \draw[con] (-0.3,5.15)--(6.3,1.85) node[right,font=\scriptsize]{冰};
  \draw[con] (4,-0.2)--(4,5.3) node[above,font=\scriptsize]{柠檬};
  \draw[shiftl] (-0.3,5.65)--(6.3,2.35);
  \node[optpt] at (4,3) {}; \node[optpt,fill=shiftc] at (4,3.5) {};
  \node[lab,right] at (4.2,2.8) {$(4,3)$};
  \node[font=\scriptsize,align=center] at (3,-1.6) {冰 $+1$：最优点上移半杯\\利润 $+1.5$};
\end{tikzpicture}
\fig{生活场景 $\leftrightarrow$ 影子价格}
\end{minipage}
\end{lifebox}

\begin{checkbox}
\begin{enumerate}[leftmargin=1.5em]
  \item 影子价格的定义式是什么？在最优表上怎样直接读出？
  \item 为什么 $Ax\le b$ 型问题的影子价格一定非负？
  \item 主例中材料 $B$ 的影子价格为什么是 0？用图和用表各解释一次。
  \item 用 $y\T b$ 验证主例的最优值。
  \item 影子价格与市场价格有哪些不同？
  \item 为什么不能说“多买 100 kg 材料 $A$ 可多赚 12.5 元”？
\end{enumerate}
\end{checkbox}

\paragraph{到这里，我们已经解决了什么？}我们会读、会算、会解释影子价格，并知道它只在“最优基不变”的范围内是常数。
\paragraph{下一步还需要解决什么？}一旦资源变化超出这个范围，最优表中的 $b$ 列会出现负数——此时检验数仍然合格，但解不可行。能不能从这张“半坏”的表接着算，而不从头再来？这就需要下一章的对偶单纯形法。

\begin{figure}[!htbp]\centering
\begin{tikzpicture}[node distance=0.35cm and 0.4cm,k/.style={box,minimum height=0.8cm,font=\small}]
  \node[k] (a) {最优表};
  \node[k,right=of a] (b) {读 $B^{-1}$};
  \node[k,right=of b] (c) {$y\T=c_B\T B^{-1}$};
  \node[k,right=of c] (d) {$y_i=-\sigma(s_i)\ge0$};
  \node[k,below=0.55cm of b] (e) {$y_i=\partial z^*/\partial b_i$};
  \node[k,right=of e] (f) {三种用法};
  \node[rbox,right=of f] (g) {只在范围内有效\\$\to$ 第 6 章};
  \draw[arr] (a)--(b); \draw[arr] (b)--(c); \draw[arr] (c)--(d); \draw[arr] (c)--(e); \draw[arr] (e)--(f); \draw[arr] (f)--(g);
\end{tikzpicture}
\fig{本章知识地图}
\end{figure}

% ============================================================
%  第 5 章  对偶单纯形法（K6）
% ============================================================
\section{对偶单纯形法：从“最优但不可行”走向可行}

前面几章里，单纯形法的每一步都站在一个\textbf{可行}的顶点上，一路改进目标值，直到检验数全部合格。但上一章末尾留下了一种新局面：资源变化后，最优表中的 $b$ 列出现负数，而检验数仍然全部合格。另一种常见局面是：$\min$ 问题带一组 $\ge$ 约束，化成标准型后右端全为负。这两种表的共同特点是\textbf{“检验数合格、解不可行”}。本章介绍专门处理这种表的方法。

\noindent\textbf{本章回答三个问题：}(1) 起始基解不可行时，单纯形法为什么走不下去？(2) 检验数已经合格而基解不可行的表能不能用？(3) 这样的表按什么规则迭代？\par
\noindent\textbf{本章知识路线：}问题引入 $\to$ 原始可行与对偶可行 $\to$ 出基判据 $\to$ 入基判据 $\to$ 算法 $\to$ 例题与几何路径 $\to$ 与原问题单纯形表的对照。


\subsection{问题引入：起始表不可行，单纯形法无从起步}

\begin{figure}[!htbp]\centering
\begin{tikzpicture}[x=0.46cm,y=0.4cm]
  \fill[feasreg] (0,14.5)--(0,12)--(0.121,11.515)--(8.889,0.556)--(10,0)--(12.5,0)--(12.5,14.5)--cycle;
  \draw[gridl] (0,0) grid[step=2] (12,14);
  \draw[ax] (-0.3,0)--(13,0) node[right,font=\scriptsize]{$y_1$};
  \draw[ax] (0,-0.3)--(0,15) node[above,font=\scriptsize]{$y_2$};
  \foreach \i in {2,4,...,12} \node[below,font=\tiny,black!60] at (\i,0) {\i};
  \foreach \i in {2,4,...,14} \node[left,font=\tiny,black!60] at (0,\i) {\i};
  \draw[con] (-0.6,5.3)--(12.5,-1.25);
  \draw[con] (0,11.667)--(10.5,-1.45);
  \draw[con] (-0.3,13.2)--(3.3,-1.2);
  \node[badpt] at (0,0) {};
  \node[font=\scriptsize,lab,text=badc] at (2.6,1.2) {$y=(0,0)$};
  \node[font=\scriptsize,green!35!black] at (8,10) {可行域};
\end{tikzpicture}
\fig{起始基解 $y=(0,0)$ 落在可行域之外}\label{fig:ds0}
\end{figure}
\begin{examplebox}[主例　\srcTag{课件《对偶单纯形法》P2--P20}]
\[
\begin{aligned}
\min w&=2000y_1+1000y_2\\
\st &y_1+2y_2\ge10\\
&\tfrac32y_1+\tfrac65y_2\ge14\\
&4y_1+y_2\ge12\\
&y_1,y_2\ge0
\end{aligned}
\]
$y_1$ 是每单位原材料的价格，$y_2$ 是每台时的价格；三条约束来自第二章的产品搭配题（三种产品单位利润 10、14、12），目标是资源总价最低。\par
\textbf{化标准型：}目标乘 $-1$ 得 $\max w'=-2000y_1-1000y_2$；三条 $\ge$ 约束各减一个剩余变量 $y_3,y_4,y_5$。剩余变量的列是 $-I$，再把三条等式两边同乘 $-1$，就得到单位阵，代价是右端项全部变负：
\[
\begin{aligned}
-y_1-2y_2+y_3&=-10\\
-\tfrac32y_1-\tfrac65y_2+y_4&=-14\\
-4y_1-y_2+y_5&=-12
\end{aligned}
\]

\begin{center}\small
\begin{tabular}{@{}ccc|ccccc@{}}
\toprule
& $c_j$ & & $-2000$ & $-1000$ & 0 & 0 & 0\\
$c_B$ & $x_B$ & $b$ & $y_1$ & $y_2$ & $y_3$ & $y_4$ & $y_5$\\\midrule
0 & $y_3$ & \cellcolor{lightred}$-10$ & $-1$ & $-2$ & 1 & 0 & 0\\
0 & $y_4$ & \cellcolor{lightred}$-14$ & $-\frac32$ & $-\frac65$ & 0 & 1 & 0\\
0 & $y_5$ & \cellcolor{lightred}$-12$ & $-4$ & $-1$ & 0 & 0 & 1\\\midrule
\multicolumn{3}{c|}{$\sigma_j$} & \cellcolor{lightgreen}$-2000$ & \cellcolor{lightgreen}$-1000$ & \cellcolor{lightgreen}0 & \cellcolor{lightgreen}0 & \cellcolor{lightgreen}0\\\bottomrule
\end{tabular}\par\nopagebreak[4]
\captionof{table}{起始表：$b$ 列全负（红），检验数全部 $\le0$（绿）}\label{tab:ds0}
\end{center}
\end{examplebox}

表~\ref{tab:ds0} 的基是 $\{y_3,y_4,y_5\}$，令非基变量 $y_1=y_2=0$，得 $y_3=-10$、$y_4=-14$、$y_5=-12$——三个基变量都是负数，基解不可行。图~\ref{fig:ds0} 中这个点就是原点：价格全取 0 时，三种产品的资源成本都不足以抵偿利润，三条约束一条也不满足。

\paragraph{单纯形法为什么走不下去？}\srcTag{课件 P17}\ 单纯形法每一步都建立在“当前解可行”之上：$\theta$ 比值用 $b$ 列除以入基列，$b$ 为负时比值失去“最多能增加多少”的含义；换基后要保证新解仍可行，而当前解本来就不可行；目标值单调改进的性质也以可行为前提。已有的两条出路是大 $M$ 法和两阶段法，但都要再添 3 个人工变量，表从 5 列变成 8 列。

\paragraph{但这张表有一个好性质。}看检验数行：$c_B$ 全为 0，$\sigma\T=c\T-c_B\T B^{-1}A$ 的被减项为 0，检验数就等于 $c_j=(-2000,-1000,0,0,0)$，全部 $\le0$。也就是说，\textbf{最优性判据已经满足}，只是解还不可行。这是因为目标由 min 改写成 max 时，两个正的价格系数都变了号。

\subsection{两种“可行”：同一张表的两条判据}

\begin{definitionbox}[定义 5.1　原始可行与对偶可行 \srcTag{课件 P20--P24}]
对 max 型问题的一张单纯形表：
\begin{itemize}
  \item 若右端项 $B^{-1}b\ge0$，称该表\textbf{原始可行}（primal feasible）：当前基解满足原问题的非负条件；
  \item 若检验数 $\sigma\le0$，称该表\textbf{对偶可行}（dual feasible）。
\end{itemize}
一张表两条都满足时，当前基可行解即最优解。
\end{definitionbox}

“对偶可行”这个名字来自第 3 章：检验数 $\sigma\le0$ 等价于 $y\T=c_B\T B^{-1}$ 满足对偶约束且 $y\ge0$（强对偶定理证明的第二、三步）。所以 $\sigma\le0$ 恰好说明“由这张表读出的对偶解是可行的”。

\begin{figure}[!htbp]\centering
\begin{tikzpicture}[x=1cm,y=1cm]
  \fill[lightgreen] (0,0) rectangle (5,2.2); \fill[lightred] (5,0) rectangle (10,2.2);
  \draw[navy] (0,0) rectangle (10,2.2); \draw[navy] (5,0)--(5,2.2);
  \node[font=\small\bfseries] at (2.5,1.85) {单纯形法};
  \node[font=\small\bfseries] at (7.5,1.85) {对偶单纯形法};
  \node[font=\small,align=center] at (2.5,0.9) {起点：$b\ge0$（原始可行）\\保住 $b\ge0$\\一路寻找 $\sigma\le0$};
  \node[font=\small,align=center] at (7.5,0.9) {起点：$\sigma\le0$（对偶可行）\\保住 $\sigma\le0$\\一路寻找 $b\ge0$};
  \node[gbox] at (5,-0.8) {终点相同：$b\ge0$ 且 $\sigma\le0$};
  \draw[arr] (2.5,0)--(4.2,-0.6); \draw[arr] (7.5,0)--(5.8,-0.6);
\end{tikzpicture}
\fig{两条路的对照（课件 P22--P24）：两者停在同一个条件上，只是到达的次序相反}\label{fig:twopath}
\end{figure}

图~\ref{fig:twopath} 是本章的核心思想：单纯形法“保住原始可行，去找对偶可行”；对偶单纯形法把两项职责对调，“\textbf{保住对偶可行，去找原始可行}”。两者的终点完全相同。笔记 P14 写得更直白：单纯形法“在保持可行性前提下，通过迭代求出最优解”；对偶单纯形法“在保持最优性前提下，通过迭代改善可行性，从而找到最优解”。

\begin{warningbox}[名称的一处提醒（课件 P25、笔记 P14）]
对偶单纯形法解的是\textbf{当前这个线性规划本身}，不是它的对偶问题。“切不可理解为求解对偶问题时使用的单纯形法”。任何一个起始表对偶可行的线性规划都可以用它，无论这个规划是不是谁的对偶。本章主例恰好是第二章那道题的对偶，只是为了后面便于逐格对照。
\end{warningbox}

\subsection{迭代规则：先定出基，再定入基}

每一轮要定三件事：出基变量、入基变量、旋转运算。次序与单纯形法相反：单纯形法先定入基再定出基，本法\textbf{先定出基再定入基}。

\paragraph{一、出基判据。}要消除的是 $b$ 中的负数，因此在取负的基变量中选取：
\[
\boxed{\;r=\arg\min_i\{(B^{-1}b)_i:(B^{-1}b)_i<0\}\;}\qquad\text{取右端项最小（最负）的那一行出基。}
\]
\srcTag{课件 P28}\ 任取一个取负的基变量换出去，非退化时算法都会终止；取最负的一行只是收敛通常更快——负得越多，离可行越远。这与单纯形法“取检验数最大的列入基”是同一性质：不影响正确性，只影响步数。

\paragraph{二、主元必须为负。}旋转之后出基行的右端变为 $b_r/a_{rk}$。$b_r<0$，要使它变为正数，\textbf{主元 $a_{rk}$ 必须为负}。因此入基列只能在出基行中系数为负的各列中选取。

\paragraph{三、入基判据：保证旋转后检验数仍 $\le0$。}

\begin{proofbox}[推导 \srcTag{课件 P29--P33}]
设第 $r$ 行出基、第 $k$ 列入基，出基行记作 $a_r=(a_{r1},\dots,a_{rn})$。旋转时检验数行减去主元行的 $\frac{\sigma_k}{a_{rk}}$ 倍，得
\[
\sigma_j'=\sigma_j-\frac{\sigma_k}{a_{rk}}\,a_{rj}.
\]
要求对每个 $j$ 都有 $\sigma_j'\le0$。记 $t=\frac{\sigma_k}{a_{rk}}$，因 $\sigma_k\le0$、$a_{rk}<0$，有 $t\ge0$。
\begin{itemize}
  \item $a_{rj}\ge0$ 的列：$\sigma_j'=\sigma_j-t\,a_{rj}\le\sigma_j\le0$，自动成立（减去一个不小于 0 的数，只会更小）；
  \item $a_{rj}<0$ 的列：$\sigma_j-t\,a_{rj}\le0\iff t\le\frac{\sigma_j}{a_{rj}}$（两边除以负数 $a_{rj}$ 要变号）。
\end{itemize}
$t$ 必须同时不超过所有这些上界，所以取最小者：
\[
\boxed{\;\theta=\min_j\Bigl\{\frac{\sigma_j}{a_{rj}}:a_{rj}<0\Bigr\}\;}\qquad\text{取到最小值的列入基。}
\]
这样取，$t$ 恰好等于上界，所有 $\sigma_j'\le0$ 得以保住。
\end{proofbox}

\begin{theorembox}[结论 5.2　无可行解判据 \srcTag{笔记 P15}]
若出基行 $r$ 中所有 $a_{rj}\ge0$，则原问题无可行解。
\end{theorembox}
\noindent\textbf{证明 \fillTag：}第 $r$ 行对应的方程是 $x_{B_r}+\sum_{j\in N}a_{rj}x_j=b_r$。若所有 $a_{rj}\ge0$，则对任意 $x\ge0$，左边 $\ge0$，而右边 $b_r<0$，方程不可能成立。\hfill$\square$

\paragraph{四、旋转运算完全相同。}\srcTag{课件 P35}\ 主元定下之后：主元所在行整体除以主元；其余各行减去主元行的适当倍数，使主元所在列的其余各项化为 0。三条判据里只有前两条是新的。

\begin{table}[!htbp]\centering
\small
\begin{tabularx}{\linewidth}{@{}>{\bfseries}lXX@{}}
\toprule
\rowcolor{lightblue} & 单纯形法 & 对偶单纯形法\\\midrule
先定 & 入基列 $k$（$\sigma_k$ 最大的正数） & 出基行 $r$（$b_r$ 最负）\\
比值由谁除以谁 & $b$ 列除以入基列 & $\sigma$ 行除以出基行\\
只对哪些项取比值 & 入基列系数\textbf{为正}的行 & 出基行系数\textbf{为负}的列\\
取最小值者 & 出基 & 入基\\
主元符号 & 正 & 负\\
无解判据 & 入基列全 $\le0$：无界 & 出基行全 $\ge0$：无可行解\\\bottomrule
\end{tabularx}\par\nopagebreak[4]
\captionof{table}{两种算法的 $\theta$ 规则对照（课件 P34）：两张表沿对角线互换——行与列对调，$b$ 列与 $\sigma$ 行对调，正与负对调}\label{tab:theta}
\end{table}

\begin{figure}[!htbp]\centering
\begin{tikzpicture}[node distance=0.4cm,
  s/.style={box,minimum width=5.0cm,minimum height=0.8cm},
  d/.style={diamond,aspect=2.6,draw=navy,fill=white,font=\small,inner sep=1pt,align=center}]
  \node[s] (a) {起始表满足 $\sigma\le0$（对偶可行）};
  \node[d,below=of a] (b) {$b$ 列\\全部 $\ge0$?};
  \node[gbox,right=1.1cm of b] (opt) {最优，停止};
  \node[s,below=of b] (c) {$b_r$ 最负的行出基};
  \node[d,below=of c] (d) {出基行有\\$a_{rj}<0$?};
  \node[rbox,right=1.1cm of d] (inf) {无可行解，停止};
  \node[s,below=of d] (e) {$\theta=\min\{\sigma_j/a_{rj}:a_{rj}<0\}$ 定入基};
  \node[s,below=of e] (f) {以 $a_{rk}$ 为主元旋转};
  \draw[arr] (a)--(b); \draw[arr] (b)--node[above,font=\scriptsize]{是}(opt);
  \draw[arr] (b)--node[right,font=\scriptsize]{否}(c); \draw[arr] (c)--(d);
  \draw[arr] (d)--node[above,font=\scriptsize]{否}(inf);
  \draw[arr] (d)--node[right,font=\scriptsize]{是}(e); \draw[arr] (e)--(f);
  \draw[arr] (f.west)--++(-1.0,0)|-(b.west);
\end{tikzpicture}
\fig{对偶单纯形法流程图（笔记 P15 五步操作的图示）}\label{fig:dsflow}
\end{figure}

\subsection{例题：两轮迭代与几何路径}

\begin{examplebox}[例 5.1　用对偶单纯形法求解主例 \srcTag{课件 P36--P48}]
\textbf{第一轮。}由表~\ref{tab:ds0}，$b=(-10,-14,-12)$，最负的是 $-14$，$y_4$ 出基。出基行 $a_r=(-\frac32,-\frac65,0,1,0)$，负系数在 $y_1,y_2$ 两列：
\[
y_1:\ \frac{-2000}{-3/2}=\frac{4000}{3},\qquad y_2:\ \frac{-1000}{-6/5}=\frac{2500}{3}.
\]
$\frac{2500}{3}$ 较小，$y_2$ 入基，主元 $-\frac65$。旋转：$y_4$ 行除以 $-\frac65$；$y_3$ 行加上新行的 2 倍，$y_5$ 行加上新行的 1 倍；检验数行减去新行的 $-1000$ 倍。

\begin{center}\small
\begin{tabular}{@{}ccc|ccccc@{}}
\toprule
$c_B$ & $x_B$ & $b$ & $y_1$ & $y_2$ & $y_3$ & $y_4$ & $y_5$\\\midrule
0 & $y_3$ & $\frac{40}{3}$ & $\frac32$ & 0 & 1 & $-\frac53$ & 0\\
$-1000$ & $y_2$ & $\frac{35}{3}$ & $\frac54$ & 1 & 0 & $-\frac56$ & 0\\
0 & $y_5$ & \cellcolor{lightred}$-\frac13$ & $-\frac{11}{4}$ & 0 & 0 & $-\frac56$ & 1\\\midrule
\multicolumn{3}{c|}{$\sigma_j$} & $-750$ & 0 & 0 & $-\frac{2500}{3}$ & 0\\\bottomrule
\end{tabular}
\end{center}
检验数仍全部 $\le0$；$b$ 中只剩 $-\frac13$ 一个负数。

\textbf{第二轮。}$y_5$ 出基，出基行 $(-\frac{11}{4},0,0,-\frac56,1)$，负系数在 $y_1,y_4$：
\[
y_1:\ \frac{-750}{-11/4}=\frac{3000}{11}\approx272.7,\qquad y_4:\ \frac{-2500/3}{-5/6}=1000 .
\]
$y_1$ 入基，主元 $-\frac{11}{4}$。旋转后得：

\begin{center}\small
\begin{tabular}{@{}ccc|ccccc@{}}
\toprule
$c_B$ & $x_B$ & $b$ & $y_1$ & $y_2$ & $y_3$ & $y_4$ & $y_5$\\\midrule
0 & $y_3$ & $\frac{434}{33}$ & 0 & 0 & 1 & $-\frac{70}{33}$ & $\frac{6}{11}$\\
$-1000$ & $y_2$ & $\frac{380}{33}$ & 0 & 1 & 0 & $-\frac{40}{33}$ & $\frac{5}{11}$\\
$-2000$ & $y_1$ & $\frac{4}{33}$ & 1 & 0 & 0 & $\frac{10}{33}$ & $-\frac{4}{11}$\\\midrule
\multicolumn{3}{c|}{$\sigma_j$} & 0 & 0 & 0 & $-\frac{20000}{33}$ & $-\frac{3000}{11}$\\\bottomrule
\end{tabular}
\end{center}
$b$ 全部非负、$\sigma$ 全部非正，最优。\par
\textbf{结果：}$y_1=\frac{4}{33}$，$y_2=\frac{380}{33}$，$y_3=\frac{434}{33}$，$y_4=y_5=0$；$w^*=2000\cdot\frac{4}{33}+1000\cdot\frac{380}{33}=\frac{8000+380000}{33}=\frac{388000}{33}\approx11757.58$。\par
\textbf{检查：}$y_1+2y_2=\frac{4+760}{33}=\frac{764}{33}\approx23.2\ge10$ ✓（剩余 $\frac{434}{33}$，即 $y_3$）；$\frac32\cdot\frac{4}{33}+\frac65\cdot\frac{380}{33}=\frac{2+152}{11}=14$ ✓；$4\cdot\frac4{33}+\frac{380}{33}=\frac{396}{33}=12$ ✓。
\end{examplebox}

\begin{figure}[!htbp]\centering
\begin{tikzpicture}[x=0.5cm,y=0.42cm]
  \fill[feasreg] (0,14.5)--(0,12)--(0.121,11.515)--(8.889,0.556)--(10,0)--(12.5,0)--(12.5,14.5)--cycle;
  \draw[gridl] (0,0) grid[step=2] (12,14);
  \draw[ax] (-0.3,0)--(13,0) node[right,font=\scriptsize]{$y_1$};
  \draw[ax] (0,-0.3)--(0,15) node[above,font=\scriptsize]{$y_2$};
  \foreach \i in {2,4,...,12} \node[below,font=\tiny,black!60] at (\i,0) {\i};
  \foreach \i in {2,4,...,14} \node[left,font=\tiny,black!60] at (0,\i) {\i};
  \draw[con] (-0.6,5.3)--(12.5,-1.25) node[pos=0.12,above,sloped,font=\tiny]{$y_1+2y_2=10$};
  \draw[con] (0,11.667)--(10.5,-1.45) node[pos=0.55,above,sloped,font=\tiny]{$\frac32y_1+\frac65y_2=14$};
  \draw[con] (-0.3,13.2)--(3.3,-1.2) node[pos=0.75,below,sloped,font=\tiny]{$4y_1+y_2=12$};
  \draw[iso] (-0.1,11.958)--(5.95,-0.14);
  \draw[arr,badc,very thick,dashed] (0,0)--(0,11.4);
  \node[badpt] at (0,0) {}; \node[badpt] at (0,11.667) {}; \node[optpt] at (0.121,11.515) {};
  \node[font=\scriptsize,lab,text=badc,right] at (0.3,0.6) {第一张表 $(0,0)$：三条都不满足};
  \node[font=\scriptsize,lab,text=badc,right] at (0.6,13.4) {第二张表 $(0,\frac{35}{3})$：第三条差 $\frac13$};
  \node[font=\scriptsize,lab,text=optc,right] at (1.2,10.4) {第三张表 $(\frac4{33},\frac{380}{33})$：最优};
  \foreach \p/\l in {(8.889,0.556)/{(\frac{80}{9},\frac59)},(10,0)/{(10,0)}} \node[vpt] at \p {};
  \node[font=\tiny,above right] at (8.9,0.6) {$(\frac{80}{9},\frac59)$};
  % 放大图
  \begin{scope}[xshift=8.4cm,yshift=-95.05cm,x=22cm,y=8.5cm]
    \fill[feasreg] (0,12.05)--(0,12)--(0.121,11.515)--(0.24,11.367)--(0.24,12.05)--cycle;
    \draw[black!50] (0,11.3) rectangle (0.24,12.05);
    \draw[con] (0,11.667)--(0.24,11.371);
    \draw[con] (0,12)--(0.175,11.3);
    \draw[arr,optc,thick] (0,11.667)--(0.115,11.522);
    \node[badpt] at (0,11.667) {}; \node[optpt] at (0.121,11.515) {};
    \node[font=\tiny,left] at (0,11.667) {$\frac{35}{3}$};
    \node[font=\tiny,left] at (0,12) {12};
    \node[font=\tiny,right] at (0.15,11.42) {$\frac32y_1+\frac65y_2=14$};
    \node[font=\tiny,left] at (0.16,11.33) {$4y_1+y_2=12$};
    \node[font=\tiny,above right] at (0.121,11.52) {$(\frac4{33},\frac{380}{33})$};
    \node[font=\scriptsize,align=center] at (0.12,11.17) {放大左上角：两点仅相距 0.19};
  \end{scope}
\end{tikzpicture}
\fig{对偶单纯形法的几何路径（课件 P50--P58）：前两个基解在可行域之外，第三个落在顶点上；右侧为左上角放大图}\label{fig:dspath}
\end{figure}

图~\ref{fig:dspath} 中三条蓝线是三条约束的边界，绿色是可行域（向右上方无界：价格定得越高，约束越容易满足）。红色虚线箭头是第一轮：$y_2$ 入基后，点沿纵轴从原点上移到 $(0,\frac{35}{3})$，前两条约束已满足，第三条还差 $\frac13$；第二轮的位移只有 0.19，在全景图上看不清，右侧放大图显示它沿第二条边界线滑到顶点 $(\frac{4}{33},\frac{380}{33})$。灰色虚线是最优等值线 $2000y_1+1000y_2=\frac{388000}{33}$，它支撑在这个顶点上。与单纯形法“在可行域内沿顶点走、起点须先由大 $M$ 法求出”不同，对偶单纯形法\textbf{从域外走进来}。

\begin{figure}[!htbp]\centering
\begin{tikzpicture}[x=1cm,y=1cm]
  \foreach \t/\x in {第一张表/0,第二张表/4.4,第三张表/8.8} \node[font=\small\bfseries] at (\x+1.4,2.75) {\t};
  \foreach \x/\a/\b/\c in {0/-10/-14/-12, 4.4/13.333/11.667/-0.333, 8.8/13.152/11.515/0.121} {
    \begin{scope}[xshift=\x cm]
      \draw[ax] (-0.2,0)--(3,0);
      \foreach \v/\i in {\a/0,\b/1,\c/2} {
        \pgfmathsetmacro{\h}{\v/7}
        \ifdim \h pt<0pt \fill[badc!70] (0.3+\i*0.9,0) rectangle (0.9+\i*0.9,\h); \else \fill[greenline!70] (0.3+\i*0.9,0) rectangle (0.9+\i*0.9,\h); \fi
      }
      \foreach \v/\i in {\a/0,\b/1,\c/2} {\pgfmathsetmacro{\h}{\v/7}\pgfmathsetmacro{\lbl}{\v}
        \ifdim \h pt<0pt \node[font=\tiny,below] at (0.6+\i*0.9,\h) {\pgfmathprintnumber[precision=2]{\lbl}};\else \node[font=\tiny,above] at (0.6+\i*0.9,\h) {\pgfmathprintnumber[precision=2]{\lbl}};\fi}
    \end{scope}
  }
  \node[font=\scriptsize,left] at (-0.3,1.2) {$b$ 列};
  \node[font=\scriptsize,align=left,right] at (12.3,0.3) {红：仍不可行\\绿：已可行};
\end{tikzpicture}
\fig{三张表 $b$ 列的变化：取负的项逐轮减少，减到一个也没有时停止（每组柱依次为三个基变量）}\label{fig:dsbars}
\end{figure}

\paragraph{与原问题那三张表逐格对照。}\srcTag{课件 P59--P67}\ 本例恰好是第二章产品搭配题 $\max 10x_1+14x_2+12x_3$，$x_1+\frac32x_2+4x_3\le2000$，$2x_1+\frac65x_2+x_3\le1000$ 的对偶。原问题用单纯形法也是 3 张表、2 次换基，而且\textbf{每一步一一对应}。

\begin{table}[!htbp]\centering
\small
\begin{tabular}{@{}lll|cccl@{}}
\toprule
\rowcolor{lightblue}表 & 原问题的基 & 对偶的基 & 配对 & 原 $\sigma$ & 对偶解 & 说明\\\midrule
一 & $\{x_4,x_5\}$ & $\{y_3,y_4,y_5\}$ & $x_2\leftrightarrow y_4$ & $14$ & $-14$ & $x_2$ 入 $\leftrightarrow$ $y_4$ 出\\
二 & $\{x_4,x_2\}$ & $\{y_3,y_2,y_5\}$ & $x_3\leftrightarrow y_5$ & $\frac13$ & $-\frac13$ & $x_3$ 入 $\leftrightarrow$ $y_5$ 出\\
三 & $\{x_3,x_2\}$ & $\{y_3,y_2,y_1\}$ & $x_1\leftrightarrow y_3$ & $-\frac{434}{33}$ & $\frac{434}{33}$ & 目标值同为 $\frac{388000}{33}$\\\bottomrule
\end{tabular}\par\nopagebreak[4]
\captionof{table}{原问题单纯形表与对偶单纯形表的对照（摘录）。变量配对：$x_1,x_2,x_3,x_4,x_5\leftrightarrow y_3,y_4,y_5,y_1,y_2$}\label{tab:pair}
\end{table}

表~\ref{tab:pair} 说明了两条等价关系：\textbf{原问题的检验数行，就是对偶的基解取相反数；原问题的基解，就是对偶的检验数行取相反数。}于是“原始可行（$b\ge0$）”与“对偶可行（$\sigma\le0$）”在两张表之间互换身份。第一轮两个 $\theta$ 比值也完全相同（$\frac{4000}{3}$ 与 $\frac{2500}{3}$）。这解释了三条判据为什么这样定：出基取 $b$ 最负的一行，对应原问题取 $\sigma$ 最大的一列；入基取 $\sigma$ 行除以出基行的最小比值，对应原问题取 $b$ 列除以入基列的最小比值。\textbf{这一对照是解释，不是步骤}：实际计算只用摆在面前的那一张表。

\begin{examplebox}[例 5.2　\srcTag{笔记 P15--P16}（已按纠错 4 修正最后一步）]
用对偶单纯形法求解 $\min z=2x_1+3x_2+4x_3$，$\st x_1+2x_2+x_3\ge3$，$2x_1-x_2+3x_3\ge4$，$x_j\ge0$。\par
\textbf{化标准型：}$\max z'=-2x_1-3x_2-4x_3$，两条约束乘 $-1$ 后加松弛变量：$-x_1-2x_2-x_3+x_4=-3$，$-2x_1+x_2-3x_3+x_5=-4$。起始表 $\sigma=(-2,-3,-4,0,0)\le0$，满足前提。

\begin{center}\small
\begin{tabular}{@{}ccc|ccccc|l@{}}
\toprule
$c_B$ & $x_B$ & $b$ & $x_1$ & $x_2$ & $x_3$ & $x_4$ & $x_5$ & 本轮判断\\\midrule
0 & $x_4$ & $-3$ & $-1$ & $-2$ & $-1$ & 1 & 0 & \\
0 & $x_5$ & \cellcolor{lightred}$-4$ & \cellcolor{shiftc!30}$-2$ & 1 & $-3$ & 0 & 1 & $x_5$ 出；$\theta$：$x_1$ 为 1，$x_3$ 为 $\frac43$\\\midrule
\multicolumn{3}{c|}{$\sigma_j$} & $-2$ & $-3$ & $-4$ & 0 & 0 & $\Rightarrow x_1$ 入基\\\midrule
0 & $x_4$ & \cellcolor{lightred}$-1$ & 0 & \cellcolor{shiftc!30}$-\frac52$ & $\frac12$ & 1 & $-\frac12$ & $x_4$ 出；$\theta$：$x_2$ 为 $\frac85$，$x_5$ 为 2\\
$-2$ & $x_1$ & 2 & 1 & $-\frac12$ & $\frac32$ & 0 & $-\frac12$ & \\\midrule
\multicolumn{3}{c|}{$\sigma_j$} & 0 & $-4$ & $-1$ & 0 & $-1$ & $\Rightarrow x_2$ 入基\\\midrule
$-3$ & $x_2$ & $\frac25$ & 0 & 1 & $-\frac15$ & $-\frac25$ & $\frac15$ & \\
$-2$ & $x_1$ & $\frac{11}{5}$ & 1 & 0 & $\frac75$ & $-\frac15$ & $-\frac25$ & $b\ge0$，最优\\\midrule
\multicolumn{3}{c|}{$\sigma_j$} & 0 & 0 & $-\frac95$ & $-\frac85$ & $-\frac15$ & \\\bottomrule
\end{tabular}
\end{center}
\textbf{结果：}$x^*=(\frac{11}{5},\frac25,0,0,0)\T$，$\max z'=-2\cdot\frac{11}{5}-3\cdot\frac25=-\frac{28}{5}$，所以 $\min z=\frac{28}{5}$。\par
\textbf{检查：}$\frac{11}{5}+\frac45=3\ge3$ ✓，$\frac{22}{5}-\frac25=4\ge4$ ✓。\par
\textbf{本题真正想说明什么：}凡是“min 问题、价值系数全非负、约束全是 $\ge$”的模型，乘 $-1$ 后起始表天然对偶可行，可以直接用对偶单纯形法，不必引入人工变量。
\end{examplebox}

表中橙色单元格是每一轮的主元，它们都是负数；红色是本轮出基行的右端。

\begin{warningbox}
\begin{enumerate}[leftmargin=1.5em]
  \item \textbf{起始表不满足 $\sigma\le0$ 也硬用。}对偶单纯形法的前提是对偶可行；笔记：“前提不满足，不可使用（对偶）单纯形法。”
  \item \textbf{比值算成 $b/a$。}对偶单纯形法的比值是“$\sigma$ 行 $\div$ 出基行”，不是“$b$ 列 $\div$ 入基列”。
  \item \textbf{在正系数列中取比值。}只能在出基行中 $a_{rj}<0$ 的列中选；选正系数会使出基行右端仍为负，同时破坏 $\sigma\le0$（练习第 28 题有图示）。
  \item \textbf{漏掉某个负系数列的比值。}笔记 P16 第二轮把 $x_5$ 一列打了叉，虽然这次没影响结论，但漏比较可能选错入基变量。
\end{enumerate}
\end{warningbox}

\begin{lifebox}[凑学分：先选“最省时”的，再按性价比补齐]
\begin{minipage}[t]{0.55\linewidth}
\textbf{生活场景：}毕业要求“专业选修 $\ge4$ 学分、通识 $\ge2$ 学分”，你想用最少的学习时间满足要求。最开始一门都没选：时间最少（“最优”），但学分不够（不可行）。\par
\textbf{对应关系：}学分要求 $\leftrightarrow$ $\ge$ 约束；还差的学分 $\leftrightarrow$ $b$ 列的负数；先补差得最多的一类 $\leftrightarrow$ 出基判据；在能补这一类的课程里选“每学分耗时最少”的 $\leftrightarrow$ $\theta=\min\{\sigma_j/a_{rj}\}$。\par
\textbf{小估算：}专业选修还差 4 学分。课程甲 2 学分、每周 3 小时，每学分 1.5 小时；课程乙 3 学分、每周 6 小时，每学分 2 小时。按 $\theta$ 规则先选甲（$1.5<2$）。\par
\textbf{边界说明：}真实选课时一门课只能整门选、还有时间冲突，属于整数规划；对偶单纯形法的“比值选择”只描述了每一步怎样以最小代价逼近可行。
\end{minipage}\hfill
\begin{minipage}[t]{0.41\linewidth}\vspace{0pt}\centering
\begin{tikzpicture}[x=1cm,y=1cm]
  \draw[ax] (0,0)--(3.8,0) node[right,font=\scriptsize]{};
  \fill[auxblue!70] (0.5,0) rectangle (1.3,1.5); \node[font=\scriptsize,above] at (0.9,1.5) {1.5};
  \fill[auxblue!35] (2.1,0) rectangle (2.9,2.0); \node[font=\scriptsize,above] at (2.5,2.0) {2.0};
  \node[font=\scriptsize,below] at (0.9,0) {课程甲}; \node[font=\scriptsize,below] at (2.5,0) {课程乙};
  \node[font=\scriptsize,rotate=90] at (-0.35,1.1) {每学分耗时};
  \node[gbox,font=\scriptsize] at (1.7,-1.0) {比值最小者入基：选甲};
  \node[box,font=\scriptsize] at (1.7,2.95) {$\theta_j=\sigma_j/a_{rj}$};
\end{tikzpicture}
\fig{生活场景 $\leftrightarrow$ 入基比值}
\end{minipage}
\end{lifebox}

\begin{checkbox}
\begin{enumerate}[leftmargin=1.5em]
  \item 什么叫原始可行、对偶可行？单纯形法和对偶单纯形法各保住哪一条？
  \item 对偶单纯形法的出基、入基规则分别是什么？为什么主元必须为负？
  \item 推导入基比值 $\theta=\min\{\sigma_j/a_{rj}:a_{rj}<0\}$。
  \item 出基行没有负系数时说明什么？为什么？
  \item 什么样的模型适合直接用对偶单纯形法求解？
  \item “对偶单纯形法就是求解对偶问题的单纯形法”，这句话错在哪里？
\end{enumerate}
\end{checkbox}

\paragraph{到这里，我们已经解决了什么？}我们有了第二把“钥匙”：一张检验数合格、解不可行的表，可以用对偶单纯形法几步修好，不必加人工变量。
\paragraph{下一步还需要解决什么？}原问题的参数（资源量、价格、工艺、新产品、新约束）改变之后，怎样判断最优表还能不能用？如果不能用，该用单纯形法还是对偶单纯形法续算？这就是下一章的灵敏度分析。

\begin{figure}[!htbp]\centering
\begin{tikzpicture}[node distance=0.35cm and 0.4cm,k/.style={box,minimum height=0.8cm,font=\small}]
  \node[k] (a) {$b$ 有负、$\sigma\le0$};
  \node[k,right=of a] (b) {对偶可行};
  \node[k,right=of b] (c) {出基：$b_r$ 最负};
  \node[k,right=of c] (d) {入基：$\min\sigma_j/a_{rj}$};
  \node[k,below=0.55cm of c] (e) {旋转（同单纯形法）};
  \node[gbox,below=0.55cm of d] (f) {$b\ge0$ 时最优};
  \node[gbox,below=0.55cm of a] (g) {对偶解释：\\与原问题表一一对应};
  \draw[arr] (a)--(b); \draw[arr] (b)--(c); \draw[arr] (c)--(d); \draw[arr] (d)--(e); \draw[arr] (e)--(f); \draw[arr] (b)--(g);
\end{tikzpicture}
\fig{本章知识地图}
\end{figure}

% ============================================================
%  第 6 章  灵敏度分析（K7、K8、K9）
% ============================================================
\section{灵敏度分析：数据变了，最优解还成立吗}

前面几章已经把一个线性规划解到最优：最优解、最优值、最优基都已得出。然而模型里的每一个数都只是“当前状态的估计值”——资源可用量会变，产品利润会变，工艺改进会改变消耗定额，市场上还会出现新的产品和新的限制。本章要回答的核心问题是（课件 P2）：\textbf{参数改动之后，原先算得的最优解是否仍然最优？}

\noindent\textbf{本章回答五个问题}（课件 P3）：
\begin{enumerate*}[label=(\arabic*)]
  \item 资源边际贡献：某种资源多给一个单位，最优值增加多少？（第 4 章已答：影子价格）
  \item 资源量的稳定范围：资源可用量在什么范围内变动，最优基保持不变？
  \item 利润系数的稳定范围：某个产品的单位利润在什么范围内变动，最优解保持不变？
  \item 新增产品：新增一个产品，是否值得投产？
  \item 新增约束：新增一条限制，最优解如何调整？
\end{enumerate*}\par
\noindent\textbf{本章知识路线：}三种情形 $\to$ 两条判据 $\to$ 四个步骤 $\to$ 右端项 $b$ $\to$ 价值系数 $c$ $\to$ 新增变量 $\to$ 新增约束。全章使用第 4 章的主例与表~\ref{tab:lec}。

\subsection{总体思路：只核验两条判据}

\begin{theorembox}[结论 6.1　最优表成立的两条判据 \srcTag{课件 P46--P48}]
一个基 $B$ 给出最优表，当且仅当同时满足：
\[
\text{可行性判据：}\ x_B=B^{-1}b\ge0;\qquad\qquad
\text{最优性判据：}\ \sigma\T=c\T-c_B\T B^{-1}A\le0 .
\]
可行性公式里没有 $c$，最优性公式里没有 $b$。所以\textbf{改 $b$ 只可能破坏可行性，改 $c$ 只可能破坏最优性}——两条判据可以分开核验，这是灵敏度分析能逐条单独处理的根本原因。
\end{theorembox}

\begin{figure}[!htbp]\centering
\begin{tikzpicture}[x=1cm,y=1cm,
  p/.style={box,minimum width=2.6cm,minimum height=0.75cm,font=\small},
  c1/.style={draw=auxblue,fill=grayblue,rounded corners=2pt,minimum width=2.6cm,minimum height=0.75cm,font=\small,align=center},
  o/.style={gbox,minimum width=2.9cm,minimum height=0.75cm,font=\small}]
  \node[font=\small\bfseries,navy] at (0,1.0) {改动的参数};
  \node[font=\small\bfseries,navy] at (4.4,1.0) {需核验的判据};
  \node[font=\small\bfseries,navy] at (9.2,1.0) {越界后的续解};
  \node[p] (b) at (0,0) {右端项 $b_i$};
  \node[p] (cn) at (0,-1.0) {非基变量 $c_k$};
  \node[p] (cb) at (0,-2.0) {基变量 $c_j$};
  \node[p] (nv) at (0,-3.0) {新增变量};
  \node[p] (nc) at (0,-4.0) {新增约束};
  \node[c1] (fe) at (4.4,-0.5) {可行性\\$B^{-1}b\ge0$};
  \node[c1] (op) at (4.4,-2.5) {最优性\\$\sigma\le0$};
  \node[o] (ds) at (9.2,-0.5) {对偶单纯形法\\{\scriptsize（$b$ 列出现负数）}};
  \node[o] (sm) at (9.2,-2.5) {单纯形法\\{\scriptsize（$\sigma$ 行出现正数）}};
  \draw[arr] (b.east)--(fe.west); \draw[arr] (nc.east)--(fe.west);
  \draw[arr] (cn.east)--(op.west); \draw[arr] (cb.east)--(op.west); \draw[arr] (nv.east)--(op.west);
  \draw[arr] (fe)--(ds); \draw[arr] (op)--(sm);
  \node[rbox,font=\scriptsize,text width=3.0cm] at (9.2,-4.0) {两条同时不满足：\\引入人工变量重解};
\end{tikzpicture}
\fig{灵敏度分析的判定图（课件 P47、P50）：每类改动只需核验一条判据；越界后按缺失的判据选择续解算法}\label{fig:satree}
\end{figure}

图~\ref{fig:satree} 从左往右读：先看改了哪类参数，再看它会影响哪条判据，最后看判据失效时用什么算法补救。改动之后最终只会落入三种情形之一（课件 P45）：
\begin{enumerate}[label=\textbf{情形\arabic*}，,leftmargin=4em]
  \item 最优基与最优解都不变，原最优表照用；
  \item 最优基不变，但基变量的取值改变，按新右端项重算；
  \item 最优基改变，原表失效，需继续迭代（对偶单纯形法或单纯形法）。
\end{enumerate}

\begin{summarybox}[灵敏度分析的四个步骤（课件 P49）]
\begin{tikzpicture}[x=0.97cm,y=1cm,s/.style={box,text width=3.1cm,minimum height=1.25cm,font=\small}]
  \node[s] (a) at (0,0) {\textbf{① 写出 $\Delta$ 形式}\\{\scriptsize 改了哪个参数、改多少}};
  \node[s] (b) at (4.1,0) {\textbf{② 追踪传递路径}\\{\scriptsize 经 $B^{-1}$ 传到 $b$ 列还是 $\sigma$ 行}};
  \node[s] (c) at (8.2,0) {\textbf{③ 建立不等式组}\\{\scriptsize 改动后仍满足判据}};
  \node[s] (d) at (12.3,0) {\textbf{④ 求解集得范围}\\{\scriptsize 越界则按缺失判据续解}};
  \draw[arr] (a)--(b); \draw[arr] (b)--(c); \draw[arr] (c)--(d);
\end{tikzpicture}
\end{summarybox}

\subsection{右端项 $b$ 的变化}

\begin{proofbox}[推导：改动怎样传到 $b$ 列 \srcTag{课件 P52--P53}]
设第 $r$ 项右端由 $b_r$ 改为 $b_r+\Delta b_r$，其余不变，即 $\Delta b=\Delta b_r\,e_r$。最优基不变时
\[
x_B'=B^{-1}(b+\Delta b)=B^{-1}b+\Delta b_r\,B^{-1}e_r=x_B+\Delta b_r\,\beta_r ,
\]
其中 $\beta_r=B^{-1}e_r$ 是 $B^{-1}$ 的\textbf{第 $r$ 列}，也就是最优表中第 $r$ 个约束的松弛变量所在的列。检验数公式中不含 $b$，所以 $\sigma$ 行不变，最优性自动保住，只需核验可行性：
\[
\boxed{\;x_{B_i}+\Delta b_r\,\beta_{ir}\ge0,\qquad i=1,\dots,m\;}
\]
逐行解这个不等式：$\beta_{ir}>0$ 给出\textbf{下界} $\Delta b_r\ge-x_{B_i}/\beta_{ir}$；$\beta_{ir}<0$ 给出\textbf{上界} $\Delta b_r\le-x_{B_i}/\beta_{ir}$；$\beta_{ir}=0$ 不构成条件。所有下界取最大、所有上界取最小，即得允许范围。
\end{proofbox}

\begin{examplebox}[例 6.1　材料 $A$ 的允许变动范围 \srcTag{课件 P54--P57}]
\textbf{已知：}$\beta_1=(\frac12,-\frac18,\frac14)\T$（最优表 $x_3$ 列），$x_B=(4,2,4)\T$（基变量次序 $x_4,x_2,x_1$）。\quad\textbf{求：}$b_1$ 的允许范围。\par
\textbf{逐行列出并求解：}
\[
\begin{array}{lll}
x_4\text{ 行：} & 4+\frac12\Delta b_1\ge0 & \Rightarrow\ \Delta b_1\ge-8\quad(\text{系数为正，下界})\\[2pt]
x_2\text{ 行：} & 2-\frac18\Delta b_1\ge0 & \Rightarrow\ \Delta b_1\le16\quad(\text{系数为负，上界})\\[2pt]
x_1\text{ 行：} & 4+\frac14\Delta b_1\ge0 & \Rightarrow\ \Delta b_1\ge-16\quad(\text{下界，被第一行覆盖})
\end{array}
\]
\textbf{结果：}$-8\le\Delta b_1\le16$，即 $b_1\in[8,32]$。在此范围内最优基 $\{x_4,x_2,x_1\}$ 不变，材料 $A$ 的影子价格恒为 $\frac18$；最优解按 $x_B'=x_B+\Delta b_1\beta_1$ 重算，最优值按 $z^*+\frac18\Delta b_1$ 更新。
\end{examplebox}

\begin{figure}[!htbp]\centering
\begin{tikzpicture}[x=0.36cm,y=1cm]
  \draw[ax] (-20,0)--(21,0) node[right,font=\scriptsize]{$\Delta b_1$};
  \foreach \x in {-16,-8,0,16} {\draw (\x,0.08)--(\x,-0.08); \node[below,font=\tiny] at (\x,-0.08) {\x};}
  \draw[greenline,very thick,-{Stealth}] (-8,0.45)--(20.5,0.45); \node[font=\scriptsize,left] at (-8.3,0.45) {$x_4$ 行：$\ge-8$};
  \draw[redline,very thick,-{Stealth}] (16,0.85)--(-19.5,0.85); \node[font=\scriptsize,right] at (16.3,0.85) {$x_2$ 行：$\le16$};
  \draw[greenline!60,very thick,-{Stealth}] (-16,1.25)--(20.5,1.25); \node[font=\scriptsize,left] at (-16.3,1.25) {$x_1$ 行：$\ge-16$};
  \fill[shiftc!35] (-8,-0.55) rectangle (16,-0.35);
  \node[font=\scriptsize,shiftc!80!black] at (4,-0.85) {交集 $[-8,16]$，即 $b_1\in[8,32]$};
\end{tikzpicture}
\fig{例 6.1 的不等式在数轴上：三条射线的公共部分就是允许范围}\label{fig:db1}
\end{figure}

\begin{figure}[!htbp]\centering
\begin{tikzpicture}[x=0.55cm,y=0.55cm]
  \begin{scope}
    \lecRegion
    \axes{9}{5}{x_1}{x_2}
    \draw[con] (-0.2,3)--(9.3,3);
    \draw[con] (-0.4,4.2)--(9.3,-0.65);
    \draw[con,black!35] (4,-0.2)--(4,5.2);
    \draw[shiftl] (2,-0.2)--(2,5.2) node[above,font=\scriptsize]{$b_1=8$};
    \draw[shiftl] (8,-0.2)--(8,5.2) node[above,font=\scriptsize]{$b_1=32$};
    \node[optpt,fill=shiftc] at (2,3) {}; \node[optpt,fill=shiftc] at (8,0) {}; \node[optpt] at (4,2) {};
    \node[lab,above left] at (2,3) {$(2,3)$}; \node[lab,above right] at (8,0) {$(8,0)$}; \node[lab,right] at (4.1,2.2) {$(4,2)$};
    \draw[arr,optc,thick] (2.1,2.95)--(3.85,2.07); \draw[arr,optc,thick] (4.15,1.93)--(7.85,0.07);
    \node[font=\scriptsize,align=center] at (4.5,-1.5) {(a) 材料 $A$：$b_1\in[8,32]$\\最优点沿台时斜线从 $(2,3)$ 滑到 $(8,0)$};
  \end{scope}
  \begin{scope}[xshift=7.2cm]
    \lecRegion
    \axes{9}{5}{x_1}{x_2}
    \draw[con] (4,-0.2)--(4,5.2);
    \draw[con] (-0.2,3)--(9.3,3);
    \draw[con,black!35] (-0.4,4.2)--(9.3,-0.65);
    \draw[shiftl] (-0.4,2.2)--(5.4,-0.7) node[right,font=\scriptsize]{$b_3=4$};
    \draw[shiftl] (-0.4,5.2)--(9.3,0.35) node[above,font=\scriptsize,pos=0.85]{$b_3=10$};
    \node[optpt,fill=shiftc] at (4,0) {}; \node[optpt,fill=shiftc] at (4,3) {}; \node[optpt] at (4,2) {};
    \node[lab,below right] at (4,0) {$(4,0)$}; \node[lab,above right] at (4,3) {$(4,3)$}; \node[lab,right] at (4.1,1.8) {$(4,2)$};
    \node[font=\scriptsize,align=center] at (4.5,-1.5) {(b) 台时：$b_3\in[4,10]$\\最优点沿材料 $A$ 竖线从 $(4,0)$ 升到 $(4,3)$};
  \end{scope}
\end{tikzpicture}
\fig{右端项允许范围的几何直观（课件 P57--P58）：两条橙色虚线是范围的两个端点，端点之间最优顶点始终由同一对约束交出}\label{fig:brange}
\end{figure}

图~\ref{fig:brange}(a) 中，材料 $A$ 的边界（竖线）在 $x_1=2$ 与 $x_1=8$ 之间平移时，最优顶点始终是“材料 $A$ 线与台时线的交点”，只是沿台时斜线滑动；到 $b_1=8$ 时滑到 $(2,3)$，碰上材料 $B$ 的边界；到 $b_1=32$ 时滑到 $(8,0)$，碰上 $x_2=0$。再往外，交点所在的约束组合就变了，最优基随之改变。图 (b) 同理：台时线在 $b_3=4$ 与 $b_3=10$ 之间平移，最优顶点沿材料 $A$ 竖线从 $(4,0)$ 升到 $(4,3)$。

对另外两项右端作同样计算（课件 P56、P59）：$\beta_3=(-2,\frac12,0)\T$ 给出 $4-2\Delta b_3\ge0$、$2+\frac12\Delta b_3\ge0$，即 $b_3\in[4,10]$；$\beta_2=(1,0,0)\T$ 只有第一行 $4+\Delta b_2\ge0$ 起作用，上界不存在，$b_2\in[8,+\infty)$——材料 $B$ 本来就有 4 kg 富余，再多也不改变最优基，只有减少超过 4 kg 才会使它变为紧约束。

\begin{table}[!htbp]\centering
\small
\begin{tabular}{@{}lcccl@{}}
\toprule
\rowcolor{lightblue}资源 & 现值 & 允许范围 & 影子价格 & 含义\\\midrule
材料 $A$ & 16 & $[8,32]$ & $\frac18$ & 范围内每多 1 kg 增利 $\frac18$\\
材料 $B$ & 12 & $[8,+\infty)$ & 0 & 有富余，多给无益\\
设备台时 & 8 & $[4,10]$ & $\frac32$ & 最紧缺的资源\\\bottomrule
\end{tabular}\par\nopagebreak[4]
\captionof{table}{三条右端项范围汇总（课件 P60）}\label{tab:brange}
\end{table}

\begin{examplebox}[例 6.2　越出范围：材料 $A$ 增至 36 kg \fillTag]
$\Delta b_1=20>16$，超出范围。按传递公式 $x_B'=(4,2,4)\T+20\cdot(\frac12,-\frac18,\frac14)\T=(14,\,-\frac12,\,9)\T$，$x_2$ 行变负。此时 $\sigma$ 行不变（仍 $\le0$），表处于“对偶可行而原始不可行”状态，直接用对偶单纯形法：\par
$x_2$ 行出基，该行为 $(0,1,-\frac18,0,\frac12)$，负系数只有 $x_3$ 列，$\theta=\frac{-1/8}{-1/8}=1$，$x_3$ 入基。旋转后：
\begin{center}\small
\begin{tabular}{@{}ccc|ccccc@{}}\toprule
$c_B$ & $x_B$ & $b$ & $x_1$ & $x_2$ & $x_3$ & $x_4$ & $x_5$\\\midrule
0 & $x_4$ & 12 & 0 & 4 & 0 & 1 & 0\\
0 & $x_3$ & 4 & 0 & $-8$ & 1 & 0 & $-4$\\
2 & $x_1$ & 8 & 1 & 2 & 0 & 0 & 1\\\midrule
\multicolumn{3}{c|}{$\sigma_j$} & 0 & $-1$ & 0 & 0 & $-2$\\\bottomrule
\end{tabular}
\end{center}
$b\ge0$，最优：$x^*=(8,0)$，$z^*=16$。新的影子价格为 $y=(0,0,2)$——材料 $A$ 不再紧缺（剩 4 kg），价格降为 0；与图~\ref{fig:zb1} 中 $b_1>32$ 后折线变平完全吻合。\textbf{只用了一次旋转。}
\end{examplebox}

\begin{warningbox}
\textbf{错误做法：}“多买 100 kg 材料 $A$，利润增加 $100\times\frac18=12.5$。”\quad
\textbf{正确做法：}先核对 $\Delta b_1=100$ 是否落在 $[-8,16]$ 内；超出范围时最优基已变，必须重新分析（课件 P61）。按例 6.2 的方法可知，$b_1\ge32$ 后最优值恒为 16，最多增利 2。
\end{warningbox}

\begin{lifebox}[阶梯电价：每度电的“边际价格”在档内不变]
\begin{minipage}[t]{0.55\linewidth}
\textbf{生活场景：}我国许多地区居民用电实行阶梯电价：月用电量在第一档内按较低单价计费，超过档位后，超出部分按更高单价计费。\par
\textbf{对应关系：}月用电量 $\leftrightarrow$ 资源量 $b$；当前所在档的单价 $\leftrightarrow$ 影子价格（多用 1 度的边际花费）；档位分界点 $\leftrightarrow$ 允许范围的端点。\par
\textbf{小估算}（示意数字，各地标准不同）：第一档 $0$--$200$ 度每度 0.5 元，第二档每度 0.55 元。本月已用 150 度，再多用 30 度仍在第一档，多花 $30\times0.5=15$ 元；若再多用 80 度，其中 50 度在第一档、30 度进入第二档，多花 $25+16.5=41.5$ 元，而不是 $80\times0.5=40$ 元。\par
\textbf{边界说明：}电费是“花费”，用得越多边际价格越高（折线向上弯）；生产问题的最优利润是“收益”，资源越多影子价格越低（图~\ref{fig:zb1} 的折线向下弯）。两者共同点是：\textbf{边际值只在一段范围内是常数}。
\end{minipage}\hfill
\begin{minipage}[t]{0.41\linewidth}\vspace{0pt}\centering
\begin{tikzpicture}
\begin{axis}[width=6.3cm,height=4.8cm,xmin=0,xmax=300,ymin=0,ymax=160,axis lines=left,
  xlabel={月用电量（度）},ylabel={电费（元）},xtick={0,150,200,300},ytick={0,100,155},
  label style={font=\scriptsize},tick label style={font=\scriptsize},grid=major,grid style={black!8}]
  \addplot[navy,very thick] coordinates {(0,0)(200,100)(300,155)};
  \addplot[only marks,mark=*,mark options={fill=optc,draw=white}] coordinates {(150,75)};
  \draw[shiftc,dashed] (axis cs:200,0)--(axis cs:200,100);
  \node[font=\scriptsize,anchor=south east] at (axis cs:120,62) {斜率 0.5};
  \node[font=\scriptsize,anchor=north west] at (axis cs:215,118) {斜率 0.55};
\end{axis}
\end{tikzpicture}
\fig{生活场景 $\leftrightarrow$ 允许范围}
\end{minipage}
\end{lifebox}

\subsection{价值系数 $c$ 的变化}

价值系数只出现在检验数公式中，基变量取值 $x_B=B^{-1}b$ 与 $c$ 无关。因此改动 $c_j$ 时，最优解的\textbf{取值}不变，可能改变的只有最优性（课件 P63）。要分两类讨论。

\begin{proofbox}[推导 \srcTag{课件 P64--P67}]
\textbf{(1) 非基变量 $x_k$ 的 $c_k$ 改为 $c_k+\Delta c_k$。}$c_B$ 中不含 $c_k$，$y\T=c_B\T B^{-1}$ 不变，其余各列的检验数不变，只有 $x_k$ 自身：
\[
\sigma_k'=(c_k+\Delta c_k)-y\T P_k=\sigma_k+\Delta c_k\le0\quad\Longleftrightarrow\quad\boxed{\Delta c_k\le-\sigma_k}.
\]
因原来 $\sigma_k\le0$，故 $-\sigma_k\ge0$：非基变量的价值系数只有上界，向下没有限制。\par
\textbf{(2) 基变量 $x_j$（位于最优表第 $r$ 行）的 $c_j$ 改为 $c_j+\Delta c_j$。}$c_B$ 的第 $r$ 个分量随之改变，新对偶变量 $y^{\prime\mathsf T}=y\T+\Delta c_j\,(B^{-1})_{r\cdot}$，其中 $(B^{-1})_{r\cdot}$ 是 $B^{-1}$ 的第 $r$ \textbf{行}。对任一非基列 $k$：
\[
\sigma_k'=c_k-y^{\prime\mathsf T} P_k=\sigma_k-\Delta c_j\,(B^{-1})_{r\cdot}P_k=\sigma_k-\Delta c_j\,t_{rk},
\]
$t_{rk}$ 是最优表第 $r$ 行、第 $k$ 列的元素。要求 $\sigma_k-\Delta c_j\,t_{rk}\le0$：$t_{rk}>0$ 给出下界 $\Delta c_j\ge\sigma_k/t_{rk}$；$t_{rk}<0$ 给出上界 $\Delta c_j\le\sigma_k/t_{rk}$；$t_{rk}=0$ 不构成条件。$x_j$ 自身是基变量，检验数恒为 0，不参与。
\end{proofbox}

\begin{examplebox}[例 6.3　产品二的价值系数 $c_2$ 的范围 \srcTag{课件 P68--P69}（已改正 P69 的排版笔误）]
\textbf{已知：}$x_2$ 在最优表第 2 行（$r=2$），非基列为 $x_3,x_5$：$t_{23}=-\frac18$，$t_{25}=\frac12$，$\sigma_3=-\frac18$，$\sigma_5=-\frac32$。\par
\textbf{列式求解：}
\[
\begin{array}{lll}
x_3\text{ 列：} & -\frac18-\Delta c_2\cdot(-\frac18)\le0 & \Rightarrow\ \Delta c_2\le1\quad(t_{23}<0\text{，上界})\\[2pt]
x_5\text{ 列：} & -\frac32-\Delta c_2\cdot\frac12\le0 & \Rightarrow\ \Delta c_2\ge-3\quad(t_{25}>0\text{，下界})
\end{array}
\]
\textbf{结果：}$-3\le\Delta c_2\le1$，即 $c_2\in[0,4]$。在此范围内最优解仍为 $(4,2)$，最优值按新的 $c_2$ 重算 $z^*=8+2c_2$。\par
同理，产品一（第 3 行）：$t_{33}=\frac14$ 给出 $-\frac18-\frac14\Delta c_1\le0\Rightarrow\Delta c_1\ge-\frac12$；$t_{35}=0$ 不构成条件。故 $c_1\in[\frac32,+\infty)$（课件 P71）。
\end{examplebox}

\begin{figure}[!htbp]\centering
\begin{tikzpicture}[x=0.62cm,y=0.62cm]
  \lecRegion
  \axes{9}{5}{x_1}{x_2}
  \draw[con,black!40] (-0.2,3)--(9.3,3);
  \draw[con,black!40] (4,-0.2)--(4,5.2);
  \draw[con,black!40] (-0.4,4.2)--(9.3,-0.65);
  % c2=3 当前等值线 2x1+3x2=14
  \draw[iso,black] (0,4.667)--(7,0) node[below,font=\scriptsize]{$c_2=3$（当前）};
  % c2=0: 2x1=8 竖线 x1=4
  \draw[very thick,dashed,redline] (4.0,-0.3)--(4.0,5.3) node[above,font=\scriptsize]{$c_2=0$};
  % c2=4: 2x1+4x2=16 与台时平行
  \draw[very thick,dashed,shiftc] (-0.4,4.2)--(8.5,-0.25);
  \node[font=\scriptsize,shiftc!80!black,anchor=west] at (5.6,2.0) {$c_2=4$：与台时线重合};
  \fill[shiftc!50,opacity=0.6] (2,3)--(4,2)--(4.1,2.15)--(2.1,3.15)--cycle;
  \fill[redline!40,opacity=0.6] (3.9,0)--(4.1,0)--(4.1,2)--(3.9,2)--cycle;
  \node[optpt] at (4,2) {};
  \node[vpt,fill=shiftc] at (2,3) {}; \node[vpt,fill=redline] at (4,0) {};
  \node[lab,left] at (1.9,3.2) {$(2,3)$}; \node[lab,below right] at (4,0) {$(4,0)$};
  \draw[arr,thick,optc] (6.4,3.8) arc[start angle=10,end angle=60,radius=1.6];
  \node[font=\scriptsize,optc] at (7.4,4.6) {$c_2$ 增大，等值线变缓};
\end{tikzpicture}
\fig{价值系数范围的几何直观（课件 P70--P71）：等值线随 $c_2$ 旋转，在两个端点处与某条边界重合，此时整条棱上的点都是最优解}\label{fig:c2}
\end{figure}

图~\ref{fig:c2} 中黑色虚线是当前等值线 $2x_1+3x_2=14$。$c_2$ 变化时等值线绕最优点旋转：$c_2$ 减到 0，等值线变成竖直线，与材料 $A$ 的边界 $x_1=4$ 重合，$(4,0)$ 到 $(4,2)$ 整条竖棱（红色）都是最优解；$c_2$ 增到 4，等值线与台时线 $x_1+2x_2=8$ 平行，$(4,2)$ 到 $(2,3)$ 整条斜棱（橙色）都是最优解。$c_2$ 在 $[0,4]$ 内，等值线的斜率夹在两条边界线的斜率之间，最优顶点始终是 $(4,2)$。产品一的范围无上界：$x_1$ 已被材料 $A$ 顶到上限 4，利润再高也没有资源再扩产。

\begin{examplebox}[例 6.4　越出范围：$c_2$ 增至 5 \fillTag]
$\Delta c_2=2>1$。按公式：$\sigma_3'=-\frac18-2\cdot(-\frac18)=\frac18>0$，$\sigma_5'=-\frac32-2\cdot\frac12=-\frac52$。$b$ 列不变（仍可行），$\sigma$ 行出现正数，表处于“原始可行而未达最优”状态，续走单纯形法（课件 P73）：$x_3$ 入基，$x_3$ 列为 $(\frac12,-\frac18,\frac14)\T$，正系数行的比值 $x_4$：$\frac{4}{1/2}=8$，$x_1$：$\frac{4}{1/4}=16$，$x_4$ 出基。旋转后 $x_3=8$，$x_2=2+\frac18\cdot8=3$，$x_1=4-\frac14\cdot8=2$，即最优点移到 $(2,3)$，$z^*=2\cdot2+5\cdot3=19$。这与图~\ref{fig:c2} 的预言一致：$c_2>4$ 后最优顶点跳到相邻顶点 $(2,3)$。
\end{examplebox}

\begin{lifebox}[地图导航：小幅拥堵不改路线]
\begin{minipage}[t]{0.55\linewidth}
\textbf{生活场景：}导航推荐你走高架（全程 20 分钟），备选的地面道路全程 24 分钟。高架上突然轻微拥堵，导航并没有立刻改道。\par
\textbf{对应关系：}各路段用时 $\leftrightarrow$ 价值系数 $c_j$；当前推荐路线 $\leftrightarrow$ 最优基；“改走地面道路”的临界拥堵时间 $\leftrightarrow$ 价值系数允许范围的端点。最短路问题本身就可以写成线性规划。\par
\textbf{小估算：}高架上那一段原来用时 8 分钟。只要这一段的用时增加不超过 $24-20=4$ 分钟（即不超过 12 分钟），高架路线仍然最快；超过 12 分钟，最优路线就会切换——这就是“范围的端点”。\par
\textbf{边界说明：}真实导航基于预测路况，并考虑红绿灯、换道等因素，还会避免频繁改道；线性模型只刻画“哪条更快”。
\end{minipage}\hfill
\begin{minipage}[t]{0.41\linewidth}\vspace{0pt}\centering
\begin{tikzpicture}[x=1cm,y=1cm]
  \node[circle,fill=navy,inner sep=2pt,label=left:{\scriptsize 出发}] (s) at (0,0) {};
  \node[circle,fill=navy,inner sep=2pt,label=right:{\scriptsize 到达}] (t) at (4,0) {};
  \draw[very thick,auxblue] (s) to[bend left=35] node[above,font=\scriptsize,align=center]{高架 20 分\\（其中一段 8 分）} (t);
  \draw[very thick,black!40] (s) to[bend right=35] node[below,font=\scriptsize]{地面 24 分} (t);
  \draw[ax] (0,-1.9)--(4.2,-1.9) node[right,font=\tiny]{该段用时};
  \fill[greenline!50] (0.4,-1.95) rectangle (2.4,-1.85);
  \node[font=\tiny,below] at (0.4,-1.95) {8}; \node[font=\tiny,below] at (2.4,-1.95) {12};
  \node[font=\scriptsize,greenline] at (1.4,-1.55) {路线不变};
\end{tikzpicture}
\fig{生活场景 $\leftrightarrow$ $c_j$ 的范围}
\end{minipage}
\end{lifebox}

\subsection{新增一个变量（新产品）}

\begin{theorembox}[结论 6.2　新产品是否值得生产 \fillTag]
新增变量 $x_{n+1}$，价值系数 $c_{n+1}$，消耗列 $P_{n+1}$。它进入模型时取值为 0，原最优解仍可行，$B^{-1}$ 也不变，只需计算它的检验数：
\[
\boxed{\;\sigma_{n+1}=c_{n+1}-y\T P_{n+1}\;}
\]
若 $\sigma_{n+1}\le0$，不值得生产，原最优解不变；若 $\sigma_{n+1}>0$，值得生产，把 $B^{-1}P_{n+1}$ 作为新列写进最优表，继续单纯形法。\par
\textbf{经济含义：}$y\T P_{n+1}$ 是新产品占用的资源按影子价格折算的“机会成本”。利润超过机会成本才值得生产。
\end{theorembox}

\begin{examplebox}[例 6.5　主例新增产品三]
产品三每件耗材料 $A$ 1 kg、材料 $B$ 0、台时 1，单位利润 2。$P_6=(1,0,1)\T$，$y=(\frac18,0,\frac32)$：
\[
y\T P_6=\tfrac18\cdot1+0\cdot0+\tfrac32\cdot1=\tfrac{13}{8},\qquad \sigma_6=2-\tfrac{13}{8}=\tfrac38>0 .
\]
值得生产。$B^{-1}P_6=(\frac12-2,\,-\frac18+\frac12,\,\frac14)\T=(-\frac32,\frac38,\frac14)\T$，比值 $x_2$：$\frac{2}{3/8}=\frac{16}{3}$，$x_1$：$\frac{4}{1/4}=16$，$x_2$ 出基。旋转后 $x_6=\frac{16}{3}$，$x_1=4-\frac14\cdot\frac{16}{3}=\frac83$，$x_4=4+\frac32\cdot\frac{16}{3}=12$，$z^*=2\cdot\frac83+2\cdot\frac{16}{3}=16$。新的对偶解 $y=(0,0,2)$ 使 $\sigma_2=3-2\cdot2=-1\le0$，确为最优。\par
\textbf{本题真正想说明什么：}判断“是否值得生产”只需一次内积，不必重解整个问题。
\end{examplebox}

\begin{figure}[!htbp]\centering
\begin{tikzpicture}[x=1cm,y=1cm]
  \draw[ax] (0,0)--(7.6,0);
  \fill[auxblue!60] (0.6,0) rectangle (1.6,0.25); \node[font=\scriptsize,above] at (1.1,0.25) {材料 $A$ $\frac18$};
  \fill[auxblue!60] (0.6,0.25) rectangle (1.6,3.25); \node[font=\scriptsize] at (1.1,1.75) {台时 $\frac32$};
  \node[font=\scriptsize,below,align=center] at (1.1,0) {机会成本\\$y\T P_6=\frac{13}{8}$};
  \fill[greenline!70] (2.6,0) rectangle (3.6,4); \node[font=\scriptsize,below,align=center] at (3.1,0) {利润 2\\（$\sigma_6=\frac38>0$）};
  \node[gbox,font=\scriptsize] at (3.1,4.5) {值得生产};
  \fill[redline!60] (5.2,0) rectangle (6.2,3.0); \node[font=\scriptsize,below,align=center] at (5.7,0) {若利润 1.5\\（$\sigma_6<0$）};
  \node[rbox,font=\scriptsize] at (5.7,3.5) {不值得};
  \draw[dashed,black!50] (0.4,3.25)--(7.0,3.25) node[right,font=\scriptsize]{$\frac{13}{8}$};
\end{tikzpicture}
\fig{新增变量的判断：利润柱是否高过机会成本线（纵轴每格 0.5 元）}\label{fig:newvar}
\end{figure}

图~\ref{fig:newvar} 中左边的叠柱是产品三占用资源的机会成本：材料 $A$ 只贡献 $\frac18$，台时贡献 $\frac32$，合计 $\frac{13}{8}$（虚线）。绿柱是利润 2，高过虚线，所以值得生产；若利润只有 1.5（红柱），低于虚线，不值得。修改\textbf{非基变量的技术系数}（某个现在不生产的产品改进了工艺）与此完全相同：把新的 $P_k$ 代入 $\sigma_k=c_k-y\T P_k$ 即可；修改\textbf{基变量的技术系数}会同时改变 $B$ 与 $B^{-1}$，两条判据都要核验，是最复杂的情形（课件 P47），本课程一般不要求。

\subsection{新增一条约束}

\begin{theorembox}[结论 6.3　新增约束的处理步骤 \srcTag{课件《对偶单纯形法》P73--P83}]
\begin{enumerate}[label=\textbf{第\arabic*步}，,leftmargin=4.2em]
  \item 把原最优解代入新约束。若满足，原最优解仍最优，结束；
  \item 若不满足，给新约束加松弛变量，作为新的一行添入最优表；
  \item 用表中基变量所在行把新行里基变量的系数消为 0（使基列重新成为单位阵）；
  \item 此时新行右端为负，检验数不变（$\le0$）：表对偶可行而原始不可行，用对偶单纯形法续解。
\end{enumerate}
\end{theorembox}

\begin{examplebox}[例 6.6　产品搭配题新增约束 $x_2\le500$ \srcTag{课件 P73--P82}]
第 5 章例 5.1 的原问题（$\max z=10x_1+14x_2+12x_3$）最优表为：
\begin{center}\small
\begin{tabular}{@{}ccc|ccccc@{}}\toprule
$c_B$ & $x_B$ & $b$ & $x_1$ & $x_2$ & $x_3$ & $x_4$ & $x_5$\\\midrule
12 & $x_3$ & $\frac{3000}{11}$ & $-\frac{6}{11}$ & 0 & 1 & $\frac{4}{11}$ & $-\frac{5}{11}$\\
14 & $x_2$ & $\frac{20000}{33}$ & $\frac{70}{33}$ & 1 & 0 & $-\frac{10}{33}$ & $\frac{40}{33}$\\\midrule
\multicolumn{3}{c|}{$\sigma_j$} & $-\frac{434}{33}$ & 0 & 0 & $-\frac{4}{33}$ & $-\frac{380}{33}$\\\bottomrule
\end{tabular}
\end{center}
\textbf{第 1 步：}$x_2=\frac{20000}{33}\approx606.06>500$，不满足。\par
\textbf{第 2--3 步：}加松弛变量 $x_6$：$x_2+x_6=500$，即行 $(0,1,0,0,0,1\mid500)$。$x_2$ 是基变量，减去 $x_2$ 行消去它的系数：
\[
(0,1,0,0,0,1\mid500)-(\tfrac{70}{33},1,0,-\tfrac{10}{33},\tfrac{40}{33},0\mid\tfrac{20000}{33})
=(-\tfrac{70}{33},0,0,\tfrac{10}{33},-\tfrac{40}{33},1\mid-\tfrac{3500}{33}).
\]
\textbf{第 4 步：}$x_6$ 行右端 $-\frac{3500}{33}<0$，$x_6$ 出基。负系数在 $x_1,x_5$：比值 $x_1$：$\frac{-434/33}{-70/33}=\frac{31}{5}=6.2$，$x_5$：$\frac{-380/33}{-40/33}=\frac{19}{2}=9.5$，$x_1$ 入基。旋转后得
\begin{center}\small
\begin{tabular}{@{}ccc|cccccc@{}}\toprule
$c_B$ & $x_B$ & $b$ & $x_1$ & $x_2$ & $x_3$ & $x_4$ & $x_5$ & $x_6$\\\midrule
12 & $x_3$ & 300 & 0 & 0 & 1 & $\frac27$ & $-\frac17$ & $-\frac{9}{35}$\\
14 & $x_2$ & 500 & 0 & 1 & 0 & 0 & 0 & 1\\
10 & $x_1$ & 50 & 1 & 0 & 0 & $-\frac17$ & $\frac47$ & $-\frac{33}{70}$\\\midrule
\multicolumn{3}{c|}{$\sigma_j$} & 0 & 0 & 0 & $-2$ & $-4$ & $-\frac{31}{5}$\\\bottomrule
\end{tabular}
\end{center}
$b\ge0$、$\sigma\le0$，最优：$x_1=50$，$x_2=500$，$x_3=300$，$z=500+7000+3600=11100$。\par
\textbf{检查：}$50+\frac32\cdot500+4\cdot300=2000$ ✓，$2\cdot50+\frac65\cdot500+300=1000$ ✓。
\end{examplebox}

\begin{figure}[!htbp]\centering
\begin{tikzpicture}
\begin{axis}[ybar,width=11cm,height=4.6cm,bar width=13pt,ymin=0,ymax=700,
  symbolic x coords={$x_1$,$x_2$,$x_3$},xtick=data,ylabel={产量},
  legend style={font=\scriptsize,at={(0.5,1.03)},anchor=south,legend columns=2},
  label style={font=\scriptsize},tick label style={font=\scriptsize},
  nodes near coords,every node near coord/.append style={font=\tiny},enlarge x limits=0.25]
  \addplot[fill=auxblue!60,draw=none] coordinates {($x_1$,0) ($x_2$,606.06) ($x_3$,272.73)};
  \addplot[fill=shiftc!70,draw=none] coordinates {($x_1$,50) ($x_2$,500) ($x_3$,300)};
  \legend{新增约束前（$z\approx11757.58$）,新增约束后（$z=11100$）}
  \draw[redline,dashed,thick] (axis cs:$x_1$,500) -- (axis cs:$x_3$,500);
\end{axis}
\end{tikzpicture}
\fig{新增约束前后的生产方案对比（课件 P82）：$x_2$ 被压到上限 500，腾出的资源转给 $x_3$ 和原来不生产的 $x_1$}\label{fig:newcon}
\end{figure}

图~\ref{fig:newcon} 中红色虚线是新约束 $x_2\le500$。约束迫使 $x_2$ 从 606 降到 500，腾出的资源一部分给 $x_3$（272.7 $\to$ 300），一部分让原来不生产的 $x_1$ 开始生产（0 $\to$ 50）。总利润从约 11757.58 降到 11100——\textbf{加约束只会让可行域变小，最优值不会变好}。整个过程只做了一次旋转，而重新建模求解要从头算起。

\begin{lifebox}[航空公司新规：随身行李限重 7 kg]
\begin{minipage}[t]{0.55\linewidth}
\textbf{生活场景：}你已经按“最想带的东西”收拾好随身包，共 8 kg。登机前得知新规：随身行李不得超过 7 kg。\par
\textbf{对应关系：}新规 $\leftrightarrow$ 新增约束；原装包方案 $\leftrightarrow$ 原最优解（代入新约束后超重 1 kg，即新行右端为负）；从当前方案出发调整，而不是全部倒出来重装 $\leftrightarrow$ 对偶单纯形法从旧表续解；拿出“每千克价值最低”的东西 $\leftrightarrow$ 比值最小者入基。\par
\textbf{小判断：}可拿出的有：充电宝 0.5 kg（很重要）、两本书 1 kg（每千克价值较低）、外套 1.2 kg（可以穿在身上）。按每千克价值损失最小的原则，先考虑把外套穿上或把书放进托运。\par
\textbf{边界说明：}物品不能拆分（整数），且“价值”是主观的；对偶单纯形法只描述“以最小代价恢复可行”的思想。
\end{minipage}\hfill
\begin{minipage}[t]{0.41\linewidth}\vspace{0pt}\centering
\begin{tikzpicture}[x=1cm,y=1cm]
  \draw[ax] (0,0)--(3.6,0);
  \fill[auxblue!60] (0.5,0) rectangle (1.3,3.2); \node[font=\scriptsize,above] at (0.9,3.2) {8 kg};
  \fill[shiftc!70] (2.1,0) rectangle (2.9,2.8); \node[font=\scriptsize,above] at (2.5,2.8) {7 kg};
  \draw[redline,dashed,thick] (0.2,2.8)--(3.4,2.8) node[right,font=\scriptsize]{新规};
  \node[font=\scriptsize,below] at (0.9,0) {原方案}; \node[font=\scriptsize,below] at (2.5,0) {调整后};
  \node[font=\scriptsize,rbox] at (0.9,3.9) {超 1 kg};
\end{tikzpicture}
\fig{生活场景 $\leftrightarrow$ 新增约束}
\end{minipage}
\end{lifebox}

\begin{table}[!htbp]\centering
\small
\begin{tabularx}{\linewidth}{@{}>{\bfseries}l>{\raggedright\arraybackslash}X>{\raggedright\arraybackslash}X>{\raggedright\arraybackslash}X@{}}
\toprule
\rowcolor{lightblue}比较维度 & 右端项 $b$ 的范围 & 价值系数 $c$ 的范围 & 新增约束\\\midrule
改动传到 & $b$ 列：$x_B+\Delta b_r\beta_r$ & $\sigma$ 行 & 新增一行\\
用到的数 & $B^{-1}$ 的第 $r$ \textbf{列} & 最优表第 $r$ \textbf{行}（基变量）或本列（非基） & 基变量所在行\\
几何图像 & 边界线平移（图~\ref{fig:brange}） & 等值线旋转（图~\ref{fig:c2}） & 切去一块可行域\\
范围内 & 最优基不变，解的取值变 & 最优解不变，最优值变 & ---\\
越界后 & 对偶单纯形法 & 单纯形法 & 对偶单纯形法\\
常见错误 & 用最优表的 $b$ 代替原 $b$ & 基变量误用“只改本列”的公式 & 忘记消去新行中的基变量系数\\\bottomrule
\end{tabularx}\par\nopagebreak[4]
\captionof{table}{相近概念辨析：三类灵敏度分析}\label{tab:sacmp}
\end{table}

\begin{checkbox}
\begin{enumerate}[leftmargin=1.5em]
  \item 为什么改 $b$ 只需核验可行性、改 $c$ 只需核验最优性？
  \item 右端项范围的上下界分别来自 $\beta_r$ 中什么符号的分量？
  \item 非基变量和基变量的价值系数范围，公式有何不同？为什么？
  \item 新产品是否值得生产，怎样一步判断？$y\T P_{n+1}$ 的经济含义是什么？
  \item 新增约束后为什么一定用对偶单纯形法续解？
  \item 用图~\ref{fig:c2} 解释 $c_2$ 范围的两个端点。
\end{enumerate}
\end{checkbox}

\paragraph{到这里，我们已经解决了什么？}我们能对资源量、价格、新产品、新约束的改动做出完整判断，并在原表失效时用一两次迭代恢复最优。
\paragraph{整个专题解决了什么？}从第 1 章的 $B^{-1}$ 出发，我们依次得到：修正单纯形法（只维护 $B^{-1}$）、对偶问题与对偶定理（$y\T=c_B\T B^{-1}$ 是对偶最优解）、影子价格（$y_i=\partial z^*/\partial b_i$）、对偶单纯形法（保住 $\sigma\le0$ 找可行）、灵敏度分析（两条判据分开核验）。下一章把这些内容汇总成便于复习的总表。

\begin{figure}[!htbp]\centering
\begin{tikzpicture}[node distance=0.35cm and 0.4cm,k/.style={box,minimum height=0.8cm,font=\small}]
  \node[k] (a) {两条判据};
  \node[k,right=of a] (b) {改 $b$：可行性};
  \node[k,right=of b] (c) {$x_B+\Delta b\beta_r\ge0$};
  \node[gbox,right=of c] (d) {越界：对偶单纯形};
  \node[k,below=0.5cm of b] (e) {改 $c$：最优性};
  \node[k,right=of e] (f) {$\sigma_k-\Delta c\,t_{rk}\le0$};
  \node[gbox,right=of f] (g) {越界：单纯形法};
  \node[k,below=0.5cm of e] (h) {新增变量 / 约束};
  \node[k,right=of h] (i) {$\sigma_{n+1}=c-y\T P$ / 新行};
  \draw[arr] (a)--(b); \draw[arr] (b)--(c); \draw[arr] (c)--(d);
  \draw[arr] (a)|-(e); \draw[arr] (e)--(f); \draw[arr] (f)--(g);
  \draw[arr] (a)|-(h); \draw[arr] (h)--(i);
\end{tikzpicture}
\fig{本章知识地图}
\end{figure}

% ============================================================
%  第 7 章  总结与复习
% ============================================================
\section{总结：把六章内容收拢成一张网}

六章内容看似分散，其实都围绕同一个对象 $B^{-1}$ 展开。本章不引入新知识，只把前面的结论汇总成便于复习和查阅的形式。

\subsection{全专题知识地图}

\begin{center}
\begin{forest}
  for tree={draw=navy,fill=lightblue,rounded corners=2pt,font=\small,align=center,
    grow'=east,parent anchor=east,child anchor=west,edge={auxblue,thick},
    l sep=5mm,s sep=1.2mm,inner sep=2.5pt,anchor=west,font=\footnotesize}
  [{最优基的逆 $B^{-1}$}
    [{表的矩阵表示\\$x_B=B^{-1}b$，$\sigma=c-c_B\T B^{-1}A$}
      [修正单纯形法]]
    [{单纯形乘子 $y\T=c_B\T B^{-1}$}
      [{对偶问题\\（转置；正常/反常）}
        [{弱对偶 $c\T x\le b\T y$}]
        [{强对偶 $z^*=w^*$}[互补松弛]]]
      [{影子价格\\$y_i=\partial z^*/\partial b_i$}[{三种用法；只在范围内有效}]]]
    [{两条判据\\$B^{-1}b\ge0$；$\sigma\le0$}
      [{对偶单纯形法\\保住 $\sigma\le0$}]
      [{灵敏度分析}
        [{改 $b$：$x_B+\Delta b_r\beta_r\ge0$}]
        [{改 $c$：$\sigma_k-\Delta c_j t_{rk}\le0$}]
        [{新增变量 / 约束}]]]]
\end{forest}
\fig{全专题知识地图：从根节点 $B^{-1}$ 出发的三条主干，分别通向算法、对偶理论与灵敏度分析}\label{fig:allmap}
\end{center}

\subsection{基础解题步骤总结}

\begin{center}
\small
\begin{tabularx}{\linewidth}{@{}>{\bfseries}p{2.3cm}>{\raggedright\arraybackslash}X@{}}
\toprule
\rowcolor{lightgreen}题型 & 标准步骤\\\midrule
写对偶问题 & ① 看原目标 max/min，定对偶目标；② 定查表方向（max 左 $\to$ 右）；③ 原变量的符号 $\to$ 对偶约束的方向；④ 原约束的方向 $\to$ 对偶变量的符号；⑤ 检查：变量数与约束数互换。\\
从最优表读信息 & $B^{-1}$：初始单位阵所在列；影子价格：松弛变量列检验数取负；对偶剩余变量：决策变量列检验数取负；$z^*=y\T b$ 验算。\\
互补松弛求对偶解 & ① $x_j>0\Rightarrow$ 第 $j$ 条对偶约束取等；② 资源有剩余 $\Rightarrow y_i=0$；③ 解方程组；④ 检查 $y$ 可行、$w=z$。\\
对偶单纯形法 & ① 化 max、乘 $-1$ 得单位阵；② 确认 $\sigma\le0$；③ $b$ 最负的行出基；④ 负系数列中 $\min\sigma_j/a_{rj}$ 入基；⑤ 旋转；⑥ 直到 $b\ge0$。\\
$b_r$ 的范围 & ① 读 $\beta_r$（$x_{s_r}$ 列）与 $x_B$；② 逐行 $x_{B_i}+\Delta b_r\beta_{ir}\ge0$；③ 正系数给下界、负系数给上界；④ 取交集，加回 $b_r$。\\
$c_j$ 的范围 & 非基：$\Delta c_k\le-\sigma_k$。基（第 $r$ 行）：对每个非基列 $\sigma_k-\Delta c_jt_{rk}\le0$，$t_{rk}>0$ 给下界、$<0$ 给上界，取交集。\\
新增变量 & $\sigma_{n+1}=c_{n+1}-y\T P_{n+1}$；$>0$ 则把 $B^{-1}P_{n+1}$ 写入表、续单纯形法。\\
新增约束 & ① 原解代入；② 不满足则加松弛变量成新行；③ 消去新行中基变量的系数；④ 对偶单纯形法。\\\bottomrule
\end{tabularx}\par\nopagebreak[4]
\captionof{table}{基础解题步骤总表}
\end{center}

\subsection{高频易错点总表}

\begin{center}
\small
\begin{tabularx}{\linewidth}{@{}>{\raggedright\arraybackslash}p{4.2cm}>{\raggedright\arraybackslash}X>{\raggedright\arraybackslash}p{1.6cm}@{}}
\toprule
\rowcolor{lightred}错误想法 & 正确想法 & 出处\\\midrule
$x_B=B^{-1}b$ 中用最优表的 $b$ 列 & $b$ 是原模型右端项 & 第 1、4 章\\
$B$ 的列随意排列 & 与基变量的行次序一致，$c_B$ 同序 & 第 1 章\\
写对偶时 $b$、$c$ 放反 & $b\to$ 对偶目标系数，$c\to$ 对偶约束右端 & 第 2 章\\
min 问题也从左往右查表 & min 问题从右往左查，不等号全反 & 第 2 章\\
弱对偶证明不提 $y\ge0$ & 乘不等式必须用非负数 & 第 3 章\\
对偶无可行解 $\Rightarrow$ 原问题无界 & 也可能原问题同样无可行解 & 第 3 章\\
影子价格为 0 $\Rightarrow$ 资源有剩余 & 只有反方向成立；退化时可能“用尽但价格为 0” & 第 3、4 章\\
影子价格在任何范围都可外推 & 只在最优基不变的范围内为常数 & 第 4、6 章\\
影子价格等于市场价格 & 影子价格是企业内部的机会成本 & 第 4 章\\
对偶单纯形法求的是对偶问题 & 它求解当前问题本身 & 第 5 章\\
对偶单纯形法比值用 $b/a$ & 用 $\sigma_j/a_{rj}$，且只取 $a_{rj}<0$ 的列 & 第 5 章\\
基变量 $c_j$ 改变只影响本列 & 影响全部非基列的检验数 & 第 6 章\\
改 $c$ 会改变最优解的取值 & 范围内取值不变，只有最优值变 & 第 6 章\\
新增约束直接写进表 & 先消去新行中基变量的系数 & 第 6 章\\\bottomrule
\end{tabularx}\par\nopagebreak[4]
\captionof{table}{高频易错点总表}
\end{center}

\subsection{相近概念总对比}

\begin{center}
\small
\begin{tabularx}{\linewidth}{@{}>{\bfseries\raggedright\arraybackslash}p{3.4cm}>{\raggedright\arraybackslash}X>{\raggedright\arraybackslash}X@{}}
\toprule
\rowcolor{lightblue}概念对 & 本质区别 & 判断方法\\\midrule
单纯形法 / 修正单纯形法 & 判据相同；后者只维护 $B^{-1}$ & 是否写出整张表\\
单纯形法 / 对偶单纯形法 & 保住 $b\ge0$ 找 $\sigma\le0$ / 保住 $\sigma\le0$ 找 $b\ge0$ & 看起始表哪条判据已满足\\
弱对偶 / 强对偶 & 任意可行解的不等式 / 最优解的等式 & 是否要求“最优”\\
对偶变量 / 影子价格 / 单纯形乘子 & 同一组数 $y\T=c_B\T B^{-1}$ 的三个名字 & 看语境：理论、经济、计算\\
影子价格 / 市场价格 & 内部机会成本 / 外部交易价格 & 是否依赖本企业的技术结构\\
原始可行 / 对偶可行 & $B^{-1}b\ge0$ / $\sigma\le0$ & 看 $b$ 列还是 $\sigma$ 行\\
松弛变量 / 剩余变量 & $\le$ 约束加上 / $\ge$ 约束减去 & 系数为 $+1$ 还是 $-1$\\
改 $b$ / 改 $c$ & 动可行性（边界平移）/ 动最优性（等值线旋转） & 改动出现在哪个公式里\\\bottomrule
\end{tabularx}\par\nopagebreak[4]
\captionof{table}{相近概念总对比}
\end{center}

\subsection{后续课程连接}

\begin{center}
\begin{tikzpicture}[x=1cm,y=1cm,k/.style={box,text width=2.9cm,minimum height=1.0cm,font=\small},
  n/.style={wbox,text width=3.6cm,minimum height=1.0cm,font=\scriptsize}]
  \node[k] (c) at (0,0) {本专题：\\对偶与灵敏度};
  \node[n] (a) at (-4.6,1.6) {\textbf{整数规划}\\分支定界每加一条约束，用对偶单纯形法重新优化};
  \node[n] (b) at (0,2.4) {\textbf{运输问题}\\位势法中的 $u_i,v_j$ 就是对偶变量};
  \node[n] (d) at (4.6,1.6) {\textbf{非线性规划}\\KKT 条件、拉格朗日乘子是影子价格的推广};
  \node[n] (e) at (-4.6,-1.6) {\textbf{博弈论}\\零和博弈的极小极大定理与 LP 对偶等价};
  \node[n] (f) at (0,-2.4) {\textbf{微观经济学}\\边际收益、机会成本、要素定价};
  \node[n] (g) at (4.6,-1.6) {\textbf{电力系统 / 自动控制}\\经济调度、节点电价、模型预测控制中的 LP/QP};
  \foreach \x in {a,b,d,e,f,g} \draw[arr] (c)--(\x);
\end{tikzpicture}
\fig{本专题与后续课程的联系（只列出真实存在的直接联系）}
\end{center}

\begin{connectionbox}[拓展联想：以下内容只用于建立知识联系，当前阶段不要求掌握]
对一般的约束优化问题 $\max f(x)$，$g_i(x)\le b_i$，可以构造拉格朗日函数 $L(x,y)=f(x)+\sum_iy_i\,(b_i-g_i(x))$。线性规划时 $f=c\T x$、$g=Ax$，对 $x\ge0$ 取上确界再对 $y\ge0$ 取下确界，得到的恰好就是本专题的对偶问题；$y_i$ 依然表示“第 $i$ 个约束放松一单位时最优值的变化率”。笔记 P7 旁注的“拉格朗日乘子”指的就是这件事。
\end{connectionbox}

\subsection{知识点—题目对应表}

下表说明 38 道正式练习分别检查哪个知识点。做完练习后，可以按此表找到自己的薄弱环节，回到对应章节复习。

\begin{center}
\small
\begin{tabular}{@{}lllll@{}}
\toprule
\rowcolor{lightblue}核心知识点 & 章节 & 基础巩固（1--25） & 中等偏下（26--35） & 跨学科（36--38）\\\midrule
K1 矩阵视角、修正单纯形法 & 第 1 章 & 5, 18, 22 & 35 & \\
K2 对偶问题的构造 & 第 2 章 & 1, 2, 11, 16, 21 & 34 & \\
K3 对偶的基本性质 & 第 3 章 & 3, 4, 12 & 27 & 37, 38\\
K4 互补松弛 & 第 3 章 & 7, 14 & 27, 34 & 37\\
K5 影子价格 & 第 4 章 & 6, 13, 17, 23 & 26, 29, 31, 34 & 36, 37, 38\\
K6 对偶单纯形法 & 第 5 章 & 8, 19, 24 & 28, 32, 35 & \\
K7 右端项灵敏度 & 第 6 章 & 9, 15, 25 & 26, 29, 31 & \\
K8 价值系数灵敏度 & 第 6 章 & 10, 15, 20 & 29, 33 & \\
K9 新增变量 / 约束 & 第 6 章 & 15 & 30, 32, 35 & \\\bottomrule
\end{tabular}\par\nopagebreak[4]
\captionof{table}{知识点—题目对应表}\label{tab:cover}
\end{center}

% ============================================================
%  正式练习（38 题）：题目区只给题意图，不泄露答案
% ============================================================
\clearpage
\section*{正式练习}
\addcontentsline{toc}{section}{正式练习（38 题）}
\markboth{正式练习}{}

\begin{summarybox}[练习说明]
共 38 题：\tagBase{A. 基础巩固 25 题}\ 单选 10、多选 5、填空 5、解答 5；\tagUp{B. 中等偏下提升 10 题}\ 单选 3、多选 2、填空 3、综合大题 2；\tagCross{C. 跨学科迁移 3 题}\ 单选、多选、填空各 1。\par
请先独立完成全部题目，再查“答案速查”，最后读“详细解析”。题目区的图只画题目已经给出的信息；解析区的图另行绘制，画出解题时添加的辅助信息。
\end{summarybox}

\begin{definitionbox}[公共数据：练习模型 E（第 6、7、9、10、13、14、17、20、22、23、25、26、29--33 题使用）]
\begin{minipage}[t]{0.52\linewidth}
某车间生产两种零件，单位利润分别为 3 和 2：
\[
\begin{aligned}
\max z&=3x_1+2x_2\\
\st &x_1+x_2\le4\qquad\text{\ding1 人工（工时）}\\
&x_1+3x_2\le9\qquad\text{\ding2 材料（kg）}\\
&x_1\le3\qquad\qquad\ \ \text{\ding3 设备（台时）}\\
&x_1,x_2\ge0
\end{aligned}
\]
加松弛变量 $x_3,x_4,x_5$ 后，单纯形法得最优表（基变量次序 $x_1,x_2,x_4$）：
\begin{center}\small
\begin{tabular}{@{}ccc|ccccc@{}}\toprule
$c_B$ & $x_B$ & $b$ & $x_1$ & $x_2$ & $x_3$ & $x_4$ & $x_5$\\\midrule
3 & $x_1$ & 3 & 1 & 0 & 0 & 0 & 1\\
2 & $x_2$ & 1 & 0 & 1 & 1 & 0 & $-1$\\
0 & $x_4$ & 3 & 0 & 0 & $-3$ & 1 & 2\\\midrule
\multicolumn{3}{c|}{$\sigma_j$} & 0 & 0 & $-2$ & 0 & $-1$\\\bottomrule
\end{tabular}
\end{center}
\end{minipage}\hfill
\begin{minipage}[t]{0.44\linewidth}\vspace{0pt}\centering
\begin{tikzpicture}[x=0.85cm,y=0.85cm]
  \exRegion
  \axes{6}{4}{x_1}{x_2}
  \exLines
  \foreach \p in {(0,0),(3,0),(3,1),(1.5,2.5),(0,3)} \node[vpt] at \p {};
  \node[lab] at (1.2,1.0) {可行域};
\end{tikzpicture}
\fig{练习模型 E 的可行域与三条约束边界}\label{fig:exE}
\end{minipage}
\end{definitionbox}

% ============================================================
\subsection*{A. 基础巩固训练（25 题）}
\addcontentsline{toc}{subsection}{A. 基础巩固训练（25 题）}

\needspace{10\baselineskip}\noindent\tagBase{A1. 单项选择题（每题只有一个正确选项）}

\Q[\diff{基础·K2}] 练习模型 E 的对偶问题是（\quad）。
\begin{center}
\begin{tikzpicture}[x=0.75cm,y=0.6cm]
  \foreach \r/\a/\b/\bb in {0/1/1/4,1/1/3/9,2/1/0/3} {
    \node[draw=navy,fill=formulablue,minimum width=0.75cm,minimum height=0.6cm,font=\small] at (0,-\r) {$\a$};
    \node[draw=navy,fill=formulablue,minimum width=0.75cm,minimum height=0.6cm,font=\small] at (1,-\r) {$\b$};
    \node[draw=navy,fill=warm,minimum width=0.75cm,minimum height=0.6cm,font=\small] at (2.3,-\r) {$\bb$};
  }
  \node[draw=navy,fill=lightgreen,minimum width=0.75cm,minimum height=0.6cm,font=\small] at (0,1.2) {3};
  \node[draw=navy,fill=lightgreen,minimum width=0.75cm,minimum height=0.6cm,font=\small] at (1,1.2) {2};
  \node[font=\scriptsize] at (0.5,2.0) {$c\T$}; \node[font=\scriptsize] at (2.3,0.8) {$b$};
  \node[font=\scriptsize] at (0.5,-3.0) {$A$};
\end{tikzpicture}
\fig{第 1 题题意图：模型 E 的数据块}
\end{center}
\opt{$\min w=4y_1+9y_2+3y_3$，$\st y_1+y_2+y_3\ge3$，$y_1+3y_2\ge2$，$y\ge0$}
    {$\min w=3y_1+2y_2$，$\st y_1+y_2\ge4$，$y_1+3y_2\ge9$，$y_1\ge3$，$y\ge0$}
    {$\max w=4y_1+9y_2+3y_3$，$\st y_1+y_2+y_3\le3$，$y_1+3y_2\le2$，$y\ge0$}
    {$\min w=4y_1+9y_2+3y_3$，$\st y_1+y_2\ge3$，$y_1+3y_2+y_3\ge2$，$y\ge0$}

\Q[\diff{基础·K2}] 原问题为 $\max$ 型，它的第 $i$ 个约束是“$\ge$”型，则对偶变量 $y_i$ 应满足（\quad）。
\begin{center}
\begin{tikzpicture}[x=1cm,y=1cm]
  \node[box] (p) at (0,0) {原问题（max）\\第 $i$ 个约束：$a_i\T x\ \ge\ b_i$};
  \node[box] (d) at (5.4,0) {对偶问题（min）\\第 $i$ 个变量：$y_i\ \ ?$};
  \draw[arr] (p)--(d);
\end{tikzpicture}
\fig{第 2 题题意图}
\end{center}
\opt{$y_i\ge0$}{$y_i\le0$}{$y_i$ 无符号限制}{$y_i=0$}

\Q[\diff{基础·K3}] 某 max 型线性规划有一个可行解，目标值为 12；它的对偶问题有一个可行解，目标值为 15。关于原问题的最优值 $z^*$，正确的是（\quad）。
\begin{center}
\begin{tikzpicture}[x=0.7cm,y=1cm]
  \draw[ax] (8,0)--(19,0) node[right,font=\scriptsize]{目标值};
  \foreach \x in {10,12,15,17} {\draw (\x,0.08)--(\x,-0.08); \node[below,font=\tiny] at (\x,-0.08) {\x};}
  \fill[navy] (12,0) circle (2.5pt) node[above,font=\scriptsize,align=center]{原问题某可行解};
  \fill[greenline] (15,0) circle (2.5pt) node[above,font=\scriptsize,align=center]{对偶某可行解};
\end{tikzpicture}
\fig{第 3 题题意图：两个已知的可行目标值}
\end{center}
\opt{$z^*\le12$}{$12\le z^*\le15$}{$z^*\ge15$}{$z^*=13.5$}

\Q[\diff{基础·K3}] 线性规划 $\max z=2x_1+x_2$，$\st -x_1+x_2\le2$，$x_1,x_2\ge0$ 的可行域如图，目标函数在可行域上无上界。则它的对偶问题（\quad）。
\begin{center}
\begin{tikzpicture}[x=0.55cm,y=0.55cm]
  \fill[feasreg] (0,0)--(5,0)--(5,5)--(3,5)--(0,2)--cycle;
  \axes{5}{5}{x_1}{x_2}
  \draw[con] (0,2)--(3.3,5.3) node[right,font=\scriptsize]{$-x_1+x_2=2$};
  \draw[iso] (0,2)--(1,0); \draw[iso] (0,4)--(2,0); \draw[iso] (1.5,5)--(4,0);
  \draw[arr,shiftc] (3,2)--(4.6,2.8);
  \node[font=\scriptsize,shiftc!80!black] at (5.6,3.4) {$z$ 增大};
\end{tikzpicture}
\fig{第 4 题题意图：可行域向右无限延伸，虚线为等值线}
\end{center}
\opt{也无界}{有最优解}{没有可行解}{无法判断}

\Q[\diff{基础·K1}] 用修正单纯形法计算时，当前基 $B$ 的逆为 $B^{-1}=\begin{pmatrix}\frac12&0\\-\frac12&1\end{pmatrix}$，原模型右端项 $b=(8,7)\T$，则基变量的取值 $x_B$ 为（\quad）。
\begin{center}
\begin{tikzpicture}
  \node at (0,0) {$\underbrace{\begin{pmatrix}\frac12&0\\[2pt]-\frac12&1\end{pmatrix}}_{B^{-1}}\ \underbrace{\begin{pmatrix}8\\7\end{pmatrix}}_{b}=\ \begin{pmatrix}?\\?\end{pmatrix}$};
\end{tikzpicture}
\fig{第 5 题题意图}
\end{center}
\opt{$(4,3)\T$}{$(8,7)\T$}{$(4,-1)\T$}{$(3,4)\T$}

\Q[\diff{基础·K5}] 在练习模型 E 的最优表中，三种资源的影子价格 $(y_1,y_2,y_3)$ 为（\quad）。
\begin{center}
\small
\begin{tabular}{@{}l|ccc@{}}\toprule
最优表中的列 & $x_3$（\ding1 的松弛） & $x_4$（\ding2 的松弛） & $x_5$（\ding3 的松弛）\\\midrule
检验数 $\sigma_j$ & $-2$ & 0 & $-1$\\\bottomrule
\end{tabular}
\captionof{figure}{第 6 题题意图：最优表中松弛变量的检验数}
\end{center}
\opt{$(-2,0,-1)$}{$(2,0,1)$}{$(0,2,1)$}{$(1,0,2)$}

\Q[\diff{基础·K4}] 练习模型 E 的最优解中，资源 \ding2（材料）还剩 3 kg 未用。则资源 \ding2 的影子价格 $y_2$（\quad）。
\begin{center}
\begin{tikzpicture}[x=1cm,y=1cm]
  \fill[auxblue!60] (0,0) rectangle (4,0.5); \draw[navy] (0,0) rectangle (6,0.5);
  \node[font=\scriptsize] at (2,0.25) {已用 6};
  \node[font=\scriptsize] at (5,0.25) {剩余 3};
  \node[font=\scriptsize,left] at (0,0.25) {材料 9 kg};
\end{tikzpicture}
\fig{第 7 题题意图：资源 \ding2 的使用情况}
\end{center}
\opt{必为正数}{必为 0}{可能为负数}{等于 3}

\Q[\diff{基础·K6}] 用对偶单纯形法迭代时，每一轮\textbf{首先}确定的是（\quad）。
\opt{入基变量}{出基变量}{检验数最大的列}{$b$ 列除以入基列的最小比值所在行}

\Q[\diff{基础·K7}] 练习模型 E 中，资源 \ding1 的影子价格 $y_1=2$，其允许变动范围为 $b_1\in[3,5]$。若 $b_1$ 由 4 增至 4.5，则新的最优值为（\quad）。
\begin{center}
\begin{tikzpicture}[x=0.75cm,y=0.75cm]
  \exRegion
  \axes{6}{4}{x_1}{x_2}
  \exLines
  \draw[shiftl] (0.3,4.2)--(5.0,-0.5) node[below,font=\scriptsize]{$x_1+x_2=4.5$};
\end{tikzpicture}
\fig{第 9 题题意图：橙色虚线为资源 \ding1 增加后的新边界}
\end{center}
\opt{11}{12}{11.5}{13}

\Q[\diff{基础·K8}] 练习模型 E 中，$c_2$ 的允许变动范围为 $[0,3]$。若 $c_2$ 由 2 改为 4，则（\quad）。
\opt{最优解不变}{最优解改变，应从原最优表出发继续用单纯形法迭代}{应从原最优表出发用对偶单纯形法迭代}{问题变为无界}

\medskip
\needspace{10\baselineskip}\noindent\tagBase{A2. 多项选择题（每题至少两个正确选项）}

\Q[\diff{基础·K2}] 关于对偶问题的构造（原问题为 max 型），下列说法正确的有（\quad）。
\begin{center}
\begin{tikzpicture}[x=1cm,y=1cm]
  \node[box,minimum width=3cm] (p) at (0,0) {原问题\\$m$ 个约束，$n$ 个变量};
  \node[box,minimum width=3cm] (d) at (5,0) {对偶问题\\$?$ 个约束，$?$ 个变量};
  \draw[arr] (p)--node[above,font=\scriptsize]{构造}(d);
\end{tikzpicture}
\fig{第 11 题题意图}
\end{center}
\begin{enumerate}[label=\Alph*.,leftmargin=2em]
  \item 原问题的第 $i$ 个约束对应对偶问题的第 $i$ 个变量；
  \item 原问题的变量 $x_j\ge0$ 对应对偶问题的第 $j$ 个约束为“$\ge$”型；
  \item 原问题的右端项 $b$ 成为对偶问题的目标函数系数；
  \item 对偶问题的对偶不一定是原问题。
\end{enumerate}

\Q[\diff{基础·K3}] 设 $\bar x$ 是 max 型原问题的可行解，$\bar y$ 是其对偶问题的可行解。下列说法正确的有（\quad）。
\begin{enumerate}[label=\Alph*.,leftmargin=2em]
  \item $c\T\bar x\le b\T\bar y$；
  \item 若 $c\T\bar x=b\T\bar y$，则 $\bar x$、$\bar y$ 分别是两个问题的最优解；
  \item 若原问题有最优解，则对偶问题也有最优解，且两者最优值相等；
  \item 若对偶问题没有可行解，则原问题一定无界。
\end{enumerate}

\Q[\diff{基础·K5}] 关于影子价格，下列说法正确的有（\quad）。
\begin{enumerate}[label=\Alph*.,leftmargin=2em]
  \item 在最优基不变的范围内，$y_i$ 表示资源 $i$ 每增加 1 单位时最优值的增量；
  \item 影子价格就是该资源的市场价格；
  \item 若资源的市场价格低于其影子价格，买入该资源（在允许范围内）可以增加利润；
  \item 影子价格对任意大小的资源变化都保持不变。
\end{enumerate}

\Q[\diff{基础·K4}] 练习模型 E 的最优解为 $x^*=(3,1)$。根据互补松弛条件，下列说法正确的有（\quad）。
\begin{center}
\begin{tikzpicture}[x=0.75cm,y=0.75cm]
  \exRegion
  \axes{6}{4}{x_1}{x_2}
  \exLines
  \node[optpt] at (3,1) {};
  \node[lab,right] at (3.15,1.25) {$x^*=(3,1)$};
\end{tikzpicture}
\fig{第 14 题题意图：已知最优点}
\end{center}
\begin{enumerate}[label=\Alph*.,leftmargin=2em]
  \item $y_2=0$；
  \item 因为 $x_1^*>0$，对偶问题的第一个约束 $y_1+y_2+y_3\ge3$ 取等号；
  \item 若 $y_1>0$，则约束 $x_1+x_2\le4$ 在最优点取等号；
  \item 一般地，若某资源的影子价格为 0，则该资源在最优解中一定有剩余。
\end{enumerate}

\Q[\diff{基础·K7/K8/K9}] 关于灵敏度分析，下列说法正确的有（\quad）。
\begin{enumerate}[label=\Alph*.,leftmargin=2em]
  \item 只改变右端项 $b$ 时，最优表的检验数行不变；
  \item 只改变非基变量 $x_k$ 的价值系数时，只有 $\sigma_k$ 会改变；
  \item 改变基变量的价值系数 $c_j$（在允许范围内）时，基变量的取值 $x_B$ 会随之改变；
  \item 新增一条约束时，若原最优解满足这条约束，则原最优解仍是最优解。
\end{enumerate}

\medskip
\needspace{10\baselineskip}\noindent\tagBase{A3. 填空题}

\Q[\diff{基础·K2}] 线性规划 $\min z=x_1+2x_2$，$\st x_1+x_2\ge2$，$x_1+3x_2\ge3$，$x_1,x_2\ge0$ 的对偶问题为：$\max w=$\underline{\hspace{2.5cm}}，$\st$ \underline{\hspace{2.5cm}}，\underline{\hspace{2.5cm}}，$y_1,y_2\ge0$。
\begin{center}
\begin{tikzpicture}[x=0.6cm,y=0.6cm]
  \fill[feasreg] (0,4)--(0,2)--(1.5,0.5)--(3,0)--(5,0)--(5,4)--cycle;
  \axes{5}{4}{x_1}{x_2}
  \draw[con] (0,2)--(2.3,-0.3) node[below,font=\scriptsize]{$x_1+x_2=2$};
  \draw[con] (0,1)--(5,-0.667);
  \node[font=\scriptsize,above] at (4.6,-0.2) {$x_1+3x_2=3$};
\end{tikzpicture}
\fig{第 16 题题意图：原问题的可行域（向右上方无界）}
\end{center}

\Q[\diff{基础·K5}] 练习模型 E 的影子价格为 $y=(2,0,1)$，右端项 $b=(4,9,3)\T$。利用 $y\T b$ 可得最优值 $z^*=$\underline{\hspace{2cm}}。
\begin{center}
\begin{tikzpicture}[x=1cm,y=1cm]
  \foreach \i/\bb/\yy in {0/4/2,1/9/0,2/3/1} {
    \node[draw=navy,fill=warm,minimum width=1.2cm,font=\small] at (\i*1.5,0) {$b_{\the\numexpr\i+1\relax}=\bb$};
    \node[draw=navy,fill=lightblue,minimum width=1.2cm,font=\small] at (\i*1.5,-0.7) {$y_{\the\numexpr\i+1\relax}=\yy$};
  }
\end{tikzpicture}
\fig{第 17 题题意图}
\end{center}

\Q[\diff{基础·K1}] 基矩阵 $B=\begin{pmatrix}2&1\\1&3\end{pmatrix}$ 的逆矩阵 $B^{-1}=$\underline{\hspace{3cm}}。

\Q[\diff{基础·K6}] 用对偶单纯形法求解 $\min z=x_1+2x_2$（约束同第 16 题），化为 max 型并乘 $-1$ 后的起始表如下。本轮出基变量为\underline{\hspace{1.2cm}}，入基变量为\underline{\hspace{1.2cm}}，$\theta=$\underline{\hspace{1.2cm}}。
\begin{center}
\small
\begin{tabular}{@{}ccc|cccc@{}}\toprule
$c_B$ & $x_B$ & $b$ & $x_1$ & $x_2$ & $x_3$ & $x_4$\\\midrule
0 & $x_3$ & $-2$ & $-1$ & $-1$ & 1 & 0\\
0 & $x_4$ & $-3$ & $-1$ & $-3$ & 0 & 1\\\midrule
\multicolumn{3}{c|}{$\sigma_j$} & $-1$ & $-2$ & 0 & 0\\\bottomrule
\end{tabular}
\captionof{figure}{第 19 题题意图：起始表}
\end{center}

\Q[\diff{基础·K8}] 练习模型 E 的最优表中，非基变量 $x_3$ 的检验数 $\sigma_3=-2$。若把 $x_3$ 的价值系数由 0 改为 $\Delta$，要保持原最优解仍最优，$\Delta$ 应满足 $\Delta\le$\underline{\hspace{1.5cm}}。

\medskip
\needspace{10\baselineskip}\noindent\tagBase{A4. 解答题}

\Q[\diff{基础·K2}] 写出下列线性规划的对偶问题：
\[
\max z=x_1+2x_2-x_3,\qquad
\begin{cases}x_1+x_2+x_3\le6\\ 2x_1-x_2\ge2\\ x_2+x_3=4\\ x_1\ge0,\ x_2\le0,\ x_3\text{ 无约束}\end{cases}
\]
\begin{center}
\begin{tikzpicture}[x=1cm,y=1cm,c/.style={draw=navy,fill=lightblue,minimum width=1.7cm,minimum height=0.55cm,font=\small}]
  \node[font=\scriptsize\bfseries] at (0,0.7) {约束类型};
  \node[c] at (0,0) {$\le6$}; \node[c] at (0,-0.65) {$\ge2$}; \node[c] at (0,-1.3) {$=4$};
  \node[font=\scriptsize\bfseries] at (3.2,0.7) {变量类型};
  \node[c,fill=warm] at (3.2,0) {$x_1\ge0$}; \node[c,fill=warm] at (3.2,-0.65) {$x_2\le0$}; \node[c,fill=warm] at (3.2,-1.3) {$x_3$ 自由};
\end{tikzpicture}
\fig{第 21 题题意图：原问题的约束与变量类型}
\end{center}

\Q[\diff{基础·K1}] 对练习模型 E，以松弛变量 $\{x_3,x_4,x_5\}$ 为初始基（$B=I$）做一轮修正单纯形法迭代：
(1) 计算初始检验数，确定入基变量与出基变量；
(2) 写出新基矩阵 $B$ 及其逆 $B^{-1}$，求新的 $x_B$；
(3) 计算新的单纯形乘子 $y\T=c_B\T B^{-1}$ 和非基变量检验数，判断是否已最优。
\begin{center}
\begin{tikzpicture}[x=0.7cm,y=0.7cm]
  \exRegion
  \axes{6}{4}{x_1}{x_2}
  \exLines
  \node[badpt,fill=navy] at (0,0) {};
  \node[lab,above right] at (0.1,0.1) {初始基解 $(0,0)$};
\end{tikzpicture}
\fig{第 22 题题意图：初始基解位于原点}
\end{center}

\Q[\diff{基础·K5}] 练习模型 E 中，资源 \ding3（设备）的影子价格 $y_3=1$，允许范围 $b_3\in[1.5,4]$。现在可以从外部租用设备，每台时租金 0.8。问：是否值得租？最多租几台时（使影子价格仍然有效）？净增利润是多少？租用后的最优解是什么？

\Q[\diff{基础·K6}] 用对偶单纯形法求解 $\min z=x_1+2x_2$，$\st x_1+x_2\ge2$，$x_1+3x_2\ge3$，$x_1,x_2\ge0$（起始表见第 19 题）。

\Q[\diff{基础·K7}] 求练习模型 E 中资源 \ding3 的右端项 $b_3$ 的允许变动范围，并用图说明两个端点的含义。

% ============================================================
\subsection*{B. 中等偏下提升训练（10 题）}
\addcontentsline{toc}{subsection}{B. 中等偏下提升训练（10 题）}

\needspace{10\baselineskip}\noindent\tagUp{B1. 单项选择题}

\Q[\diff{提升·K5+K7}] 练习模型 E 中 $y_1=2$，$b_1$ 的允许范围为 $[3,5]$。若资源 \ding1 由 4 增至 6，最优值实际增加（\quad）。
\begin{center}
\begin{tikzpicture}[x=0.75cm,y=0.75cm]
  \exRegion
  \axes{6}{4}{x_1}{x_2}
  \exLines
  \draw[shiftl] (1.8,4.2)--(6.3,-0.3) node[right,font=\scriptsize]{$x_1+x_2=6$};
\end{tikzpicture}
\fig{第 26 题题意图：资源 \ding1 增至 6 后的新边界（橙色虚线）}
\end{center}
\opt{4}{2}{0}{6}

\Q[\diff{提升·K4+K3}] 线性规划 $\max z=4x_1+3x_2$，$\st x_1+x_2\le5$，$2x_1+x_2\le8$，$x_1,x_2\ge0$ 的最优解为 $x^*=(3,2)$。其对偶问题的最优解 $y^*$ 为（\quad）。
\begin{center}
\begin{tikzpicture}[x=0.6cm,y=0.6cm]
  \fill[feasreg] (0,0)--(4,0)--(3,2)--(0,5)--cycle;
  \axes{6}{5}{x_1}{x_2}
  \draw[con] (-0.2,5.2)--(5.4,-0.4);
  \draw[con] (4.2,-0.4)--(1.4,5.2);
  \node[optpt] at (3,2) {}; \node[lab,right] at (3.15,2.25) {$x^*=(3,2)$};
\end{tikzpicture}
\fig{第 27 题题意图：原问题的可行域与已知最优点}
\end{center}
\opt{$(2,1)$}{$(1,2)$}{$(0,4)$}{$(3,0)$}

\Q[\diff{提升·K6}] 对偶单纯形法某一轮中，出基行与检验数行如下（右端 $b_r=-6$）。应选（\quad）入基。
\begin{center}
\small
\begin{tabular}{@{}l|ccccc@{}}\toprule
 & $x_1$ & $x_2$ & $x_3$ & $x_4$ & $x_5$\\\midrule
出基行 $a_r$ & $-2$ & 1 & $-4$ & 0 & 1\\
检验数 $\sigma$ & $-3$ & $-1$ & $-2$ & 0 & 0\\\bottomrule
\end{tabular}
\captionof{figure}{第 28 题题意图}
\end{center}
\opt{$x_1$}{$x_2$}{$x_3$}{问题无可行解}

\medskip
\needspace{10\baselineskip}\noindent\tagUp{B2. 多项选择题}

\Q[\diff{提升·K5+K7+K8}] 对练习模型 E（$y=(2,0,1)$；$b_1\in[3,5]$，$b_2\in[6,+\infty)$，$b_3\in[1.5,4]$；$c_1\in[2,+\infty)$，$c_2\in[0,3]$），下列说法正确的有（\quad）。
\begin{enumerate}[label=\Alph*.,leftmargin=2em]
  \item $b_2$ 由 9 降为 7，最优解仍为 $(3,1)$，最优值不变；
  \item $c_1$ 由 3 降为 2.5，最优解仍为 $(3,1)$，最优值变为 9.5；
  \item $b_3$ 由 3 增至 5，最优值增加 $2\times1=2$；
  \item $c_2$ 由 2 增至 2.5，最优解不变。
\end{enumerate}

\Q[\diff{提升·K9+K5}] 在练习模型 E 中新增一种零件 $x_6$，每件消耗资源 \ding1、\ding2、\ding3 分别为 1、2、0，单位利润为 $c_6$。下列说法正确的有（\quad）。
\begin{center}
\begin{tikzpicture}[x=1cm,y=1cm]
  \node[draw=navy,fill=lightblue,font=\small] at (0,0) {$P_6=\begin{pmatrix}1\\2\\0\end{pmatrix}$};
  \node[draw=navy,fill=lightblue,font=\small] at (3.2,0) {$y=(2,0,1)$};
  \node[draw=navy,fill=warm,font=\small] at (6.0,0) {利润 $c_6$};
\end{tikzpicture}
\fig{第 30 题题意图：新产品的数据}
\end{center}
\begin{enumerate}[label=\Alph*.,leftmargin=2em]
  \item 若 $c_6=1.5$，新零件不值得生产；
  \item 若 $c_6=3$，新零件值得生产，应把它的列写入最优表后继续单纯形法；
  \item 判断新零件是否值得生产，必须重新求解整个问题；
  \item 新增变量后，原最优解（令 $x_6=0$）仍然满足原始可行性。
\end{enumerate}

\medskip
\needspace{10\baselineskip}\noindent\tagUp{B3. 填空题}

\Q[\diff{提升·K5+K7}] 练习模型 E 中，资源 \ding3 的影子价格为\underline{\hspace{1.2cm}}；在最优基不变的前提下，$b_3$ 最多可以增至\underline{\hspace{1.2cm}}，此时最优值为\underline{\hspace{1.2cm}}。

\Q[\diff{提升·K6+K9}] 在练习模型 E 中新增约束 $x_2\le0.5$。用对偶单纯形法求得的新最优解为 $x^*=$\underline{\hspace{2cm}}，新最优值为\underline{\hspace{1.5cm}}。
\begin{center}
\begin{tikzpicture}[x=0.75cm,y=0.75cm]
  \exRegion
  \axes{6}{4}{x_1}{x_2}
  \exLines
  \draw[shiftl] (-0.2,0.5)--(6.3,0.5) node[right,font=\scriptsize]{$x_2=0.5$};
\end{tikzpicture}
\fig{第 32 题题意图：新增约束的边界（橙色虚线）}
\end{center}

\Q[\diff{提升·K8}] 练习模型 E 中，基变量 $x_1$ 的价值系数 $c_1$ 的允许变动范围为\underline{\hspace{2.5cm}}。

\medskip
\needspace{10\baselineskip}\noindent\tagUp{B4. 综合大题}

\Q[\diff{提升·K2+K4+K5+K7}] 一家小面包坊生产面包（$x_1$ 批）和蛋糕（$x_2$ 批），每批利润分别为 5 和 4（百元）。面粉每天 8 袋，烤箱每天 12 小时，奶油每天 3.5 盒：
\[
\max z=5x_1+4x_2,\qquad
\begin{cases}x_1+2x_2\le8&(\text{面粉})\\ 3x_1+2x_2\le12&(\text{烤箱})\\ x_2\le3.5&(\text{奶油})\\ x_1,x_2\ge0\end{cases}
\]
已知最优方案为 $x^*=(2,3)$，最优表中基变量次序为 $x_1,x_2,x_5$，
$B^{-1}=\begin{pmatrix}-\frac12&\frac12&0\\[2pt]\frac34&-\frac14&0\\[2pt]-\frac34&\frac14&1\end{pmatrix}$（列依次对应面粉、烤箱、奶油）。
\begin{enumerate}[label=(\arabic*)]
  \item 写出对偶问题；
  \item 用互补松弛条件求三种资源的影子价格；
  \item 若可以额外租用烤箱，每小时租金 1（百元），是否值得？租 1 小时净增利润多少？（需说明为何可以直接使用影子价格）
  \item 求面粉可用量 $b_1$ 的允许变动范围。
\end{enumerate}
\begin{center}
\begin{tikzpicture}[x=0.7cm,y=0.7cm]
  \fill[feasreg] (0,0)--(4,0)--(2,3)--(1,3.5)--(0,3.5)--cycle;
  \axes{6}{5}{x_1}{x_2}
  \draw[con] (-0.4,4.2)--(6.3,0.85) node[right,font=\scriptsize]{面粉};
  \draw[con] (4.13,-0.2)--(0.8,4.8) node[above,font=\scriptsize]{烤箱};
  \draw[con] (-0.2,3.5)--(6.3,3.5) node[right,font=\scriptsize]{奶油};
  \node[optpt] at (2,3) {}; \node[lab,right] at (2.15,2.7) {$x^*=(2,3)$};
\end{tikzpicture}
\fig{第 34 题题意图：面包坊问题的可行域与已知最优点}
\end{center}

\Q[\diff{提升·K6+K9+K1}] （续第 34 题）面包坊增加一条包装限制：每天最多包装 4.5 批，即 $x_1+x_2\le4.5$。最优表中 $x_1,x_2$ 两行为
$x_1-\frac12x_3+\frac12x_4=2$，\ $x_2+\frac34x_3-\frac14x_4=3$，检验数 $\sigma_3=-\frac12$，$\sigma_4=-\frac32$（$x_3,x_4,x_5$ 分别为面粉、烤箱、奶油的松弛变量）。
\begin{enumerate}[label=(\arabic*)]
  \item 原最优方案是否满足新限制？
  \item 把新约束加松弛变量 $x_6$ 后写成最优表中的一行（消去基变量的系数）；
  \item 用对偶单纯形法求新的最优方案；
  \item 新的最优利润是多少？比原来减少多少？用新约束的影子价格解释这一减少量。
\end{enumerate}
\begin{center}
\begin{tikzpicture}[x=0.7cm,y=0.7cm]
  \fill[feasreg] (0,0)--(4,0)--(2,3)--(1,3.5)--(0,3.5)--cycle;
  \axes{6}{5}{x_1}{x_2}
  \draw[con] (-0.4,4.2)--(6.3,0.85);
  \draw[con] (4.13,-0.2)--(0.8,4.8);
  \draw[con] (-0.2,3.5)--(6.3,3.5);
  \draw[shiftl] (0,4.5)--(4.8,-0.3) node[below,font=\scriptsize]{$x_1+x_2=4.5$};
  \node[optpt] at (2,3) {};
\end{tikzpicture}
\fig{第 35 题题意图：新增包装限制（橙色虚线）}
\end{center}

% ============================================================
\subsection*{C. 跨学科迁移训练（3 题）}
\addcontentsline{toc}{subsection}{C. 跨学科迁移训练（3 题）}

\needspace{10\baselineskip}\noindent\tagCross{C1. 跨学科单项选择题}

\Q[\diff{跨学科·电力系统·K5}] \textbf{背景：}电网调度要用最低成本满足负荷。设两台发电机：G1 发电成本 20 元/MWh，最大出力 60 MW；G2 发电成本 30 元/MWh，最大出力 100 MW。负荷为 80 MW。调度模型为
$\min C=20p_1+30p_2$，$\st p_1+p_2=80$（供需平衡），$0\le p_1\le60$，$0\le p_2\le100$。
供需平衡约束的对偶变量（影子价格）称为“边际电价”，表示负荷再增加 1 MW 时最低总成本的增加量。该边际电价为（\quad）元/MWh。
\begin{center}
\begin{tikzpicture}[x=0.06cm,y=0.085cm]
  \draw[ax] (0,0)--(175,0) node[right,font=\scriptsize]{累计出力 / MW};
  \draw[ax] (0,0)--(0,45) node[above,font=\scriptsize]{成本 元/MWh};
  \fill[auxblue!40] (0,0) rectangle (60,20); \node[font=\scriptsize] at (30,10) {G1：20};
  \fill[auxblue!70] (60,0) rectangle (160,30); \node[font=\scriptsize] at (110,15) {G2：30};
  \foreach \x in {60,160} \node[below,font=\tiny] at (\x,0) {\x};
  \foreach \y in {20,30} \node[left,font=\tiny] at (0,\y) {\y};
  \draw[redline,thick,dashed] (80,0)--(80,40) node[above,font=\scriptsize]{负荷 80};
\end{tikzpicture}
\fig{第 36 题题意图：两台机组按成本由低到高排列}
\end{center}
\opt{20}{25}{30}{50}

\medskip
\needspace{10\baselineskip}\noindent\tagCross{C2. 跨学科多项选择题}

\Q[\diff{跨学科·营养学·K3+K4+K5}] \textbf{背景：}营养配餐要用最少的钱满足营养需求。食物 $a$ 每份 3 元，含蛋白质 2 单位、维生素 1 单位；食物 $b$ 每份 2 元，含蛋白质 1 单位、维生素 2 单位。每天至少需要蛋白质 8 单位、维生素 7 单位：
$\min C=3a+2b$，$\st 2a+b\ge8$，$a+2b\ge7$，$a,b\ge0$。
其对偶问题可理解为：一家营养片厂商给每单位蛋白质、维生素分别定价 $u,v$，在“营养片不比食物贵”（$2u+v\le3$，$u+2v\le2$）的前提下使收入 $8u+7v$ 最大。已知最优配餐为 $a=3$，$b=2$。下列说法正确的有（\quad）。
\begin{center}
\begin{tikzpicture}[x=0.42cm,y=0.42cm]
  \fill[feasreg] (0,9)--(0,8)--(3,2)--(7,0)--(9,0)--(9,9)--cycle;
  \axes{9}{9}{a}{b}
  \draw[con] (0,8)--(4.4,-0.8) node[below,font=\scriptsize]{蛋白质};
  \draw[con] (-0.6,3.8)--(9.4,-1.2) node[below,font=\scriptsize]{维生素};
  \node[optpt] at (3,2) {}; \node[lab,above right] at (3.1,2.1) {$(3,2)$};
\end{tikzpicture}
\fig{第 37 题题意图：配餐问题的可行域与已知最优点}
\end{center}
\begin{enumerate}[label=\Alph*.,leftmargin=2em]
  \item 每天最低花费为 13 元；
  \item 若蛋白质需求增加 1 单位（9 单位），最低花费增加 $\frac43$ 元；
  \item 对偶最优解 $(u^*,v^*)$ 可理解为营养片厂商对每单位营养素的最高合理定价，且 $8u^*+7v^*=13$；
  \item 因为 $a=3>0$、$b=2>0$，两种营养约束的影子价格都必为 0。
\end{enumerate}

\medskip
\needspace{10\baselineskip}\noindent\tagCross{C3. 跨学科填空题}

\Q[\diff{跨学科·计算机·K3+K5}] \textbf{背景：}云平台把 CPU 和内存分配给多类计算任务，以最大化收益为目标建立线性规划（约束为“资源用量 $\le$ 容量”）。求解后得到：CPU 容量 40 核时，影子价格 0.5 元/核时；内存容量 100 GB$\cdot$时，影子价格 0.2 元/(GB$\cdot$时)。则该平台的最大收益为\underline{\hspace{1.5cm}}元；若再以 0.3 元/核时租入 5 核时 CPU（已知在允许范围内），净收益增加\underline{\hspace{1.5cm}}元。
\begin{center}
\begin{tikzpicture}[x=1cm,y=1cm]
  \node[draw=navy,fill=warm,font=\small,minimum width=3cm] at (0,0) {CPU：40 核时};
  \node[draw=navy,fill=warm,font=\small,minimum width=3cm] at (0,-0.7) {内存：100 GB$\cdot$时};
  \node[draw=navy,fill=lightblue,font=\small,minimum width=2.6cm] at (3.6,0) {$y_{\text{CPU}}=0.5$};
  \node[draw=navy,fill=lightblue,font=\small,minimum width=2.6cm] at (3.6,-0.7) {$y_{\text{内存}}=0.2$};
\end{tikzpicture}
\fig{第 38 题题意图}
\end{center}

% ============================================================
%  答案速查（只查结果，不讲过程）
% ============================================================
\clearpage
\section*{全部答案速查}
\addcontentsline{toc}{section}{全部答案速查}
\markboth{答案速查}{}

\begin{center}
\small
\renewcommand{\arraystretch}{1.25}
\begin{tabularx}{\linewidth}{@{}>{\bfseries}p{2.6cm}X@{}}
\toprule
\rowcolor{tagbase}基础单选（1--10） & 1.~A\quad 2.~B\quad 3.~B\quad 4.~C\quad 5.~A\quad 6.~B\quad 7.~B\quad 8.~B\quad 9.~B\quad 10.~B\\
\rowcolor{tagbase}基础多选（11--15） & 11.~ABC\quad 12.~ABC\quad 13.~AC\quad 14.~ABC\quad 15.~ABD\\
\rowcolor{tagbase}基础填空（16--20） & 16.~$2y_1+3y_2$；$y_1+y_2\le1$；$y_1+3y_2\le2$\quad 17.~11\quad 18.~$\frac15\begin{pmatrix}3&-1\\-1&2\end{pmatrix}$\quad 19.~$x_4$；$x_2$；$\frac23$\quad 20.~2\\
\rowcolor{tagbase}基础解答（21--25） & 21.~$\min w=6y_1+2y_2+4y_3$，$y_1+2y_2\ge1$，$y_1-y_2+y_3\le2$，$y_1+y_3=-1$，$y_1\ge0,y_2\le0,y_3$ 自由\newline
22.~$x_1$ 入基、$x_5$ 出基；$B^{-1}=\begin{pmatrix}1&0&-1\\0&1&-1\\0&0&1\end{pmatrix}$，$x_B=(1,6,3)\T$；$y=(0,0,3)$，$\sigma_2=2>0$，未最优\newline
23.~值得租；最多租 1 台时；净增 0.2；新最优解 $(4,0)$，$z=12$\newline
24.~$x^*=(\frac32,\frac12)$，$z^*=\frac52$\newline
25.~$b_3\in[1.5,4]$\\
\rowcolor{tagup}提升单选（26--28） & 26.~B\quad 27.~A\quad 28.~C\\
\rowcolor{tagup}提升多选（29--30） & 29.~ABD\quad 30.~ABD\\
\rowcolor{tagup}提升填空（31--33） & 31.~1；4；12\quad 32.~$(3,0.5)$；10\quad 33.~$[2,+\infty)$\\
\rowcolor{tagup}提升大题（34--35） & 34.~(1) $\min w=8y_1+12y_2+3.5y_3$，$y_1+3y_2\ge5$，$2y_1+2y_2+y_3\ge4$，$y\ge0$；(2) $y=(0.5,1.5,0)$；(3) 值得，净增 0.5；(4) $b_1\in[4,\frac{26}{3}]$\newline
35.~(1) 不满足；(2) $x_6-\frac14x_3-\frac14x_4=-\frac12$；(3) $x^*=(3,1.5)$；(4) $z^*=21$，减少 1，$=y_6\times0.5=2\times0.5$\\
\rowcolor{tagcross}跨学科（36--38） & 36.~C\quad 37.~ABC\quad 38.~40 元；1 元\\
\bottomrule
\end{tabularx}
\end{center}

% ============================================================
%  基础巩固 25 题详细解析
% ============================================================
\section*{基础巩固训练详细解析}
\addcontentsline{toc}{section}{详细解析：基础巩固（1--25）}
\markboth{详细解析：基础巩固}{}

\subsection*{基础单选题解析}

% ---------------- 1 ----------------
\A{1}
\SolHead{A}{基础}{K2 对称形式的对偶}
\Lb{对应公式}$\max c\T x,\ Ax\le b,\ x\ge0\ \Longleftrightarrow\ \min b\T y,\ A\T y\ge c,\ y\ge0$。
\begin{figure}[!htbp]\centering
\begin{tikzpicture}[x=0.75cm,y=0.6cm]
  \foreach \r/\a/\b/\bb in {0/1/1/4,1/1/3/9,2/1/0/3} {
    \node[draw=navy,fill=formulablue,minimum width=0.75cm,minimum height=0.6cm,font=\small] at (0,-\r) {$\a$};
    \node[draw=navy,fill=formulablue,minimum width=0.75cm,minimum height=0.6cm,font=\small] at (1,-\r) {$\b$};
    \node[draw=navy,fill=warm,minimum width=0.75cm,minimum height=0.6cm,font=\small] at (2.3,-\r) {$\bb$};
  }
  \foreach \c/\v in {0/3,1/2} \node[draw=navy,fill=lightgreen,minimum width=0.75cm,minimum height=0.6cm,font=\small] at (\c,1.2) {$\v$};
  \draw[-{Stealth[length=3mm]},very thick,shiftc] (3.3,-0.5)--(5.2,-0.5) node[midway,above,font=\scriptsize]{转置};
  \begin{scope}[xshift=4.8cm]
    \foreach \r/\a/\b/\cc/\cv in {0/1/1/1/3,1/1/3/0/2} {
      \node[draw=navy,fill=formulablue,minimum width=0.75cm,minimum height=0.6cm,font=\small] at (0,-\r) {$\a$};
      \node[draw=navy,fill=formulablue,minimum width=0.75cm,minimum height=0.6cm,font=\small] at (1,-\r) {$\b$};
      \node[draw=navy,fill=formulablue,minimum width=0.75cm,minimum height=0.6cm,font=\small] at (2,-\r) {$\cc$};
      \node[draw=navy,fill=lightgreen,minimum width=0.75cm,minimum height=0.6cm,font=\small] at (3.3,-\r) {$\cv$};
    }
    \foreach \c/\v in {0/4,1/9,2/3} \node[draw=navy,fill=warm,minimum width=0.75cm,minimum height=0.6cm,font=\small] at (\c,1.2) {$\v$};
    \node[font=\scriptsize] at (1,2.0) {对偶目标系数 $b\T$};
    \node[font=\scriptsize,align=center] at (3.3,0.8) {右端 $c$};
    \node[font=\scriptsize] at (1,-2.0) {$A\T$：第 1 行 $(1,1,1)$，第 2 行 $(1,3,0)$};
  \end{scope}
\end{tikzpicture}
\fig{第 1 题解析图：数据块整体转置后，按行读出两条对偶约束}
\end{figure}
\Lb{解题思路}模型 E 是标准的对称形式（max、$\le$、$x\ge0$），直接转置。\par
\Lb{判断过程}由图：对偶目标系数是原右端项 $(4,9,3)$；对偶约束的系数是 $A$ 的列：$x_1$ 列 $(1,1,1)$、$x_2$ 列 $(1,3,0)$；右端是原目标系数 3、2；方向为 $\ge$，$y\ge0$。所以对偶为 $\min 4y_1+9y_2+3y_3$，$y_1+y_2+y_3\ge3$，$y_1+3y_2\ge2$。\par
\Lb{各选项}A 正确。B 把 $b$ 与 $c$ 放反了：目标用了 $c$，右端用了 $b$。C 把 min 写成 max、$\ge$ 写成 $\le$。D 没有转置：用 $A$ 的行 $(1,1)$、$(1,3)$ 作系数，且把 $y_3$ 错放到第二个约束。\par
\Lb{常见错误}只换了目标方向而忘记转置。\quad\Lb{一句话总结}对称形式的对偶 $=$ 数据块整体转置 $+$ max 变 min $+$ $\le$ 变 $\ge$。

% ---------------- 2 ----------------
\A{2}
\SolHead{B}{基础}{K2 一般形式的对应规则}
\Lb{对应规则}max 问题中，“$\le$”约束是正常的，对应 $y_i\ge0$；“$\ge$”是反常的，对应 $y_i\le0$；“$=$”对应 $y_i$ 无约束。
\begin{figure}[!htbp]\centering
\begin{tikzpicture}[x=1cm,y=1cm]
  \foreach \i/\L/\t/\c in {0/A/{$y_i\ge0$}/0,1/B/{$y_i\le0$}/1,2/C/{无约束}/2,3/D/{$y_i=0$}/3} {
    \begin{scope}[xshift=\i*3.9cm]
      \draw[ax] (-1.5,0)--(1.6,0) node[right,font=\tiny]{$y_i$};
      \draw (0,0.08)--(0,-0.08) node[below,font=\tiny]{0};
      \ifnum\c=0 \draw[very thick,badc,-{Stealth}] (0,0.25)--(1.5,0.25); \fill[badc] (0,0.25) circle (2pt);\fi
      \ifnum\c=1 \draw[very thick,greenline,-{Stealth}] (0,0.25)--(-1.4,0.25); \fill[greenline] (0,0.25) circle (2pt);\fi
      \ifnum\c=2 \draw[very thick,badc,{Stealth}-{Stealth}] (-1.4,0.25)--(1.5,0.25);\fi
      \ifnum\c=3 \fill[badc] (0,0.25) circle (3pt);\fi
      \node[font=\scriptsize] at (0,0.75) {\textbf{\L}.\ \t};
      \ifnum\c=1 \node[font=\scriptsize,greenline] at (0,-0.65) {✓ 反常约束 $\to$ 反常变量};\else \node[font=\scriptsize,badc] at (0,-0.65) {✗};\fi
    \end{scope}
  }
\end{tikzpicture}
\fig{第 2 题选项图：四个选项分别对应数轴上的四种取值范围}
\end{figure}
\Lb{判断过程}把“$\ge$”约束两边乘 $-1$ 化成正常的“$\le$”，它对应 $y_i'\ge0$；换回原约束时令 $y_i=-y_i'$，故 $y_i\le0$（第 2 章推导第 1 条）。图中只有 B 的射线指向负半轴。\par
\Lb{各选项}A 是“$\le$”约束对应的情形；C 是等式约束对应的情形；D 没有对应的约束类型——对偶变量从不被强制为 0。\par
\Lb{常见错误}把 min 问题的规则（$\ge$ 对应 $y_i\ge0$）套用到 max 问题上。\quad\Lb{一句话总结}正常对正常，反常对反常。

% ---------------- 3 ----------------
\A{3}
\SolHead{B}{基础}{K3 弱对偶性}
\Lb{对应定理}弱对偶：原可行值 $\le$ 对偶可行值；故任一原可行值是 $z^*$ 的下界，任一对偶可行值是 $z^*$ 的上界。
\begin{figure}[!htbp]\centering
\begin{tikzpicture}[x=0.32cm,y=1cm]
  \foreach \i/\L/\a/\b/\ok in {0/A/8/12/0,1/B/12/15/1,2/C/15/19/0,3/D/13.5/13.5/0} {
    \begin{scope}[xshift=\i*4.1cm]
      \draw[ax] (8,0)--(19.5,0);
      \foreach \x in {12,15} {\draw (\x,0.08)--(\x,-0.08); \node[below,font=\tiny] at (\x,-0.08) {\x};}
      \ifnum\ok=1 \def\cc{greenline}\else\def\cc{badc}\fi
      \ifdim \a pt=\b pt \fill[\cc] (\a,0.3) circle (2.5pt); \else \draw[very thick,\cc] (\a,0.3)--(\b,0.3);\fi
      \fill[navy] (12,0) circle (1.5pt); \fill[greenline!60!black] (15,0) circle (1.5pt);
      \node[font=\scriptsize] at (13.7,0.8) {\textbf{\L}};
    \end{scope}
  }
\end{tikzpicture}
\fig{第 3 题选项图：各选项声称的 $z^*$ 所在范围（同一数轴，绿色为正确）}
\end{figure}
\Lb{判断过程}由弱对偶，$12\le z^*$（原问题已有目标值 12 的可行解，最优值不会更小）；又 $z^*\le15$（任一原可行值都不超过对偶可行值 15）。所以 $12\le z^*\le15$。\par
\Lb{各选项}A 方向反了：12 是下界不是上界；C 方向也反了；D 没有依据：区间内任何值都可能，不能断定是中点。\par
\Lb{常见错误}以为 $z^*$ 就是两个数的平均值。\quad\Lb{一句话总结}一对可行解给出最优值的一个“夹逼区间”。

% ---------------- 4 ----------------
\A{4}
\SolHead{C}{基础}{K3 无界性推论}
\begin{figure}[!htbp]\centering
\begin{tikzpicture}[x=1cm,y=1cm]
  \foreach \i/\t in {0/有最优解,1/无界,2/无可行解} {
    \node[font=\scriptsize\bfseries,navy] at (2.2+\i*2.4,0.55) {对偶：\t};
    \node[font=\scriptsize\bfseries,navy,anchor=east] at (0.8,-\i*0.8) {原问题：\t};
  }
  \foreach \r/\c/\t/\s in {0/0/可能/lightgreen,0/1/不可能/lightred,0/2/不可能/lightred,
     1/0/{不可能（B）}/lightred,1/1/{不可能（A）}/lightred,1/2/{必然（C）}/lightgreen,
     2/0/不可能/lightred,2/1/可能/lightgreen,2/2/可能/lightgreen} {
    \node[draw=white,fill=\s,minimum width=2.35cm,minimum height=0.75cm,font=\scriptsize] at (2.2+\c*2.4,-\r*0.8) {\t};
  }
  \draw[shiftc,very thick,rounded corners] (0.95,-1.2) rectangle (8.65,-0.4);
\end{tikzpicture}
\fig{第 4 题解析图：在状态矩阵中，“原问题无界”一行（橙框）只有“对偶无可行解”这一格可能}
\end{figure}
\Lb{判断过程}反证：若对偶有可行解 $\bar y$，由弱对偶，所有原可行解满足 $c\T x\le b\T\bar y$，目标值有上界，与题图中“等值线可以无限外推”矛盾。故对偶无可行解。可以直接验证：对偶为 $\min 2y$，$\st -y\ge2$（$x_1$ 列），$y\ge1$（$x_2$ 列），$y\ge0$；$-y\ge2$ 要求 $y\le-2$，与 $y\ge1$ 矛盾。\par
\Lb{各选项}A、B 都落在橙框中的红格，不可能；D：由推论 3.3 可以判断，并非无法判断。\par
\Lb{常见错误}把逆命题当成立：对偶无可行解时原问题不一定无界。\quad\Lb{一句话总结}一方无界 $\Rightarrow$ 另一方无可行解。

% ---------------- 5 ----------------
\A{5}
\SolHead{A}{基础}{K1 修正单纯形法 $x_B=B^{-1}b$}
\begin{figure}[!htbp]\centering
\begin{tikzpicture}
  \node at (0,0) {$\begin{pmatrix}\frac12&0\\[2pt]-\frac12&1\end{pmatrix}\begin{pmatrix}8\\7\end{pmatrix}=\begin{pmatrix}\frac12\cdot8+0\cdot7\\[2pt]-\frac12\cdot8+1\cdot7\end{pmatrix}=\begin{pmatrix}4\\3\end{pmatrix}$};
  \fill[shiftc!25,opacity=0.5] (-3.05,0.05) rectangle (-1.65,0.55);
  \fill[auxblue!25,opacity=0.5] (-3.05,-0.6) rectangle (-1.65,-0.08);
  \node[font=\scriptsize,shiftc!80!black] at (-2.35,0.95) {第 1 行};
  \node[font=\scriptsize,auxblue] at (-2.35,-0.95) {第 2 行};
\end{tikzpicture}
\fig{第 5 题解析图：$B^{-1}$ 的每一行与 $b$ 做内积}
\end{figure}
\Lb{判断过程}第 1 行 $\frac12\cdot8+0\cdot7=4$；第 2 行 $-\frac12\cdot8+1\cdot7=3$。这就是例 1.1 第二轮的结果 $x_3=4$、$x_5=3$。\par
\Lb{各选项}B 是没乘 $B^{-1}$ 的原 $b$；C 把第 2 行算成 $-\frac12\cdot8+0\cdot7$，用错了元素；D 是行次序颠倒。\par
\Lb{一句话总结}$x_B=B^{-1}b$，$b$ 用原模型的右端项。

% ---------------- 6 ----------------
\A{6}
\SolHead{B}{基础}{K5 从最优表读影子价格}
\Lb{对应公式}$\sigma(s_i)=0-y\T e_i=-y_i$，即 $y_i=-\sigma(s_i)$。
\begin{figure}[!htbp]\centering
\begin{tikzpicture}[x=0.45cm,y=0.45cm]
  \foreach \i/\L/\a/\b/\c/\ok in {0/A/-2/0/-1/0,1/B/2/0/1/1,2/C/0/2/1/0,3/D/1/0/2/0} {
    \begin{scope}[xshift=\i*4.0cm]
      \draw[ax] (-0.3,0)--(4.4,0);
      \draw[black!40] (0,-2.3)--(0,2.5);
      \foreach \y in {-2,2} {\draw[gridl] (0,\y)--(4,\y); \node[left,font=\tiny] at (0,\y) {\y};}
      \ifnum\ok=1 \def\cc{greenline}\else\def\cc{badc!70}\fi
      \fill[\cc] (0.4,0) rectangle (1.2,\a); \fill[\cc] (1.6,0) rectangle (2.4,\b); \fill[\cc] (2.8,0) rectangle (3.6,\c);
      \node[font=\tiny,below] at (0.8,-2.3) {$y_1$}; \node[font=\tiny,below] at (2.0,-2.3) {$y_2$}; \node[font=\tiny,below] at (3.2,-2.3) {$y_3$};
      \node[font=\scriptsize] at (2,3.2) {\textbf{\L}.\ $(\a,\b,\c)$};
    \end{scope}
  }
\end{tikzpicture}
\fig{第 6 题选项图：四个选项的影子价格柱状图（同一纵轴）}
\end{figure}
\Lb{判断过程}检验数 $(-2,0,-1)$ 取相反数得 $(2,0,1)$。验算：$y\T b=2\cdot4+0\cdot9+1\cdot3=11=z^*$ ✓；对偶约束 $2+0+1=3\ge3$、$2+0=2\ge2$ ✓。\par
\Lb{各选项（看图）}A 的柱子朝下：影子价格为负，违反“$\le$ 型问题影子价格非负”，是忘了取负号；C 的高柱错位到 $y_2$——资源 \ding2 有剩余，柱高应为 0；D 把 $y_1,y_3$ 对调了。只有 B 满足非负且 $y_2=0$。\par
\Lb{一句话总结}影子价格 $=$ 松弛变量检验数的相反数。

% ---------------- 7 ----------------
\A{7}
\SolHead{B}{基础}{K4 互补松弛（资源侧）}
\begin{figure}[!htbp]\centering
\smallE{0.7}{\draw[con,shiftc,very thick] (-0.3,3.1)--(6.3,0.9); \node[optpt] at (3,1) {}; \draw[{Stealth}-{Stealth},redline,thick] (3,1)--(3,2); \node[font=\scriptsize,right,redline] at (3.05,1.55) {未触及 \ding2};}
\fig{第 7 题解析图：资源 \ding2 的边界（橙色）没有经过最优点 $(3,1)$}
\end{figure}
\Lb{判断过程}互补松弛：$y_2\cdot s_2=0$。已知 $s_2=3>0$，所以 $y_2=0$。几何上，\ding2 的边界线在最优点上方，最优点不在它上面，把这条线再往外推也不会改变最优点，所以多给材料毫无价值。\par
\Lb{各选项}A：只有资源用尽时影子价格才可能为正；C：$\le$ 型问题影子价格非负；D：把剩余量和价格混为一谈。\par
\Lb{一句话总结}有剩余 $\Rightarrow$ 影子价格为 0。

% ---------------- 8 ----------------
\A{8}
\SolHead{B}{基础}{K6 对偶单纯形法的迭代次序}
\begin{figure}[!htbp]\centering
\begin{tikzpicture}[x=1cm,y=1cm,s/.style={box,minimum width=2.6cm,font=\small}]
  \node[font=\small\bfseries] at (-1.6,0) {单纯形法}; \node[font=\small\bfseries] at (-1.6,-1) {对偶单纯形法};
  \node[s] (a1) at (1.5,0) {① 定入基列}; \node[s] (a2) at (4.6,0) {② 定出基行}; \node[s] (a3) at (7.7,0) {③ 旋转};
  \node[s,fill=shiftc!25] (b1) at (1.5,-1) {① 定出基行}; \node[s] (b2) at (4.6,-1) {② 定入基列}; \node[s] (b3) at (7.7,-1) {③ 旋转};
  \draw[arr] (a1)--(a2); \draw[arr] (a2)--(a3); \draw[arr] (b1)--(b2); \draw[arr] (b2)--(b3);
\end{tikzpicture}
\fig{第 8 题解析图：两种算法每一轮的次序相反}
\end{figure}
\Lb{判断过程}对偶单纯形法要消除 $b$ 列中的负数，所以先在 $b$ 为负的行中选出基行（取最负者），再在该行负系数的列中用 $\min\sigma_j/a_{rj}$ 选入基列。\par
\Lb{各选项}A、C 是单纯形法的第一步；D 是单纯形法定出基行的比值规则。\par
\Lb{一句话总结}对偶单纯形法：先出基，后入基。

% ---------------- 9 ----------------
\A{9}
\SolHead{B}{基础}{K5、K7 影子价格的使用}
\begin{figure}[!htbp]\centering
\begin{minipage}{0.45\linewidth}\centering
\smallE{0.75}{\draw[shiftl] (0.3,4.2)--(5.0,-0.5); \node[badpt] at (3,1) {}; \node[optpt,fill=shiftc] at (3,1.5) {}; \node[lab,right,font=\scriptsize] at (3.1,1.75) {$(3,1.5)$};}
\end{minipage}\hfill
\begin{minipage}{0.52\linewidth}\centering
\begin{tikzpicture}
\begin{axis}[width=7cm,height=4.6cm,xmin=0,xmax=7,ymin=0,ymax=15,axis lines=left,xlabel={$b_1$},ylabel={$z^*$},
  xtick={3,4,4.5,5},ytick={9,11,12,13},label style={font=\scriptsize},tick label style={font=\tiny},grid=major,grid style={black!8}]
  \addplot[navy,very thick] coordinates {(0,0)(3,9)(5,13)(7,13)};
  \addplot[only marks,mark=*,mark options={fill=greenline,draw=white},mark size=2.5pt] coordinates {(4.5,12)};
  \addplot[only marks,mark=x,mark options={draw=badc,thick},mark size=3pt] coordinates {(4.5,11)(4.5,11.5)(4.5,13)};
  \node[font=\tiny,anchor=west] at (axis cs:4.6,12) {B};
  \node[font=\tiny,anchor=west] at (axis cs:4.6,11) {A}; \node[font=\tiny,anchor=west] at (axis cs:4.6,11.5) {C}; \node[font=\tiny,anchor=west] at (axis cs:4.6,13.3) {D};
\end{axis}
\end{tikzpicture}
\end{minipage}
\fig{第 9 题解析图：左为新的最优顶点；右为 $z^*$ 随 $b_1$ 的折线，四个选项中只有 B 落在折线上}
\end{figure}
\Lb{判断过程}$\Delta b_1=0.5\in[-1,1]$（允许范围 $[3,5]$ 换成增量），最优基不变，$z^*$ 增加 $y_1\Delta b_1=2\times0.5=1$，新最优值 12。几何验证：新顶点由 $x_1=3$ 与 $x_1+x_2=4.5$ 交出，为 $(3,1.5)$，$z=9+3=12$ ✓。\par
\Lb{各选项}A 认为最优值不变（误以为资源 \ding1 有剩余）；C 用了 $y_1=1$；D 误把增量算成 $2\times1$。右图中三个红叉都不在折线上。\par
\Lb{一句话总结}范围内：新最优值 $=z^*+y_i\Delta b_i$。

% ---------------- 10 ----------------
\A{10}
\SolHead{B}{基础}{K8 价值系数越界后的续解}
\begin{figure}[!htbp]\centering
\begin{minipage}{0.48\linewidth}\centering
\smallE{0.62}{\draw[iso,badc] (0,3.25)--(4.333,0); \node[badpt] at (3,1) {}; \node[font=\scriptsize,badc] at (3.0,4.6) {A：最优解仍为 $(3,1)$？};}
\end{minipage}\hfill
\begin{minipage}{0.48\linewidth}\centering
\smallE{0.62}{\draw[iso,greenline!60!black] (0,3.625)--(4.833,0); \node[optpt,fill=greenline] at (1.5,2.5) {}; \node[font=\scriptsize,greenline!60!black] at (3.0,4.6) {B：最优点移到 $(1.5,2.5)$};}
\end{minipage}
\fig{第 10 题选项图：$c_2=4$ 时的等值线 $3x_1+4x_2=$ 常数。左图过 $(3,1)$ 的等值线（$z=13$）仍穿过可行域内部，说明还能更大；右图等值线在 $(1.5,2.5)$ 处支撑可行域（$z=14.5$）}
\end{figure}
\Lb{判断过程}$\Delta c_2=2>1$，超出 $[0,3]$。$x_2$ 在第 2 行，$t_{23}=1$、$t_{25}=-1$：$\sigma_3'=-2-2\cdot1=-4$，$\sigma_5'=-1-2\cdot(-1)=1>0$。$b$ 列不受影响（仍可行），$\sigma$ 行出现正数——“原始可行而未达最优”，从原表继续用单纯形法（$x_5$ 入基）。\par
\Lb{各选项（看图）}A：左图过 $(3,1)$ 的等值线下方仍有可行点被切到另一侧，$(1.5,2.5)$ 的目标值 14.5 大于 13，原解不再最优；C：对偶单纯形法用于 $b$ 出现负数，本题 $b$ 没变；D：可行域有界，不会无界。\par
\Lb{一句话总结}改 $c$ 越界 $\to$ $\sigma$ 出现正数 $\to$ 单纯形法。

\subsection*{基础多选题解析}

% ---------------- 11 ----------------
\A{11}
\SolHead{ABC}{基础}{K2 对偶问题的构造}
\begin{figure}[!htbp]\centering
\begin{tikzpicture}[x=1cm,y=1cm,p/.style={box,minimum width=1.6cm,font=\scriptsize},q/.style={gbox,minimum width=1.6cm,font=\scriptsize}]
  \begin{scope}
    \node[p] (a) at (0,0) {第 $i$ 个约束}; \node[q] (b) at (0,-1.2) {第 $i$ 个变量 $y_i$};
    \draw[arr] (a)--(b); \node[font=\scriptsize] at (0,0.75) {\textbf{A} ✓};
  \end{scope}
  \begin{scope}[xshift=3.6cm]
    \node[p] (a) at (0,0) {$x_j\ge0$}; \node[q] (b) at (0,-1.2) {第 $j$ 约束 $\ge c_j$};
    \draw[arr] (a)--(b); \node[font=\scriptsize] at (0,0.75) {\textbf{B} ✓};
  \end{scope}
  \begin{scope}[xshift=7.2cm]
    \node[p] (a) at (0,0) {右端项 $b$}; \node[q] (b) at (0,-1.2) {目标系数};
    \draw[arr] (a)--(b); \node[font=\scriptsize] at (0,0.75) {\textbf{C} ✓};
  \end{scope}
  \begin{scope}[xshift=10.8cm]
    \node[p] (a) at (0,0) {原问题 P}; \node[q] (b) at (0,-1.2) {对偶 D};
    \draw[arr] (a.south east)to[bend left=40](b.north east); \draw[arr] (b.north west)to[bend left=40]node[left,font=\tiny]{再取对偶}(a.south west);
    \node[font=\scriptsize] at (0,0.75) {\textbf{D} ✗：回到 P};
  \end{scope}
\end{tikzpicture}
\fig{第 11 题选项图：每个选项对应的映射关系}
\end{figure}
\Lb{A}正确。原问题有 $m$ 个约束，对偶有 $m$ 个变量，一一对应（$A$ 的行 $\to$ $A\T$ 的列）。\par
\Lb{B}正确。max 问题的正常变量 $x_j\ge0$ 对应 min 问题的正常约束“$\ge$”。\par
\Lb{C}正确。对偶目标为 $b\T y$。\par
\Lb{D}错误。对偶的对偶就是原问题（第 2 章的对称性证明），如图最右：取两次对偶回到起点。\par
\Lb{正确答案}ABC。\quad\Lb{一句话总结}约束 $\leftrightarrow$ 变量，$b\leftrightarrow c$，两次对偶回到原问题。

% ---------------- 12 ----------------
\A{12}
\SolHead{ABC}{基础}{K3 对偶的基本性质}
\begin{figure}[!htbp]\centering
\begin{tikzpicture}[x=0.5cm,y=0.5cm]
  \begin{scope}
    \draw[ax] (0,0)--(6.4,0); \fill[navy] (1.5,0) circle (3pt); \fill[greenline] (4.5,0) circle (3pt);
    \node[font=\tiny,below] at (1.5,0) {$c\T\bar x$}; \node[font=\tiny,below] at (4.5,0) {$b\T\bar y$};
    \node[font=\scriptsize] at (3,1.5) {\textbf{A} ✓：左 $\le$ 右};
  \end{scope}
  \begin{scope}[xshift=4cm]
    \draw[ax] (0,0)--(6.4,0); \node[optpt] at (3,0) {};
    \node[font=\tiny,below] at (3,0) {重合};
    \node[font=\scriptsize] at (3,1.5) {\textbf{B} ✓：重合即最优};
  \end{scope}
  \begin{scope}[xshift=8cm]
    \draw[ax] (0,0)--(6.4,0);
    \foreach \x in {0.6,1.4,2.2} \fill[navy!60] (\x,0) circle (2pt);
    \foreach \x in {3.8,4.6,5.6} \fill[greenline!70] (\x,0) circle (2pt);
    \node[optpt] at (3,0) {};
    \node[font=\scriptsize] at (3,1.5) {\textbf{C} ✓：两侧逼近同一点};
  \end{scope}
  \begin{scope}[xshift=12.2cm,x=0.42cm,y=0.42cm]
    \fill[feasreg] (0,0)--(0,4)--(4,4)--(4,3)--(1,0)--cycle;
    \fill[pattern=north east lines,pattern color=redline!60] (2,0)--(4,2)--(4,0)--cycle;
    \draw[ax] (0,0)--(4.4,0); \draw[ax] (0,0)--(0,4.4);
    \draw[con] (1,0)--(4.2,3.2); \draw[con,redline] (2,0)--(4.2,2.2);
    \node[font=\scriptsize,align=center] at (2.2,5.6) {\textbf{D} ✗：两区域不相交，\\原问题也不可行};
  \end{scope}
\end{tikzpicture}
\fig{第 12 题选项图：A--C 用数轴表示目标值的关系；D 用第 3 章的反例说明“对偶不可行”时原问题也可能不可行}
\end{figure}
\Lb{A}正确，弱对偶性。\quad\Lb{B}正确，推论 3.2（数轴上两点重合）。\quad\Lb{C}正确，强对偶定理。\par
\Lb{D}错误。反例（图右，第 3 章图~\ref{fig:unb}b）：$\max 2x_1-x_2$，$x_1-x_2\le1$，$-x_1+x_2\le-2$，$x\ge0$。原问题的两个半平面（绿色与红色斜线区）不相交，无可行解；其对偶同样矛盾。所以“对偶不可行”推不出“原问题无界”。\par
\Lb{正确答案}ABC。\quad\Lb{一句话总结}无界 $\Rightarrow$ 对方不可行；反过来不成立。

% ---------------- 13 ----------------
\A{13}
\SolHead{AC}{基础}{K5 影子价格的含义与使用}
\begin{figure}[!htbp]\centering
\begin{tikzpicture}
\begin{groupplot}[group style={group size=2 by 1,horizontal sep=1.6cm},width=6.2cm,height=4.2cm,xmin=0,xmax=7,ymin=0,ymax=15,
  axis lines=left,xlabel={$b_1$},ylabel={$z^*$},label style={font=\scriptsize},tick label style={font=\tiny},xtick={3,4,5},ytick={9,11,13}]
\nextgroupplot[title={\scriptsize \textbf{A} ✓：范围内斜率恒为 $y_1=2$}]
  \addplot[navy,very thick] coordinates {(0,0)(3,9)(5,13)(7,13)};
  \addplot[greenline,ultra thick] coordinates {(3,9)(5,13)};
  \addplot[only marks,mark=*,mark options={fill=optc}] coordinates {(4,11)};
\nextgroupplot[title={\scriptsize \textbf{D} ✗：越过 $b_1=5$ 后斜率变 0}]
  \addplot[navy,very thick] coordinates {(0,0)(3,9)(5,13)(7,13)};
  \addplot[badc,dashed,thick] coordinates {(5,13)(7,17)};
  \node[font=\tiny,badc] at (axis cs:5.6,14.6) {外推};
\end{groupplot}
\end{tikzpicture}

\begin{tikzpicture}[x=1cm,y=1cm]
  \draw[ax] (0,0)--(8.4,0) node[right,font=\scriptsize]{元/单位};
  \fill[shiftc] (5,0) circle (3pt) node[above=3pt,font=\scriptsize]{影子价格 $y_i$};
  \fill[greenline] (2.5,0) circle (3pt) node[above=3pt,font=\scriptsize,align=center]{市场价更低：\\买入有利（C ✓）};
  \fill[auxblue] (7.4,0) circle (3pt) node[above=3pt,font=\scriptsize,align=center]{市场价：由供求决定，\\与 $y_i$ 无必然关系（B ✗）};
\end{tikzpicture}
\fig{第 13 题选项图：上排为模型 E 的 $z^*(b_1)$ 折线（A、D），下排为价格数轴（B、C）}
\end{figure}
\Lb{A}正确，$y_i=\partial z^*/\partial b_i$，前提是最优基不变（上左图绿色段）。\par
\Lb{B}错误，影子价格是企业内部的机会成本，由本企业的技术结构决定，不同企业不同。\par
\Lb{C}正确，市场价低于影子价格时，每买 1 单位净增 $y_i-$ 市场价 $>0$（下图绿点）。\par
\Lb{D}错误，上右图：$b_1>5$ 后折线变平，按 $y_1=2$ 外推（红虚线）会高估。\par
\Lb{正确答案}AC。\quad\Lb{一句话总结}影子价格是“范围内的内部边际价值”。

% ---------------- 14 ----------------
\A{14}
\SolHead{ABC}{基础}{K4 互补松弛}
\begin{figure}[!htbp]\centering
\begin{tikzpicture}
  \node at (0,0) {\smallE{0.42}{\draw[con,greenline,very thick] (-0.3,3.1)--(6.3,0.9); \node[optpt] at (3,1) {}; \node[font=\scriptsize] at (3,5.0) {\textbf{A} ✓：\ding2 不过最优点};}};
  \node at (4.1,0) {\begin{tikzpicture}[x=1cm,y=1cm]
     \fill[auxblue!60] (0,0) rectangle (0.6,1.2); \node[font=\tiny,below] at (0.3,0) {$x_1^*=3$};
     \fill[shiftc!70] (1.2,0) rectangle (1.8,0.02); \node[font=\tiny,below] at (1.5,0) {$y_{s1}=0$};
     \node[font=\scriptsize,align=center] at (0.9,2.0) {\textbf{B} ✓：$x_1>0$\\$\Rightarrow$ 对偶剩余为 0};
   \end{tikzpicture}};
  \node at (8.2,0) {\smallE{0.42}{\draw[con,greenline,very thick] (-0.2,4.2)--(4.5,-0.5); \node[optpt] at (3,1) {}; \node[font=\scriptsize] at (3,5.0) {\textbf{C} ✓：\ding1 过最优点};}};
  \node at (12.4,0) {\begin{tikzpicture}[x=0.42cm,y=0.42cm]\fill[feasreg] (0,0)--(3,0)--(3,1)--(0,2)--cycle;
     \draw[ax] (-0.2,0)--(6.3,0); \draw[ax] (0,-0.2)--(0,4.3);
     \draw[con] (-0.2,4.2)--(4.5,-0.5); \draw[con,redline,very thick] (-0.3,2.1)--(6.3,-0.1); \draw[con] (3,-0.2)--(3,4.3);
     \node[optpt] at (3,1) {}; \node[font=\scriptsize,align=center] at (3,5.6) {\textbf{D} ✗：$b_2=6$ 时 \ding2 恰过最优点，\\可取 $y_2=0$ 而资源用尽};\end{tikzpicture}};
\end{tikzpicture}
\fig{第 14 题选项图：A、C 看约束是否经过最优点；B 看产量与对偶剩余；D 是退化反例}
\end{figure}
\Lb{A}正确：资源 \ding2 剩余 $9-6=3>0$，由 $y_2s_2=0$ 得 $y_2=0$（图 A：绿色的 \ding2 线不经过最优点）。\par
\Lb{B}正确：产品侧 $x_1\cdot y_{s1}=0$，$x_1=3>0$，所以 $y_{s1}=0$，第一个对偶约束取等号；验证 $2+0+1=3$ ✓。\par
\Lb{C}正确：由 $y_1(4-x_1-x_2)=0$，$y_1>0$ 时 $x_1+x_2=4$；图 C 中 \ding1 线恰好经过最优点。\par
\Lb{D}错误：互补松弛只保证两个因子至少一个为零。图 D 把 $b_2$ 改为 6，\ding2 线（红）也经过 $(3,1)$，此时最优基 $\{x_1,x_2,x_4\}$ 中 $x_4=0$，资源 \ding2 恰好用尽，而影子价格仍可取 $y_2=0$。\par
\Lb{正确答案}ABC。\quad\Lb{一句话总结}“有剩余 $\Rightarrow$ 价格为 0”成立，反过来不一定。

% ---------------- 15 ----------------
\A{15}
\SolHead{ABD}{基础}{K7、K8、K9 灵敏度分析的传递路径}
\begin{figure}[!htbp]\centering
\begin{tikzpicture}[x=0.42cm,y=0.42cm]
  \foreach \i/\L/\mode in {0/A/1,1/B/2,2/C/3,3/D/4} {
    \begin{scope}[xshift=\i*4.1cm]
      \ifnum\mode=1 \fill[shiftc!50] (0,1) rectangle (1,4); \fill[lightgreen] (1,0) rectangle (6,1); \fi
      \ifnum\mode=2 \fill[shiftc!50] (4,0) rectangle (5,1); \fi
      \ifnum\mode=3 \fill[shiftc!50] (1,0) rectangle (6,1); \fill[lightgreen] (0,1) rectangle (1,4); \fi
      \ifnum\mode=4 \fill[auxblue!25] (0,4) rectangle (7,5); \fill[auxblue!25] (6,0) rectangle (7,5); \fi
      \draw[navy] (0,0) grid (6,4); \ifnum\mode=4 \draw[navy,dashed] (0,4) grid (7,5); \draw[navy,dashed] (6,0) grid (7,4);\fi
      \draw[navy,thick] (1,0)--(1,4); \draw[navy,thick] (0,1)--(6,1);
      \node[font=\tiny] at (0.5,4.4) {$b$}; \node[font=\tiny] at (-0.6,0.5) {$\sigma$};
    \end{scope}
  }
  \node[font=\scriptsize,align=center] at (3,-1.6) {\textbf{A} ✓ 改 $b$：只动 $b$ 列\\$\sigma$ 行（绿）不变};
  \node[font=\scriptsize,align=center] at (12.8,-1.6) {\textbf{B} ✓ 改非基 $c_k$：\\只动 $\sigma_k$ 一格};
  \node[font=\scriptsize,align=center] at (22.6,-1.6) {\textbf{C} ✗ 改基 $c_j$：动整行 $\sigma$，\\$b$ 列（绿）不变};
  \node[font=\scriptsize,align=center] at (32.4,-1.6) {\textbf{D} ✓ 新增约束：\\添一行一列（蓝）};
\end{tikzpicture}
\fig{第 15 题选项图：同一张示意单纯形表上，橙色表示改动波及的区域，绿色表示不受影响的区域}
\end{figure}
\Lb{A}正确：$\sigma\T=c\T-c_B\T B^{-1}A$ 中没有 $b$。\par
\Lb{B}正确：$c_B$ 不含 $c_k$，$y$ 不变，只有 $\sigma_k=c_k-y\T P_k$ 变。\par
\Lb{C}错误：$x_B=B^{-1}b$ 中没有 $c$，基变量取值不变，变的是全部非基检验数与最优值。\par
\Lb{D}正确：原最优解若满足新约束，它在更小的可行域中仍可行，而原问题中它已是最好的，故仍最优（结论 6.3 第 1 步）。\par
\Lb{正确答案}ABD。\quad\Lb{一句话总结}改 $b$ 动可行性，改 $c$ 动最优性，各管一行（列）。

\subsection*{基础填空题解析}

% ---------------- 16 ----------------
\A{16}
\SolHead{$\max w=2y_1+3y_2$；$y_1+y_2\le1$；$y_1+3y_2\le2$}{基础}{K2 min 问题的对偶}
\begin{figure}[!htbp]\centering
\begin{tikzpicture}[x=2.2cm,y=2.2cm]
  \fill[feasreg] (0,0)--(1,0)--(0.5,0.5)--(0,0.667)--cycle;
  \draw[gridl,step=0.25] (0,0) grid (1.4,1.1);
  \draw[ax] (-0.05,0)--(1.5,0) node[right,font=\scriptsize]{$y_1$};
  \draw[ax] (0,-0.05)--(0,1.2) node[above,font=\scriptsize]{$y_2$};
  \foreach \x in {0.5,1} \node[below,font=\tiny] at (\x,0) {\x};
  \foreach \y in {0.5,1} \node[left,font=\tiny] at (0,\y) {\y};
  \draw[con] (0,1)--(1.1,-0.1) node[below,font=\scriptsize]{$y_1+y_2=1$};
  \draw[con] (-0.05,0.683)--(1.4,0.2) node[right,font=\scriptsize]{$y_1+3y_2=2$};
  \draw[iso] (0,0.833)--(1.25,0);
  \node[optpt] at (0.5,0.5) {}; \node[lab,above right,font=\scriptsize] at (0.5,0.52) {$(\frac12,\frac12)$，$w=\frac52$};
\end{tikzpicture}
\fig{第 16 题解析图：对偶问题的可行域（有界），最优点 $(\frac12,\frac12)$}
\end{figure}
\Lb{推理过程}原问题 min、约束 $\ge$（min 的正常约束）$\Rightarrow$ $y\ge0$；变量 $\ge0$ $\Rightarrow$ 对偶约束为 max 的正常约束“$\le$”。目标系数取原右端 $(2,3)$，约束系数取 $A$ 的列：$x_1$ 列 $(1,1)$、$x_2$ 列 $(1,3)$，右端取 $c=(1,2)$。\par
\Lb{图的作用}解析图画出了对偶可行域：它是有界的四边形，最优点在两条边界的交点 $(\frac12,\frac12)$，$w^*=1+\frac32=\frac52$。这与第 24 题用对偶单纯形法求得的 $z^*=\frac52$ 相等，正是强对偶定理。\par
\Lb{易错点}忘了 min 问题要“从右往左”查表，把约束写成 $\ge$。\quad\Lb{等价答案}约束可写为 $y_1+y_2\le1$，$y_1+3y_2\le2$，次序可换。

% ---------------- 17 ----------------
\A{17}
\SolHead{11}{基础}{K5、强对偶 $z^*=y\T b$}
\begin{figure}[!htbp]\centering
\begin{tikzpicture}[x=1cm,y=0.3cm]
  \draw[ax] (0,0)--(5.2,0);
  \fill[auxblue!70] (0.5,0) rectangle (1.5,8); \node[font=\scriptsize] at (1,4) {\color{white}$2\times4$};
  \fill[auxblue!40] (0.5,8) rectangle (1.5,8.05);
  \fill[shiftc!70] (0.5,8) rectangle (1.5,11); \node[font=\scriptsize] at (1,9.5) {$1\times3$};
  \node[font=\scriptsize,right] at (1.6,11) {$=11$};
  \node[font=\scriptsize,below] at (1,0) {$y\T b$};
  \node[font=\scriptsize,align=left,right] at (2.4,6) {资源 \ding1 贡献 8\\资源 \ding2 贡献 $0\times9=0$\\资源 \ding3 贡献 3};
\end{tikzpicture}
\fig{第 17 题解析图：总利润按影子价格分摊到各资源}
\end{figure}
\Lb{推理过程}$z^*=y\T b=2\times4+0\times9+1\times3=11$，与最优表 $z^*=3\times3+2\times1=11$ 一致。\par
\Lb{含义}资源 \ding2 有剩余，分摊为 0；利润全部来自紧缺的 \ding1 和 \ding3。\quad\Lb{易错点}用最优表的 $b$ 列 $(3,1,3)$ 代替原 $b$。

% ---------------- 18 ----------------
\A{18}
\SolHead{$B^{-1}=\frac15\,(3,\,-1;\ -1,\,2)$，即第 1 行 $(\frac35,-\frac15)$，第 2 行 $(-\frac15,\frac25)$}{基础}{K1 $2\times2$ 求逆}
\begin{figure}[!htbp]\centering
\begin{tikzpicture}[x=1cm,y=1cm]
  \node (a) at (0,0) {$\begin{pmatrix}2&1\\1&3\end{pmatrix}$};
  \node (b) at (3.4,0) {$\begin{pmatrix}3&1\\1&2\end{pmatrix}$};
  \node (c) at (6.8,0) {$\begin{pmatrix}3&-1\\-1&2\end{pmatrix}$};
  \node (d) at (10.4,0) {$\frac15\begin{pmatrix}3&-1\\-1&2\end{pmatrix}$};
  \draw[arr] (a)--node[above,font=\scriptsize]{主对角互换}(b);
  \draw[arr] (b)--node[above,font=\scriptsize]{副对角变号}(c);
  \draw[arr] (c)--node[above,font=\scriptsize]{$\div\det=5$}(d);
  \node[font=\scriptsize,gbox] at (5.2,-1.2) {检验：$\begin{pmatrix}2&1\\1&3\end{pmatrix}\cdot\frac15\begin{pmatrix}3&-1\\-1&2\end{pmatrix}=\frac15\begin{pmatrix}5&0\\0&5\end{pmatrix}=I$ ✓};
\end{tikzpicture}
\fig{第 18 题解析图：$2\times2$ 求逆的三步与乘法检验}
\end{figure}
\Lb{易错点}笔记 P3 曾把它写成 $\begin{pmatrix}-\frac25&\frac15\\\frac15&-\frac35\end{pmatrix}$（纠错 2），用乘法检验立刻能发现。

% ---------------- 19 ----------------
\A{19}
\SolHead{$x_4$；$x_2$；$\frac23$}{基础}{K6 出基与入基判据}
\begin{figure}[!htbp]\centering
\begin{tikzpicture}[x=1cm,y=1.6cm]
  \draw[ax] (0,0)--(4.5,0);
  \fill[auxblue!40] (0.6,0) rectangle (1.4,1); \node[font=\scriptsize,above] at (1,1) {$\frac{-1}{-1}=1$};
  \fill[greenline!70] (2.4,0) rectangle (3.2,0.667); \node[font=\scriptsize,above] at (2.8,0.667) {$\frac{-2}{-3}=\frac23$};
  \node[font=\scriptsize,below] at (1,0) {$x_1$}; \node[font=\scriptsize,below] at (2.8,0) {$x_2$};
  \node[font=\scriptsize,align=left,right] at (4.6,0.5) {出基行：$x_4$ 行（$b=-3$ 最负）\\只对负系数列取比值\\最小者 $x_2$ 入基};
\end{tikzpicture}
\fig{第 19 题解析图：出基行中两个负系数列的比值 $\sigma_j/a_{rj}$}
\end{figure}
\Lb{推理过程}$b=(-2,-3)$，最负的是 $-3$，$x_4$ 出基。出基行 $(-1,-3,0,1)$ 中负系数在 $x_1,x_2$：$\frac{\sigma_1}{a_{r1}}=\frac{-1}{-1}=1$，$\frac{\sigma_2}{a_{r2}}=\frac{-2}{-3}=\frac23$。最小者为 $x_2$，$\theta=\frac23$，主元 $-3$。\par
\Lb{易错点}用 $b/a$ 计算比值，或在 $x_4$ 列（系数 $+1$）上也取比值。

% ---------------- 20 ----------------
\A{20}
\SolHead{2}{基础}{K8 非基变量的价值系数}
\begin{figure}[!htbp]\centering
\begin{tikzpicture}[x=1cm,y=1cm]
  \draw[ax] (-3.4,0)--(3.4,0) node[right,font=\scriptsize]{$\Delta$};
  \foreach \x in {-2,0,2} {\draw (\x,0.08)--(\x,-0.08); \node[below,font=\tiny] at (\x,-0.08) {\x};}
  \draw[greenline,very thick,-{Stealth}] (2,0.3)--(-3.3,0.3); \fill[greenline] (2,0.3) circle (2.5pt);
  \node[font=\scriptsize,above] at (-0.8,0.35) {$\sigma_3'=-2+\Delta\le0$};
  \node[font=\scriptsize,redline] at (2.8,0.65) {$\Delta>2$：$\sigma_3'>0$};
\end{tikzpicture}
\fig{第 20 题解析图：$\Delta$ 只有上界 2，向下没有限制}
\end{figure}
\Lb{推理过程}$x_3$ 非基，$\sigma_3'=\sigma_3+\Delta=-2+\Delta\le0\Rightarrow\Delta\le2$。\par
\Lb{含义}$x_3$ 是资源 \ding1 的松弛变量。“给闲置的 \ding1 每单位一个价值 $\Delta$”超过它的影子价格 2 时，留着不用比投入生产更划算，最优解就会改变。\quad\Lb{易错点}套用基变量的公式。

\subsection*{基础解答题解析}

% ---------------- 21 ----------------
\A{21}
\Lb{题意分析}原问题为 max 型，约束与变量都不规矩，按“四步法”逐项对应。\par
\Lb{知识定位}K2 对应规则：max 问题约束“$\le$/$\ge$/$=$” $\to$ $y_i\ge0/\le0/$自由；变量“$\ge0/\le0/$自由” $\to$ 对偶约束“$\ge/\le/=$”。
\begin{figure}[!htbp]\centering
\begin{tikzpicture}[x=1cm,y=1cm,pv/.style={draw=navy,fill=lightblue,rounded corners=2pt,minimum width=1.7cm,minimum height=0.55cm,font=\small},dv/.style={draw=greenline,fill=lightgreen,rounded corners=2pt,minimum width=3.2cm,minimum height=0.55cm,font=\small}]
  \foreach \i/\a/\b in {0/{\le6}/{y_1\ge0},1/{\ge2}/{y_2\le0},2/{=4}/{y_3\text{ 自由}}} {\node[pv] (p\i) at (0,-\i*0.75) {$\a$}; \node[dv,minimum width=2cm] (d\i) at (3.0,-\i*0.75) {$\b$}; \draw[arr] (p\i)--(d\i);}
  \foreach \i/\a/\b in {0/{x_1\ge0}/{y_1+2y_2\ge1},1/{x_2\le0}/{y_1-y_2+y_3\le2},2/{x_3\text{ 自由}}/{y_1+y_3=-1}} {\node[pv] (q\i) at (7.0,-\i*0.75) {$\a$}; \node[dv] (e\i) at (10.6,-\i*0.75) {$\b$}; \draw[arr] (q\i)--(e\i);}
  \node[font=\scriptsize\bfseries] at (1.5,0.7) {第 4 步：约束 $\to$ 变量};
  \node[font=\scriptsize\bfseries] at (8.8,0.7) {第 3 步：变量 $\to$ 约束};
\end{tikzpicture}
\fig{第 21 题解析图：逐项对应的推导路线}
\end{figure}
\Lb{分步求解}\textbf{第 1 步}：对偶目标 $\min w=6y_1+2y_2+4y_3$。\textbf{第 2 步}：max 问题，左 $\to$ 右查表。\textbf{第 3 步}：$A=\begin{pmatrix}1&1&1\\2&-1&0\\0&1&1\end{pmatrix}$，按列：$x_1$ 列 $(1,2,0)$，$x_1\ge0\Rightarrow y_1+2y_2\ge1$；$x_2$ 列 $(1,-1,1)$，$x_2\le0\Rightarrow y_1-y_2+y_3\le2$；$x_3$ 列 $(1,0,1)$，$x_3$ 自由 $\Rightarrow y_1+y_3=-1$。\textbf{第 4 步}：见图左。\par
\Lb{结果}
$\min w=6y_1+2y_2+4y_3$，$\st y_1+2y_2\ge1$，$y_1-y_2+y_3\le2$，$y_1+y_3=-1$，$y_1\ge0$，$y_2\le0$，$y_3$ 无约束。\par
\Lb{结果检查}3 个约束、3 个变量 $\to$ 3 个变量、3 个约束 ✓；对偶约束右端 $(1,2,-1)$ 恰为原目标系数 ✓。\par
\Lb{易错点}$x_2$ 列的系数是 $(1,-1,1)$，注意第二个约束中 $x_2$ 前是负号。\quad\Lb{同类题以后怎么做}先画两张“连线表”，再写式子。

% ---------------- 22 ----------------
\A{22}
\Lb{题意分析}从原点 $(0,0)$ 出发做一轮修正单纯形法，不写整张表。\par
\Lb{已知}$c=(3,2,0,0,0)$，$b=(4,9,3)\T$，$P_1=(1,1,1)\T$，$P_2=(1,3,0)\T$，初始 $B=I$，$c_B=0$。\par
\Lb{分步求解}\textbf{(1)} $y\T=c_B\T B^{-1}=0$，$\sigma_1=3$，$\sigma_2=2$，$x_1$ 入基。$\bar P_1=B^{-1}P_1=(1,1,1)\T$，比值 $\frac41,\frac91,\frac31$，最小为 3，第 3 行的 $x_5$ 出基。\par
\textbf{(2)} 新基按行次序 $\{x_3,x_4,x_1\}$：$B=(P_3,P_4,P_1)=\begin{pmatrix}1&0&1\\0&1&1\\0&0&1\end{pmatrix}$，$B^{-1}=\begin{pmatrix}1&0&-1\\0&1&-1\\0&0&1\end{pmatrix}$（检验 $BB^{-1}=I$ ✓），$x_B=B^{-1}b=(4-3,\,9-3,\,3)\T=(1,6,3)\T$。\par
\textbf{(3)} $c_B=(0,0,3)\T$，$y\T=c_B\T B^{-1}=(0,0,3)$。非基变量：$\sigma_2=2-y\T P_2=2-0=2>0$；$\sigma_5=0-y\T e_3=-3$。$\sigma_2>0$，未达最优，下一轮 $x_2$ 入基。
\begin{figure}[!htbp]\centering
\begin{minipage}{0.4\linewidth}\centering
\smallE{0.7}{\node[vpt] at (0,0) {}; \draw[arr,optc,very thick] (0,0)--(2.9,0); \node[optpt] at (3,0) {}; \node[lab,above,font=\scriptsize] at (3.3,0.15) {$(3,0)$}; \draw[arr,optc,dashed] (3,0)--(3,0.9); \node[font=\scriptsize,optc] at (4.2,0.6) {下一轮};}
\end{minipage}\hfill
\begin{minipage}{0.56\linewidth}\small
\begin{tabular}{@{}lccc@{}}\toprule
 & 第一轮前 & 第一轮后 & \\\midrule
基 & $\{x_3,x_4,x_5\}$ & $\{x_3,x_4,x_1\}$ & \\
$x_B$ & $(4,9,3)$ & $(1,6,3)$ & \\
$y$ & $(0,0,0)$ & $(0,0,3)$ & \\
$\sigma_{\text{非基}}$ & $\sigma_1=3,\ \sigma_2=2$ & $\sigma_2=2,\ \sigma_5=-3$ & \\
$z$ & 0 & 9 & \\\bottomrule
\end{tabular}
\end{minipage}
\fig{第 22 题解析图：基解从原点沿 $x_1$ 轴移到 $(3,0)$，$z$ 由 0 升到 9；下一轮将沿设备约束上移}
\end{figure}
\Lb{结果检查}$(3,0)$：\ding1 剩 $1=x_3$，\ding2 剩 $6=x_4$，\ding3 用尽 ✓。\par
\Lb{易错点}$\theta$ 只对 $\bar P_1$ 中正系数行取比值；$B$ 的列次序必须与行次序 $\{x_3,x_4,x_1\}$ 一致。\par
\Lb{同类题以后怎么做}每一轮固定算四样：$x_B$、$y$、$\sigma_{\text{非基}}$、$\bar P_k$。

% ---------------- 23 ----------------
\A{23}
\Lb{题意分析}比较租金与影子价格，并核对增量是否在允许范围内。\par
\Lb{已知}$y_3=1$，$b_3=3\in[1.5,4]$，租金 0.8。\quad\Lb{知识定位}K5 影子价格的用法；K7 允许范围。\par
\Lb{分步求解}\textbf{第 1 步}：$0.8<1$，租入有利。\textbf{第 2 步}：$b_3$ 上限为 4，最多租 $4-3=1$ 台时。\textbf{第 3 步}：净增利润 $=1\times(1-0.8)=0.2$。\textbf{第 4 步}：$x_B'=x_B+\Delta b_3\beta_3=(3,1,3)\T+1\cdot(1,-1,2)\T=(4,0,5)\T$，即 $x_1=4$，$x_2=0$，$z=12$（比 11 多 1）。
\begin{figure}[!htbp]\centering
\smallE{0.75}{\draw[shiftl] (4,-0.2)--(4,4.3) node[above,font=\scriptsize]{$x_1=4$}; \fill[shiftc!20] (3,0)--(4,0)--(3,1)--cycle; \node[badpt] at (3,1) {}; \node[optpt,fill=shiftc] at (4,0) {}; \node[lab,below right,font=\scriptsize] at (4,0) {$(4,0)$，$z=12$};}
\fig{第 23 题解析图：设备边界由 $x_1=3$ 右移到 $x_1=4$，可行域多出橙色三角形，最优点移到 $(4,0)$}
\end{figure}
\Lb{结果检查}在 $(4,0)$：\ding1 $4\le4$（恰好用尽），\ding2 $4\le9$，\ding3 $4\le4$ ✓。若再多租，$b_3>4$ 后 \ding1 成为唯一瓶颈，最优点停在 $(4,0)$，设备的影子价格变为 0。\par
\Lb{易错点}不核对范围，以为租得越多越好。\quad\Lb{同类题以后怎么做}先比价格，再看范围，最后乘增量。

% ---------------- 24 ----------------
\A{24}
\Lb{题意分析}min 问题、系数非负、约束全为 $\ge$，乘 $-1$ 后起始表对偶可行。\par
\Lb{分步求解}\textbf{起始表}见第 19 题。\textbf{第一轮}：$x_4$ 出基，$x_2$ 入基（$\theta=\frac23$），主元 $-3$。$x_2$ 行 $=(\frac13,1,0,-\frac13\mid1)$；$x_3$ 行加上新行的 1 倍：$(-\frac23,0,1,-\frac13\mid-1)$；$\sigma'=\sigma-\frac23\times$（原 $x_4$ 行）：$\sigma_1=-1-\frac23\cdot(-1)=-\frac13$，$\sigma_4=0-\frac23\cdot1=-\frac23$。\par
\textbf{第二轮}：$b=(-1,1)$，$x_3$ 出基；负系数在 $x_1$（$-\frac23$）与 $x_4$（$-\frac13$）：比值 $\frac{-1/3}{-2/3}=\frac12$，$\frac{-2/3}{-1/3}=2$，$x_1$ 入基，主元 $-\frac23$。新 $x_1$ 行 $=(1,0,-\frac32,\frac12\mid\frac32)$；$x_2$ 行减去新行的 $\frac13$ 倍：$(0,1,\frac12,-\frac12\mid\frac12)$；$\sigma_3=-\frac12$，$\sigma_4=-\frac12$。\par
\textbf{结果}：$b=(\frac32,\frac12)\ge0$，最优。$x^*=(\frac32,\frac12)$，$z^*=\frac32+2\cdot\frac12=\frac52$。
\begin{figure}[!htbp]\centering
\begin{tikzpicture}[x=1.1cm,y=1.1cm]
  \fill[feasreg] (0,3)--(0,2)--(1.5,0.5)--(3,0)--(4,0)--(4,3)--cycle;
  \draw[gridl] (0,0) grid (4,3);
  \draw[ax] (-0.2,0)--(4.3,0) node[right,font=\scriptsize]{$x_1$};
  \draw[ax] (0,-0.2)--(0,3.3) node[above,font=\scriptsize]{$x_2$};
  \foreach \i in {1,2,3,4} \node[below,font=\tiny] at (\i,0) {\i};
  \foreach \i in {1,2,3} \node[left,font=\tiny] at (0,\i) {\i};
  \draw[con] (0,2)--(2.2,-0.2); \draw[con] (0,1)--(4,-0.333);
  \draw[iso] (0,1.25)--(2.5,0);
  \draw[arr,badc,very thick,dashed] (0,0)--(0,0.93);
  \draw[arr,optc,very thick] (0.05,1)--(1.43,0.53);
  \node[badpt] at (0,0) {}; \node[badpt] at (0,1) {}; \node[optpt] at (1.5,0.5) {};
  \node[lab,right,font=\scriptsize,text=badc] at (0.1,0.3) {表一 $(0,0)$};
  \node[lab,left,font=\scriptsize,text=badc] at (-0.1,1.2) {表二 $(0,1)$};
  \node[lab,above right,font=\scriptsize,text=optc] at (1.5,0.55) {表三 $(\frac32,\frac12)$};
\end{tikzpicture}
\fig{第 24 题解析图：前两个基解在可行域外（$(0,1)$ 满足 $x_1+3x_2\ge3$ 但不满足 $x_1+x_2\ge2$），第三个落在顶点上}
\end{figure}
\Lb{结果检查}$\frac32+\frac12=2$ ✓，$\frac32+\frac32=3$ ✓，两个约束都取等号；对偶最优解 $y=(\frac12,\frac12)$（$x_3,x_4$ 检验数取负），$w=1+\frac32=\frac52=z^*$ ✓（与第 16 题吻合）。\par
\Lb{易错点}最后一步忘记把 $\max z'$ 换回 $\min z$ 的符号（纠错 4）。\par
\Lb{同类题以后怎么做}每轮写出 $b$ 列中负数的个数，它应逐步减少到 0。

% ---------------- 25 ----------------
\A{25}
\Lb{题意分析}右端项范围：只核验可行性。\par
\Lb{已知}$\beta_3=x_5$ 列 $=(1,-1,2)\T$，$x_B=(3,1,3)\T$（次序 $x_1,x_2,x_4$）。\par
\Lb{分步求解}$x_1$ 行：$3+\Delta b_3\ge0\Rightarrow\Delta b_3\ge-3$；$x_2$ 行：$1-\Delta b_3\ge0\Rightarrow\Delta b_3\le1$；$x_4$ 行：$3+2\Delta b_3\ge0\Rightarrow\Delta b_3\ge-1.5$。交集 $-1.5\le\Delta b_3\le1$，即 $b_3\in[1.5,4]$。
\begin{figure}[!htbp]\centering
\smallE{0.8}{\draw[shiftl] (1.5,-0.2)--(1.5,4.3) node[above,font=\scriptsize]{$b_3=1.5$}; \draw[shiftl] (4,-0.2)--(4,4.3) node[above,font=\scriptsize]{$b_3=4$};
  \node[optpt,fill=shiftc] at (1.5,2.5) {}; \node[optpt,fill=shiftc] at (4,0) {}; \node[optpt] at (3,1) {};
  \draw[arr,optc,thick] (1.6,2.4)--(2.9,1.1); \draw[arr,optc,thick] (3.1,0.9)--(3.9,0.1);
  \node[lab,above right,font=\scriptsize] at (1.55,2.55) {$(1.5,2.5)$：碰到 \ding2}; \node[lab,below,font=\scriptsize] at (4,-0.1) {$(4,0)$：碰到 $x_2=0$};}
\fig{第 25 题解析图：设备边界在两条橙色虚线之间平移时，最优点沿 \ding1 线滑动}
\end{figure}
\Lb{端点含义}$b_3=1.5$ 时最优点滑到 $(1.5,2.5)$，恰好碰上 \ding2 的边界（$x_4=0$），再往左 \ding2 就成为紧约束；$b_3=4$ 时滑到 $(4,0)$，碰上 $x_2=0$，再往右 $x_2$ 会变负，最优基改变。\par
\Lb{易错点}$\beta_3$ 要取 $x_5$ 列（第 3 个约束的松弛变量），而不是第 3 行。\quad\Lb{同类题以后怎么做}“正系数给下界，负系数给上界，取交集”。

% ============================================================
%  提升 10 题与跨学科 3 题详细解析
% ============================================================
\section*{中等偏下提升训练详细解析}
\addcontentsline{toc}{section}{详细解析：中等偏下提升（26--35）}
\markboth{详细解析：提升训练}{}

\subsection*{提升单选题解析}

% ---------------- 26 ----------------
\A{26}
\SolHead{B}{中等偏下}{K5 影子价格 + K7 允许范围}
\Lb{知识链}核对范围 $\to$ 发现越界 $\to$ 对偶单纯形法续解 $\to$ 比较最优值。\par
\Lb{判断过程}$\Delta b_1=2$，超出 $[-1,1]$，不能直接用 $2\times2=4$。传递公式给出 $x_B'=(3,1,3)\T+2\cdot(0,1,-3)\T=(3,3,-3)\T$，$x_4$ 行变负。用对偶单纯形法：$x_4$ 行 $(0,0,-3,1,2\mid-3)$ 中只有 $x_3$ 列为负，$\theta=\frac{-2}{-3}=\frac23$，$x_3$ 入基。旋转后 $x_3=1$，$x_2=3-1=2$，$x_1=3$，检验数 $\sigma_4=-\frac23$，$\sigma_5=-\frac73$，最优。新最优点 $(3,2)$，$z^*=9+4=13$，实际增加 $13-11=2$。
\begin{figure}[!htbp]\centering
\begin{tikzpicture}
\begin{groupplot}[group style={group size=4 by 1,horizontal sep=0.55cm},width=4.3cm,height=3.6cm,xmin=3.5,xmax=6.5,ymin=10,ymax=18,
  axis lines=left,xtick={4,5,6},ytick={11,13,15,17},tick label style={font=\tiny},xlabel={\tiny $b_1$}]
\nextgroupplot[title={\scriptsize A：按 $y_1=2$ 外推}]
  \addplot[black!25,thick] coordinates {(3.5,10)(5,13)(6.5,13)};
  \addplot[badc,very thick] coordinates {(4,11)(6,15)}; \addplot[only marks,mark=*,mark options={fill=badc}] coordinates {(6,15)};
\nextgroupplot[title={\scriptsize B：实际折线 ✓}]
  \addplot[greenline,very thick] coordinates {(3.5,10)(5,13)(6.5,13)}; \addplot[only marks,mark=*,mark options={fill=greenline}] coordinates {(6,13)};
\nextgroupplot[title={\scriptsize C：误认为无增量}]
  \addplot[black!25,thick] coordinates {(3.5,10)(5,13)(6.5,13)};
  \addplot[badc,very thick] coordinates {(4,11)(6,11)}; \addplot[only marks,mark=*,mark options={fill=badc}] coordinates {(6,11)};
\nextgroupplot[title={\scriptsize D：误用斜率 3}]
  \addplot[black!25,thick] coordinates {(3.5,10)(5,13)(6.5,13)};
  \addplot[badc,very thick] coordinates {(4,11)(6,17)}; \addplot[only marks,mark=*,mark options={fill=badc}] coordinates {(6,17)};
\end{groupplot}
\end{tikzpicture}

\smallE{0.7}{\draw[shiftl] (1.8,4.2)--(6.3,-0.3); \fill[shiftc!20] (3,1)--(3,2)--(1.5,2.5)--cycle; \node[badpt] at (3,1) {}; \node[optpt,fill=shiftc] at (3,2) {}; \node[lab,right,font=\scriptsize] at (3.1,2.2) {$(3,2)$：\ding2 与 \ding3 的交点};}
\fig{第 26 题解析图：上排为四个选项隐含的 $z^*(b_1)$ 路径（灰线为真实折线，共同坐标）；下图为 $b_1=6$ 时的新最优点}
\end{figure}
\Lb{各选项（看图）}A 沿斜率 2 直线外推，忽略了 $b_1=5$ 处的折点；C 认为资源 \ding1 无价值；D 把 $c_1=3$ 当成影子价格。只有 B 的点落在真实折线上。下图说明：$b_1$ 超过 5 后，\ding1 的边界越过 \ding2 与 \ding3 的交点 $(3,2)$，不再起约束作用，所以增量封顶。\par
\Lb{常见错误}不查范围直接乘影子价格。\quad\Lb{一句话总结}越界就重算，影子价格不能外推。

% ---------------- 27 ----------------
\A{27}
\SolHead{A}{中等偏下}{K4 互补松弛 + K3 强对偶}
\Lb{判断过程}对偶：$\min w=5y_1+8y_2$，$\st y_1+2y_2\ge4$，$y_1+y_2\ge3$，$y\ge0$。$x_1=3>0$、$x_2=2>0$，两条对偶约束取等号：$y_1+2y_2=4$，$y_1+y_2=3$，解得 $y_2=1$，$y_1=2$。验算 $w=10+8=18=z=12+6$ ✓。
\begin{figure}[!htbp]\centering
\begin{tikzpicture}[x=0.45cm,y=0.45cm]
  \foreach \i/\L/\px/\py/\st in {0/A/2/1/0,1/B/1/2/1,2/C/0/4/1,3/D/3/0/2} {
    \begin{scope}[xshift=\i*3.5cm]
      \fill[feasreg] (0,5)--(0,3)--(2,1)--(4,0)--(5,0)--(5,5)--cycle;
      \draw[ax] (-0.2,0)--(5.4,0) node[right,font=\tiny]{$y_1$}; \draw[ax] (0,-0.2)--(0,5.4) node[above,font=\tiny]{$y_2$};
      \draw[con] (0,2)--(4.4,-0.2); \draw[con] (0,3)--(3.2,-0.2);
      \draw[iso] (0,2.25)--(3.6,0);
      \ifnum\st=0 \node[optpt,fill=greenline] at (\px,\py) {};\fi
      \ifnum\st=1 \node[optpt,fill=auxblue] at (\px,\py) {};\fi
      \ifnum\st=2 \node[badpt] at (\px,\py) {};\fi
      \node[font=\scriptsize,align=center] at (2.5,6.3) {\textbf{\L}.\ $(\px,\py)$};
    \end{scope}
  }
  \node[font=\scriptsize,align=center] at (2.5,-1.4) {可行且 $w=18$\\在最优等值线上 ✓};
  \node[font=\scriptsize,align=center] at (10.3,-1.4) {可行，$w=21$\\不在最优等值线上};
  \node[font=\scriptsize,align=center] at (18.1,-1.4) {可行，$w=32$\\离最优更远};
  \node[font=\scriptsize,align=center] at (25.9,-1.4) {$y_1+2y_2=3<4$\\不可行};
\end{tikzpicture}
\fig{第 27 题选项图：四个选项画在同一张对偶可行域图上（虚线为最优等值线 $5y_1+8y_2=18$）}
\end{figure}
\Lb{各选项}B、C 都对偶可行，但目标值大于 18，不是最优；D 不满足 $y_1+2y_2\ge4$，在可行域外。\par
\Lb{常见错误}把对偶解写成 $(y_2,y_1)$ 的次序。\quad\Lb{一句话总结}已知原最优解，用“产量为正 $\Rightarrow$ 对偶约束取等”列方程。

% ---------------- 28 ----------------
\A{28}
\SolHead{C}{中等偏下}{K6 入基判据}
\Lb{判断过程}出基行负系数在 $x_1$（$-2$）、$x_3$（$-4$）：$\frac{-3}{-2}=1.5$，$\frac{-2}{-4}=0.5$，取最小者 $x_3$ 入基，主元 $-4$。旋转后 $b_r'=\frac{-6}{-4}=1.5>0$，$\sigma'=\sigma-0.5\,a_r=(-2,-1.5,0,0,-0.5)\le0$ ✓。
\begin{figure}[!htbp]\centering
\begin{tikzpicture}[x=0.4cm,y=0.42cm]
  \foreach \i/\L/\vv/\a/\b/\c/\d/\e/\note/\ok in {
     0/A/x_1/0/-2.5/4/0/-1.5/{$\sigma_3'=4>0$}/0,
     1/B/x_2/-5/0/-6/0/1/{$\sigma_5'=1>0$，且 $b_r'=-6$}/0,
     2/C/x_3/-2/-1.5/0/0/-0.5/{全部 $\le0$ ✓}/1} {
    \begin{scope}[xshift=\i*4.4cm]
      \draw[ax] (-0.3,0)--(6.6,0);
      \draw[black!40] (0,-6.3)--(0,4.5);
      \foreach \y in {-6,4} {\draw[gridl] (0,\y)--(6.3,\y); \node[left,font=\tiny] at (0,\y) {\y};}
      \ifnum\ok=1 \def\cc{greenline}\else\def\cc{badc!70}\fi
      \foreach \v [count=\k from 0] in {\a,\b,\c,\d,\e} {\fill[\cc] (0.3+\k*1.2,0) rectangle (1.1+\k*1.2,\v);}
      \foreach \k in {1,...,5} \node[font=\tiny,below] at (-0.5+\k*1.2,-6.3) {$x_\k$};
      \node[font=\scriptsize] at (3.2,5.4) {\textbf{\L}：选 $\vv$ 入基后的 $\sigma'$};
      \node[font=\scriptsize] at (3.2,-7.9) {\note};
    \end{scope}
  }
  \begin{scope}[xshift=12.6cm]
    \node[font=\scriptsize,align=left,text width=3.0cm] at (4,0) {\textbf{D}：出基行中有负系数（$x_1,x_3$），存在可选的入基列，不是“无可行解”。只有当出基行全部 $\ge0$ 时才无可行解。};
  \end{scope}
\end{tikzpicture}
\fig{第 28 题选项图：分别选 $x_1,x_2,x_3$ 入基后的新检验数行（同一纵轴）}
\end{figure}
\Lb{各选项（看图）}A 的柱状图中 $x_3$ 一柱冲到正值 4：比值不是最小，旋转后破坏了对偶可行；B 选了正系数列：主元为正，$b_r'=-6$ 仍为负，而且 $\sigma_5'=1>0$；C 全部柱子都在 0 以下。\par
\Lb{一句话总结}只在负系数列中比较，取 $\sigma_j/a_{rj}$ 最小者。

\subsection*{提升多选题解析}

% ---------------- 29 ----------------
\A{29}
\SolHead{ABD}{中等偏下}{K5、K7、K8 综合}
\begin{figure}[!htbp]\centering
\begin{tikzpicture}
  \node at (0,0) {\smallE{0.42}{\draw[shiftl] (-0.3,2.43)--(6.3,0.23); \node[optpt] at (3,1) {}; \node[font=\scriptsize] at (3,5.1) {\textbf{A} ✓：\ding2 下移仍不过最优点};}};
  \node at (4.1,0) {\smallE{0.42}{\draw[iso,greenline!60!black] (0.36,4.3)--(3.8,0); \node[optpt] at (3,1) {}; \node[font=\scriptsize] at (3,5.1) {\textbf{B} ✓：$c_1=2.5$ 等值线仍支撑 $(3,1)$};}};
  \node at (8.2,0) {\smallE{0.42}{\fill[shiftc!20] (3,0)--(4,0)--(3,1)--cycle; \draw[shiftl] (5,-0.2)--(5,4.3); \node[optpt,fill=shiftc] at (4,0) {}; \node[font=\scriptsize,align=center] at (3,5.3) {\textbf{C} ✗：$b_3=5$ 时最优点 $(4,0)$\\$z=12$，只增 1};}};
  \node at (12.3,0) {\smallE{0.42}{\draw[iso,greenline!60!black] (0.25,4.3)--(3.833,0); \node[optpt] at (3,1) {}; \node[font=\scriptsize] at (3,5.1) {\textbf{D} ✓：$c_2=2.5$ 仍支撑 $(3,1)$};}};
\end{tikzpicture}
\fig{第 29 题选项图：四个选项画在同一模型 E 底图上}
\end{figure}
\Lb{A}正确。$\Delta b_2=-2\in[-3,+\infty)$，最优基不变；$y_2=0$，最优值不变；$x_B'=(3,1,3-2)=(3,1,1)$，只是材料剩余由 3 变为 1。图 A：\ding2 线下移到 $x_1+3x_2=7$，仍在 $(3,1)$ 上方（$3+3=6<7$）。\par
\Lb{B}正确。$c_1=2.5\in[2,+\infty)$，最优解不变，$z=2.5\times3+2\times1=9.5$。\par
\Lb{C}错误。$\Delta b_3=2>1$，越界。$b_3=5$ 时 \ding3 线右移到 $x_1=5$，但 \ding1 先把最优点卡在 $(4,0)$，$z=12$，只增加 1。\par
\Lb{D}正确。$c_2=2.5\in[0,3]$，等值线 $3x_1+2.5x_2=11.5$ 仍在 $(3,1)$ 处支撑可行域。\par
\Lb{哪个条件导致成立与否}A、B、D 的改动都落在各自的允许范围内；C 的改动越过了上界 4。\quad\Lb{正确答案}ABD。

% ---------------- 30 ----------------
\A{30}
\SolHead{ABD}{中等偏下}{K9 新增变量 + K5 机会成本}
\Lb{核心计算}$y\T P_6=2\cdot1+0\cdot2+1\cdot0=2$，所以 $\sigma_6=c_6-2$。
\begin{figure}[!htbp]\centering
\begin{tikzpicture}[x=1cm,y=1cm]
  \draw[ax] (0,0)--(9.4,0);
  \draw[dashed,black!50] (0,2)--(9,2) node[right,font=\scriptsize]{机会成本 2};
  \fill[auxblue!60] (0.5,0) rectangle (1.3,2); \node[font=\scriptsize,below,align=center] at (0.9,0) {$y\T P_6$\\（$=2\cdot1$）};
  \fill[badc!60] (2.3,0) rectangle (3.1,1.5); \node[font=\scriptsize,below,align=center] at (2.7,0) {\textbf{A}：$c_6=1.5$\\$\sigma_6<0$，不生产 ✓};
  \fill[greenline!70] (4.3,0) rectangle (5.1,3); \node[font=\scriptsize,below,align=center] at (4.7,0) {\textbf{B}：$c_6=3$\\$\sigma_6=1>0$，生产 ✓};
  \node[box,font=\scriptsize,text width=2.4cm] at (7.4,1.0) {\textbf{D} ✓：$x_6=0$ 时\\原 $x_B$ 不变，仍 $\ge0$};
  \node[rbox,font=\scriptsize,text width=2.4cm] at (7.4,2.9) {\textbf{C} ✗：只需算一个\\检验数};
\end{tikzpicture}
\fig{第 30 题选项图：利润柱与机会成本线的比较}
\end{figure}
\Lb{A}正确：$\sigma_6=-0.5<0$。\quad\Lb{B}正确：$\sigma_6=1>0$；把 $B^{-1}P_6$ 写进表后继续单纯形法。$B^{-1}P_6$ 的计算：$B^{-1}=\begin{pmatrix}0&0&1\\1&0&-1\\-3&1&2\end{pmatrix}$（最优表 $x_3,x_4,x_5$ 三列），$B^{-1}P_6=(0,1,-1)\T$，比值只有 $x_2$ 行 $\frac11=1$，$x_2$ 出基。\par
\Lb{C}错误：判断只需 $\sigma_6=c_6-y\T P_6$，不必重解。\quad\Lb{D}正确：新变量取 0，原约束与取值都不变。\par
\Lb{正确答案}ABD。\quad\Lb{一句话总结}新产品利润高于其资源的影子价值之和才值得生产。

\subsection*{提升填空题解析}

% ---------------- 31 ----------------
\A{31}
\SolHead{1；4；12}{中等偏下}{K5 + K7}
\Lb{推理过程}影子价格 $y_3=-\sigma_5=1$。范围（第 25 题）$b_3\in[1.5,4]$，最多增至 4，$z^*=11+1\times(4-3)=12$，最优点 $(4,0)$。
\begin{figure}[!htbp]\centering
\begin{tikzpicture}
\begin{axis}[width=7.5cm,height=4.2cm,xmin=0,xmax=6,ymin=0,ymax=14,axis lines=left,xlabel={$b_3$},ylabel={$z^*$},
  xtick={1.5,3,4},ytick={6,9.5,11,12},label style={font=\scriptsize},tick label style={font=\tiny},grid=major,grid style={black!8}]
  \addplot[navy,very thick] coordinates {(0,6)(1.5,9.5)(4,12)(6,12)};
  \addplot[greenline,ultra thick] coordinates {(1.5,9.5)(4,12)};
  \addplot[only marks,mark=*,mark options={fill=optc}] coordinates {(3,11)};
  \addplot[only marks,mark=*,mark options={fill=shiftc}] coordinates {(4,12)};
  \node[font=\tiny,anchor=north west] at (axis cs:2.2,10.2) {斜率 $y_3=1$};
\end{axis}
\end{tikzpicture}
\fig{第 31 题解析图：$z^*$ 随 $b_3$ 的折线，绿色段斜率为影子价格 1，右端点 $b_3=4$ 处 $z^*=12$}
\end{figure}
\Lb{图的说明}$b_3\le1.5$ 时最优点在 \ding2 线上 $(b_3,\frac{9-b_3}{3})$，$z^*=3b_3+6-\frac23b_3$，斜率 $\frac73$；$b_3\ge4$ 后 \ding1 成为唯一瓶颈，折线变平。\quad\Lb{易错点}答成 $b_3$ 最多增至 $3+4=7$（把范围端点当增量）。

% ---------------- 32 ----------------
\A{32}
\SolHead{$(3,0.5)$；10}{中等偏下}{K9 新增约束 + K6 对偶单纯形法}
\Lb{推理过程}原解 $x_2=1>0.5$，不满足。新约束 $x_2+x_6=0.5$；由最优表 $x_2$ 行 $x_2+x_3-x_5=1$ 消去 $x_2$：$x_6-x_3+x_5=-0.5$。$x_6$ 出基，负系数只有 $x_3$（$-1$），$\theta=\frac{-2}{-1}=2$，$x_3$ 入基：$x_3=0.5$，$x_2=1-0.5=0.5$，$x_4=3+3\times0.5=4.5$，$x_1=3$。$\sigma_3'=0$，$\sigma_5'=-1-2\times1=-3$，$\sigma_6'=-2$，最优。$z=9+1=10$。
\begin{figure}[!htbp]\centering
\smallE{0.8}{\fill[shiftc!25] (0,0)--(3,0)--(3,0.5)--(0,0.5)--cycle; \draw[shiftl] (-0.2,0.5)--(6.3,0.5) node[right,font=\scriptsize]{$x_2=0.5$}; \node[badpt] at (3,1) {}; \node[optpt,fill=shiftc] at (3,0.5) {}; \draw[arr,optc] (3,0.95)--(3,0.6); \node[lab,right,font=\scriptsize] at (3.15,0.45) {$(3,0.5)$，$z=10$};}
\fig{第 32 题解析图：新约束切掉可行域的上半部分（剩下橙色窄条），最优点沿 \ding3 线下移}
\end{figure}
\Lb{合理性检查}$z$ 减少 $11-10=1=y_{\text{新}}\times0.5=2\times0.5$；加约束只会使最优值变差 ✓。\quad\Lb{易错点}忘记消去新行中的基变量 $x_2$。

% ---------------- 33 ----------------
\A{33}
\SolHead{$[2,+\infty)$}{中等偏下}{K8 基变量的价值系数}
\Lb{推理过程}$x_1$ 在第 1 行，非基列 $t_{13}=0$、$t_{15}=1$：$x_3$ 列 $-2-\Delta c_1\cdot0\le0$ 恒成立；$x_5$ 列 $-1-\Delta c_1\cdot1\le0\Rightarrow\Delta c_1\ge-1$。所以 $c_1\ge2$。
\begin{figure}[!htbp]\centering
\smallE{0.8}{\draw[very thick,dashed,shiftc] (-0.2,4.2)--(4.5,-0.5); \node[font=\scriptsize,shiftc!80!black] at (5.4,3.4) {$c_1=2$：与 \ding1 重合}; \draw[very thick,dashed,redline] (3,-0.3)--(3,4.4); \node[font=\scriptsize,redline] at (4.3,4.1) {$c_1\to\infty$：趋于竖直}; \node[optpt] at (3,1) {}; \fill[shiftc!50,opacity=.6] (1.5,2.5)--(3,1)--(3.1,1.1)--(1.6,2.6)--cycle;}
\fig{第 33 题解析图：等值线 $c_1x_1+2x_2=$ 常数随 $c_1$ 旋转；$c_1=2$ 时与 \ding1 平行，整条棱（橙色）都最优}
\end{figure}
\Lb{含义}$c_1<2$ 时等值线比 \ding1 线更平，最优点跳到 $(1.5,2.5)$；$c_1$ 再大，等值线趋于竖直，但 $x_1$ 已被设备卡在 3，最优点不变——所以没有上界。\quad\Lb{易错点}套用非基变量公式得到“只有上界”。

\subsection*{提升综合大题解析}

% ---------------- 34 ----------------
\A{34}
\Lb{1. 题目到底在问什么}已知面包坊的最优方案，要求写对偶、求三种资源的“内部价格”、判断租烤箱是否划算、求面粉的稳定范围。\par
\Lb{2. 已知条件整理}
\begin{center}\small
\begin{tabular}{@{}lccccc@{}}\toprule
资源 & 面包 & 蛋糕 & 可用量 & 最优方案耗用 & 剩余\\\midrule
面粉 & 1 & 2 & 8 & $2+6=8$ & 0\\
烤箱 & 3 & 2 & 12 & $6+6=12$ & 0\\
奶油 & 0 & 1 & 3.5 & 3 & 0.5\\\midrule
利润 & 5 & 4 & & $z^*=10+12=22$ & \\\bottomrule
\end{tabular}
\end{center}
\Lb{3. 知识链定位}对偶构造（K2）$\to$ 互补松弛（K4）$\to$ 影子价格的使用（K5）$\to$ 右端项范围（K7）。

\Lb{4. 第(1)问}对称形式直接转置：
$\min w=8y_1+12y_2+3.5y_3$，$\st y_1+3y_2\ge5$（面包），$2y_1+2y_2+y_3\ge4$（蛋糕），$y_1,y_2,y_3\ge0$。

\Lb{5. 第(2)问为什么这样做}已知原最优解，用互补松弛最省事。$x_1=2>0$、$x_2=3>0$ $\Rightarrow$ 两条对偶约束取等；奶油剩 0.5 $\Rightarrow$ $y_3=0$。于是
\[
y_1+3y_2=5,\qquad 2y_1+2y_2=4\ \Rightarrow\ y_1+y_2=2\ \Rightarrow\ 2y_2=3,\quad y_2=1.5,\quad y_1=0.5 .
\]
检查：$w=8\times0.5+12\times1.5+0=4+18=22=z^*$ ✓；也可用 $y\T=c_B\T B^{-1}=(5,4,0)B^{-1}=(-\frac52+3,\ \frac52-1,\ 0)=(0.5,1.5,0)$ 验证。

\Lb{6. 第(3)问}烤箱影子价格 1.5（百元/小时）$>$ 租金 1，值得租。但要先确认“多 1 小时”仍在最优基不变的范围内：$\beta_2=(\frac12,-\frac14,\frac14)\T$，$x_B=(2,3,0.5)\T$：
$2+\frac12\Delta\ge0\Rightarrow\Delta\ge-4$；$3-\frac14\Delta\ge0\Rightarrow\Delta\le12$；$0.5+\frac14\Delta\ge0\Rightarrow\Delta\ge-2$。范围 $\Delta b_2\in[-2,12]$，$\Delta=1$ 在内，可以直接用影子价格：净增 $1\times(1.5-1)=0.5$（百元）。新方案 $x_B'=(2.5,\,2.75,\,0.75)$，$z=12.5+11=23.5=22+1.5$ ✓。

\Lb{7. 第(4)问}$\beta_1=(-\frac12,\frac34,-\frac34)\T$：
$2-\frac12\Delta\ge0\Rightarrow\Delta\le4$；$3+\frac34\Delta\ge0\Rightarrow\Delta\ge-4$；$0.5-\frac34\Delta\ge0\Rightarrow\Delta\le\frac23$。交集 $[-4,\frac23]$，即 $b_1\in[4,\frac{26}{3}]$。
\begin{figure}[!htbp]\centering
\begin{tikzpicture}[x=0.75cm,y=0.75cm]
  \begin{scope}
  \fill[feasreg] (0,0)--(4,0)--(2,3)--(1,3.5)--(0,3.5)--cycle;
  \axes{6}{5}{x_1}{x_2}
  \draw[con] (4.13,-0.2)--(0.8,4.8); \draw[con] (-0.2,3.5)--(6.3,3.5);
  \draw[con,black!35] (-0.4,4.2)--(6.3,0.85);
  \draw[shiftl] (-0.2,2.1)--(4.4,-0.2) node[below,font=\tiny]{$b_1=4$};
  \draw[shiftl] (-0.4,4.53)--(6.3,1.18) node[right,font=\tiny]{$b_1=\frac{26}{3}$};
  \node[optpt,fill=shiftc] at (4,0) {}; \node[optpt,fill=shiftc] at (1.667,3.5) {}; \node[optpt] at (2,3) {};
  \node[lab,font=\tiny,above left] at (1.6,3.5) {$(\frac53,\frac72)$};
  \node[font=\scriptsize,align=center] at (3,-1.4) {第 (4) 问：面粉边界在两条虚线间平移，\\最优点沿烤箱线滑动};
  \end{scope}
  \begin{scope}[xshift=6.6cm]
  \fill[feasreg] (0,0)--(4.333,0)--(2.5,2.75)--(1,3.5)--(0,3.5)--cycle;
  \axes{6}{5}{x_1}{x_2}
  \draw[con] (-0.4,4.2)--(6.3,0.85); \draw[con] (-0.2,3.5)--(6.3,3.5);
  \draw[con,black!35] (4.13,-0.2)--(0.8,4.8);
  \draw[shiftl] (4.47,-0.2)--(1.13,4.8) node[above,font=\tiny]{烤箱 13 h};
  \node[badpt] at (2,3) {}; \node[optpt,fill=shiftc] at (2.5,2.75) {};
  \node[lab,font=\tiny,right] at (2.6,2.85) {$(2.5,2.75)$};
  \node[font=\scriptsize,align=center] at (3,-1.4) {第 (3) 问：多租 1 小时烤箱，\\最优点沿面粉线移动，$z$ 增 1.5};
  \end{scope}
\end{tikzpicture}
\fig{第 34 题解析图：左为面粉的允许范围（端点为 $(4,0)$ 与 $(\frac53,\frac72)$），右为租烤箱后的新最优点}
\end{figure}
\Lb{8. 结果合理性检查}$b_1=\frac{26}{3}$ 时最优点 $(\frac53,\frac72)$ 恰好碰到奶油上限（$x_5=0$），$b_1=4$ 时到达 $(4,0)$（$x_2=0$），与左图端点一致。\par
\Lb{9. 最容易卡在哪里}第(3)问直接用影子价格而不核对范围；第(4)问把 $\beta_1$ 取成 $B^{-1}$ 的第一行。\par
\Lb{10. 如果条件变化}若面粉增加 1 袋（$\Delta=1>\frac23$），就越过了奶油的限制，不能按 $0.5$ 计算增利，需用对偶单纯形法续解。\par
\Lb{变式思考}奶油若只剩 3 盒，最优方案会变吗？（提示：$\beta_3=(0,0,1)\T$，范围 $b_3\ge3$，恰在端点，方案不变但出现退化。）

% ---------------- 35 ----------------
\A{35}
\Lb{1. 题目到底在问什么}在第 34 题的最优表上加一条约束，用对偶单纯形法重新优化，并解释利润减少量。\par
\Lb{2. 已知条件整理}$x_1=2+\frac12x_3-\frac12x_4$，$x_2=3-\frac34x_3+\frac14x_4$（由两行移项），$\sigma_3=-\frac12$，$\sigma_4=-\frac32$，$z^*=22$。\par
\Lb{3. 知识链定位}原解代入（判断）$\to$ 用基变量行消元（K1 矩阵视角）$\to$ 对偶单纯形法（K6）$\to$ 影子价格解释（K5）。

\Lb{4. 第(1)问}$x_1+x_2=2+3=5>4.5$，不满足。

\Lb{5. 第(2)问为什么这样做}新行里出现基变量 $x_1,x_2$，必须用它们所在的行消去，表中基列才保持单位阵：
\[
x_1+x_2+x_6=4.5\ \Rightarrow\ \Bigl(2+\tfrac12x_3-\tfrac12x_4\Bigr)+\Bigl(3-\tfrac34x_3+\tfrac14x_4\Bigr)+x_6=4.5
\ \Rightarrow\ \boxed{x_6-\tfrac14x_3-\tfrac14x_4=-\tfrac12}.
\]

\Lb{6. 第(3)问}$x_6$ 行右端 $-\frac12<0$，$x_6$ 出基。负系数列：$x_3$：$\frac{-1/2}{-1/4}=2$；$x_4$：$\frac{-3/2}{-1/4}=6$。$x_3$ 入基，主元 $-\frac14$。$x_3$ 行 $\times(-4)$：$x_3+x_4-4x_6=2$。代回：
$x_1=2+\frac12\cdot2=3$，$x_2=3-\frac34\cdot2=1.5$，$x_5=0.5+\frac34\cdot2=2$（$x_5$ 行 $x_5-\frac34x_3+\frac14x_4=0.5$）。新检验数 $\sigma'=\sigma-2\times(\text{原 }x_6\text{ 行})$：$\sigma_4'=-\frac32+\frac12=-1$，$\sigma_6'=-2$，$\sigma_3'=0$，全部 $\le0$；$b=(3,1.5,2,2)\ge0$，最优。
\begin{figure}[!htbp]\centering
\begin{tikzpicture}[x=0.8cm,y=0.8cm]
  \fill[feasreg] (0,0)--(4,0)--(3,1.5)--(1,3.5)--(0,3.5)--cycle;
  \fill[black!8] (3,1.5)--(2,3)--(1,3.5)--cycle;
  \axes{6}{5}{x_1}{x_2}
  \draw[con] (-0.4,4.2)--(6.3,0.85); \draw[con] (4.13,-0.2)--(0.8,4.8); \draw[con] (-0.2,3.5)--(6.3,3.5);
  \draw[shiftl] (0,4.5)--(4.8,-0.3) node[below,font=\scriptsize]{$x_1+x_2=4.5$};
  \node[badpt] at (2,3) {}; \node[optpt,fill=shiftc] at (3,1.5) {};
  \draw[arr,optc,thick] (2.08,2.88)--(2.92,1.62);
  \node[lab,font=\scriptsize,right] at (3.1,1.6) {$(3,1.5)$，$z=21$};
  \node[lab,font=\scriptsize,left] at (1.9,3.05) {原 $(2,3)$};
  \node[font=\scriptsize,black!60] at (2.3,2.75) {被切掉};
\end{tikzpicture}
\fig{第 35 题解析图：新约束切掉灰色三角形，原最优点 $(2,3)$ 被排除，最优点沿烤箱线滑到 $(3,1.5)$}
\end{figure}

\Lb{7. 最终答案}新方案：面包 3 批、蛋糕 1.5 批；面粉剩 2 袋，奶油剩 2 盒，烤箱和包装用尽；$z^*=5\times3+4\times1.5=21$。

\Lb{8. 第(4)问与合理性检查}利润减少 $22-21=1$。新约束的影子价格 $y_6=-\sigma_6'=2$，相当于“包装能力”从 5 批被压到 4.5 批，每少 1 批损失 2，少 0.5 批损失 $2\times0.5=1$ ✓。加约束后最优值只能变差 ✓。\par
\Lb{9. 最容易卡在哪里}第(2)问直接把 $(1,1,0,0,0,1\mid4.5)$ 写进表而不消元，导致基列不是单位阵。\par
\Lb{10. 如果条件变化}若包装上限为 5 批，原方案 $2+3=5$ 恰好满足，最优方案不变（结论 6.3 第 1 步）。\par
\Lb{变式思考}若包装上限降为 4 批，新行右端为 $-1$，按同样方法 $x_3$ 入基后 $x_3=4$，$x_2=0$，$x_1=4$，$z=20$。这一结果是否仍满足面粉、奶油约束？

% ============================================================
\section*{跨学科迁移训练详细解析}
\addcontentsline{toc}{section}{详细解析：跨学科迁移（36--38）}
\markboth{详细解析：跨学科迁移}{}

% ---------------- 36 ----------------
\A{36}
\SolHead{C}{基础～中等偏下}{K5 影子价格 $=\partial z^*/\partial b$}
\Lb{判断过程}按成本从低到高安排：G1 满发 60 MW，剩下 20 MW 由 G2 承担，$C^*=20\times60+30\times20=1800$。负荷增至 81 MW 时，G1 已满，多出的 1 MW 只能由 G2 提供，$C^*=1830$，增加 30。所以边际电价为 30 元/MWh。
\begin{figure}[!htbp]\centering
\begin{tikzpicture}[x=0.021cm,y=0.05cm]
  \foreach \i/\L/\p/\ok in {0/A/20/0,1/B/25/0,2/C/30/1,3/D/50/0} {
    \begin{scope}[xshift=\i*4.1cm]
      \fill[auxblue!35] (0,0) rectangle (60,20); \fill[auxblue!65] (60,0) rectangle (160,30);
      \draw[ax] (0,0)--(170,0); \draw[ax] (0,0)--(0,58);
      \draw[redline,dashed] (80,0)--(80,55);
      \ifnum\ok=1 \def\cc{greenline}\else\def\cc{badc}\fi
      \draw[\cc,very thick] (0,\p)--(165,\p);
      \ifnum\ok=1 \fill[greenline] (80,30) circle (4pt);\fi
      \node[font=\scriptsize] at (80,66) {\textbf{\L}.\ \p};
    \end{scope}
  }
\end{tikzpicture}
\fig{第 36 题选项图：机组按成本排成“阶梯”，红虚线为负荷 80 MW。只有 C 的价格线经过负荷处的阶梯高度}
\end{figure}
\Lb{各选项（看图）}A：20 是 G1 的成本，但 G1 已满发，边际上用不到它；B：25 是平均值，没有意义；D：50 是两者之和，混淆了概念。\par
\Lb{当前专题知识}影子价格 $y_i=\partial z^*/\partial b_i$（第 4 章定理 4.3）。\par
\Lb{借用的外部情境}电力系统的经济调度与边际电价。\par
\Lb{为什么不算超纲}题干给出了完整的线性规划模型和“边际电价”的定义，不需要任何电力系统知识。\par
\Lb{迁移点}把“资源 $b_i$ 多一单位 $\to$ 利润增加 $y_i$”换成“负荷多 1 MW $\to$ 最低成本增加多少”：负荷就是右端项，边际电价就是这条等式约束的影子价格。

% ---------------- 37 ----------------
\A{37}
\SolHead{ABC}{中等偏下}{K3 强对偶、K4 互补松弛、K5 影子价格}
\begin{figure}[!htbp]\centering
\begin{tikzpicture}
  \node at (0,0) {\begin{tikzpicture}[x=0.33cm,y=0.33cm]
    \fill[feasreg] (0,9.5)--(0,8)--(3,2)--(7,0)--(9.5,0)--(9.5,9.5)--cycle;
    \draw[ax] (-0.2,0)--(9.8,0) node[right,font=\tiny]{$a$}; \draw[ax] (0,-0.2)--(0,9.8) node[above,font=\tiny]{$b$};
    \draw[con] (0,8)--(4.4,-0.8); \draw[con] (-0.6,3.8)--(9.4,-1.2);
    \draw[iso] (0,6.5)--(4.333,0);
    \node[optpt] at (3,2) {};
    \node[font=\scriptsize,align=center] at (4.8,11.2) {\textbf{A} ✓：最低花费\\$3\times3+2\times2=13$};
  \end{tikzpicture}};
  \node at (4.3,0) {\begin{tikzpicture}[x=0.33cm,y=0.33cm]
    \fill[feasreg] (0,9.5)--(0,9)--(3.667,1.667)--(7,0)--(9.5,0)--(9.5,9.5)--cycle;
    \draw[ax] (-0.2,0)--(9.8,0); \draw[ax] (0,-0.2)--(0,9.8);
    \draw[con,black!30] (0,8)--(4.4,-0.8); \draw[shiftl] (0,9)--(4.9,-0.8); \draw[con] (-0.6,3.8)--(9.4,-1.2);
    \node[badpt] at (3,2) {}; \node[optpt,fill=shiftc] at (3.667,1.667) {};
    \node[font=\scriptsize,align=center] at (4.8,11.2) {\textbf{B} ✓：蛋白质 $8\to9$\\新点 $(\frac{11}{3},\frac53)$，$+\frac43$};
  \end{tikzpicture}};
  \node at (8.6,0) {\begin{tikzpicture}[x=1.5cm,y=1.5cm]
    \fill[feasreg] (0,0)--(1.5,0)--(1.333,0.333)--(0,1)--cycle;
    \draw[ax] (-0.05,0)--(1.75,0) node[right,font=\tiny]{$u$}; \draw[ax] (0,-0.05)--(0,1.65) node[above,font=\tiny]{$v$};
    \draw[con] (1.55,-0.1)--(0.75,1.5); \draw[con] (-0.1,1.05)--(1.75,0.125);
    \draw[iso] (0,1.857)--(1.625,0);
    \node[optpt] at (1.333,0.333) {};
    \node[font=\tiny,above right] at (1.34,0.34) {$(\frac43,\frac13)$};
    \node[font=\scriptsize,align=center] at (0.85,2.25) {\textbf{C} ✓：对偶最优\\$8u+7v=13$};
  \end{tikzpicture}};
  \node at (12.6,0) {\begin{tikzpicture}[x=1cm,y=1cm]
    \fill[auxblue!60] (0,0) rectangle (0.5,1.33); \node[font=\tiny,above] at (0.25,1.33) {$u=\frac43$};
    \fill[auxblue!60] (0.9,0) rectangle (1.4,0.33); \node[font=\tiny,above] at (1.15,0.33) {$v=\frac13$};
    \draw[ax] (-0.2,0)--(1.7,0);
    \node[font=\scriptsize,align=center] at (0.75,2.25) {\textbf{D} ✗：两种营养\\都用“刚好够”，\\价格都为正};
  \end{tikzpicture}};
\end{tikzpicture}
\fig{第 37 题选项图：A、B 在原问题可行域上看，C 在对偶可行域上看，D 看影子价格的符号}
\end{figure}
\Lb{A}正确。最优点 $(3,2)$ 是两条营养约束的交点，$C^*=9+4=13$；图 A 中虚线等值线 $3a+2b=13$ 在此支撑可行域。\par
\Lb{B}正确。由互补松弛：$a>0\Rightarrow2u+v=3$，$b>0\Rightarrow u+2v=2$，解得 $u=\frac43$，$v=\frac13$。蛋白质需求 $8\to9$：新顶点 $a=\frac{2\cdot9-7}{3}=\frac{11}{3}$，$b=\frac{2\cdot7-9}{3}=\frac53$，均 $\ge0$（最优基不变），$C=11+\frac{10}{3}=\frac{43}{3}=13+\frac43$ ✓。\par
\Lb{C}正确。对偶问题恰是“营养片厂商”的定价问题；强对偶给出 $8u^*+7v^*=\frac{32}{3}+\frac73=13$。\par
\Lb{D}错误。$a,b>0$ 推出的是\textbf{对偶约束}取等（产品侧互补松弛），而不是营养约束的影子价格为 0。两条营养约束在最优点都恰好满足（$2\cdot3+2=8$，$3+4=7$），影子价格 $\frac43,\frac13$ 都为正。\par
\Lb{当前专题知识}强对偶、互补松弛、影子价格。\quad\Lb{借用的外部情境}营养学中的配餐问题（min 型，$\ge$ 约束）。\par
\Lb{为什么不算超纲}题干给出了每种食物的营养含量和需求量，以及对偶问题的经济解释，不需要营养学知识。\par
\Lb{迁移点}生产问题中“资源的影子价格”在配餐问题中变成“营养素的边际成本”：需求每多 1 单位，最低花费增加多少。

% ---------------- 38 ----------------
\A{38}
\SolHead{40 元；1 元}{基础}{K3 强对偶 $z^*=y\T b$；K5 影子价格}
\begin{figure}[!htbp]\centering
\begin{tikzpicture}[x=1cm,y=0.12cm]
  \draw[ax] (0,0)--(7.4,0);
  \fill[auxblue!70] (0.6,0) rectangle (1.6,20); \node[font=\scriptsize,white] at (1.1,10) {$0.5\times40$};
  \fill[shiftc!70] (0.6,20) rectangle (1.6,40); \node[font=\scriptsize] at (1.1,30) {$0.2\times100$};
  \node[font=\scriptsize,right] at (1.7,40) {$=40$};
  \node[font=\scriptsize,below] at (1.1,0) {最大收益 $y\T b$};
  \fill[greenline!70] (4.4,0) rectangle (5.4,12.5); \node[font=\scriptsize,above] at (4.9,12.5) {增收 $0.5\times5=2.5$};
  \fill[badc!60] (5.8,0) rectangle (6.8,7.5); \node[font=\scriptsize,above] at (6.3,7.5) {租金 1.5};
  \node[font=\scriptsize,below] at (5.6,0) {净增 $2.5-1.5=1$};
\end{tikzpicture}
\fig{第 38 题解析图：左柱为收益按影子价格分摊到两种资源，右边两柱为租用 CPU 的增收与租金}
\end{figure}
\Lb{推理过程}第一空：由强对偶，最大收益 $=y\T b=0.5\times40+0.2\times100=20+20=40$ 元。第二空：在允许范围内，多 5 核时增收 $0.5\times5=2.5$ 元，租金 $0.3\times5=1.5$ 元，净增 1 元。\par
\Lb{当前专题知识}强对偶定理 $z^*=y^{*\mathsf T}b$；影子价格的“购入决策”用法。\quad\Lb{借用的外部情境}云计算平台的资源调度。\par
\Lb{为什么不算超纲}题干直接给出影子价格与容量，并说明了约束形式，不需要了解任务调度的细节。\par
\Lb{迁移点}第 4 章“影子价格加权求和等于最优值”（课件 P20）与“市场价低于影子价则买入”（课件 P29），原样搬到计算资源上。\par
\Lb{易错点}第二空忘记扣除租金，答成 2.5。

% ============================================================
%  参考资料
% ============================================================
\section*{参考资料}
\addcontentsline{toc}{section}{参考资料}
\markboth{参考资料}{}

\subsection*{原始教学依据}
\begin{enumerate}[label={[\arabic*]},leftmargin=2.2em]
  \item 运筹学课堂手写笔记，第四章“修正单纯形法与对偶”“对偶问题与灵敏度分析”，2026 年 10 月 14--21 日，共 16 页（Memo No.~532--535、557--559、564--568、571、607--608 及一页未编号照片）。
  \item 运筹学课件《对偶单纯形法：从对偶可行出发，逐步恢复原始可行——一种无需人工变量的求解路径》，共 85 页。
  \item 运筹学课件（第 17、18 节 / 共 32 节），山东大学控制科学与工程学院：影子价格与线性规划灵敏度分析，共 73 页。
\end{enumerate}

\subsection*{核心定义与公式核验}
\begin{enumerate}[label={[\arabic*]},leftmargin=2.2em,resume]
  \item 《运筹学》教材编写组．运筹学（第 4 版）．北京：清华大学出版社，2012．（第 2 章 线性规划的对偶理论与灵敏度分析）
  \item F.~S.~Hillier, G.~J.~Lieberman. \emph{Introduction to Operations Research}. New York: McGraw-Hill.（对偶理论、对偶单纯形法与灵敏度分析各章）
  \item D.~Bertsimas, J.~N.~Tsitsiklis. \emph{Introduction to Linear Optimization}. Belmont: Athena Scientific, 1997.（第 3 章修正单纯形法，第 4 章对偶理论与互补松弛，第 5 章灵敏度分析）
  \item R.~J.~Vanderbei. \emph{Linear Programming: Foundations and Extensions}. New York: Springer.（对偶定理的证明与对偶单纯形法）
\end{enumerate}

\subsection*{学科史资料}
\begin{enumerate}[label={[\arabic*]},leftmargin=2.2em,resume]
  \item G.~B.~Dantzig. \emph{Linear Programming and Extensions}. Princeton: Princeton University Press, 1963.
  \item G.~B.~Dantzig. Linear programming. \emph{Operations Research}, 2002, 50(1): 42--47.（丹齐格对单纯形法与对偶思想起源的回忆）
  \item L.~V.~Kantorovich. Mathematical methods of organizing and planning production. \emph{Management Science}, 1960, 6(4): 366--422.（1939 年俄文原著的英译本）
  \item D.~Gale, H.~W.~Kuhn, A.~W.~Tucker. Linear programming and the theory of games. In: T.~C.~Koopmans (ed.), \emph{Activity Analysis of Production and Allocation}. New York: Wiley, 1951.
  \item C.~E.~Lemke. The dual method of solving the linear programming problem. \emph{Naval Research Logistics Quarterly}, 1954, 1(1): 36--47.
  \item G.~B.~Dantzig, W.~Orchard-Hays. The product form for the inverse in the simplex method. \emph{Mathematical Tables and Other Aids to Computation}, 1954, 8: 64--67.
  \item 诺贝尔奖官方网站：1975 年诺贝尔经济学奖授予 L.~V.~Kantorovich 与 T.~C.~Koopmans，“表彰他们对资源最优配置理论的贡献”。
\end{enumerate}

\tip[说明]{第 0 章“从课本到现实”中的电力市场节点电价、航空收益管理与求解器重新优化，属于运筹学与相关工程领域的常识性应用介绍，讲义只描述其基本原理，不涉及具体市场规则或软件实现细节；生活例子中的阶梯电价数字为示意数字。}

\vspace{1cm}
\begin{center}
\begin{tikzpicture}
  \node[gbox,text width=13cm,align=center,font=\small] {\textbf{全篇一句话}\\[2pt]
  单纯形表里的 $B^{-1}$ 既是计算的工具（修正单纯形法），也是对偶的解（$y\T=c_B\T B^{-1}$），更是资源的价格（影子价格）；\\
  守住 $\sigma\le0$ 去找可行，就是对偶单纯形法；分开核验 $B^{-1}b\ge0$ 与 $\sigma\le0$，就是灵敏度分析。};
\end{tikzpicture}
\end{center}

\end{document}
